\documentclass[fs95,seriesFre]{korfarm-base}
\usepackage{korfarm-frege}

% ============================================================
% passage-note.tex (v11)
% 지문 박스 + 첫 페이지 우측 메모란 (동적 높이)
% 정책: PAGE_BREAK_POLICY.md §3, §4 참조
%
% 변경 (v11):
%  · 지문 박스 폭 확대: enlarge right by=-3.7cm  (메모 3.4cm + gap 3mm)
%  · 메모란 동적 높이: frame.north east ~ frame.south east 사이 자동 채움
%  · 메모 줄선: clip 영역으로 박스 내부만 채움 (박스 높이 모르고도 작동)
%  · 문단 들여쓰기 활성화: tcolorbox 내부 \parindent=1em, \noindent 제거
%
% 사용:
%   \kfPassageNote{제목}{지문 본문}      → 메모란 포함
%   \kfPassageNoteNT{지문 본문}          → 메모란 포함, 제목 없음
%   \kfPassageFull{제목}{본문}           → 메모란 없음 (활동 내 보기)
% ============================================================

\makeatletter

% 메모란 규격 (cm 고정)
\newlength{\kf@memoW}
\setlength{\kf@memoW}{3.4cm}
\newlength{\kf@memoGap}
\setlength{\kf@memoGap}{3mm}

% --- 메인: 제목 있는 버전 ---
\newcommand{\kf@passagenote@with}[2]{%
  \par\ifkfPassageNewPage\ifdim\pagetotal>0.4\textheight\clearpage\fi\fi\smallskip\needspace{6\baselineskip}%
  \begin{tcolorbox}[
    enhanced, breakable,
    width=\dimexpr\linewidth-4cm\relax,
    colback=kfPassageBg, colframe=kfPrimary!70,
    boxrule=0.6pt, arc=3pt,
    left=\kfBoxPad, right=\kfBoxPad, top=\dimexpr\kfBoxPad+8pt\relax, bottom=\kfBoxPad,
    title={\small\bfseries\color{kfPrimary} #1},
    coltitle=kfPrimary,
    colbacktitle=kfPrimaryLight!60,
    attach boxed title to top left={yshift=-2.4mm, xshift=8pt},
    boxed title style={colframe=kfPrimary!50, boxrule=0.5pt, arc=2pt},
    parbox=false,
    overlay unbroken and first={%
      \begin{scope}
        \draw[kfRule, line width=0.4pt, fill=kfNoteBg, rounded corners=2pt]
          ([xshift=\kf@memoGap]frame.north east) rectangle
          ([xshift=\kf@memoGap+\kf@memoW]frame.south east);
        \node[anchor=north east, font=\footnotesize\bfseries\color{kfMute}, inner sep=3pt]
          at ([xshift=\kf@memoGap+\kf@memoW, yshift=-2pt]frame.north east) {메모};
        \node[anchor=north east, fill=kfPrimary!15, text=kfPrimary,
              font=\scriptsize\bfseries, rounded corners=2pt,
              inner xsep=4pt, inner ysep=2pt]
          at ([xshift=\kf@memoGap+\kf@memoW-3pt, yshift=-19pt]frame.north east) {\kf@memocat};
        \begin{scope}
          \path[clip]
            ([xshift=\kf@memoGap]frame.north east) rectangle
            ([xshift=\kf@memoGap+\kf@memoW]frame.south east);
          \foreach \i in {1,...,40}{%
            \draw[kfRule, line width=0.25pt]
              ([xshift=\kf@memoGap+2mm, yshift=-7mm - \i*5mm]frame.north east) --
              ([xshift=\kf@memoGap+\kf@memoW-2mm, yshift=-7mm - \i*5mm]frame.north east);%
          }
        \end{scope}
      \end{scope}
    }
  ]
    \setlength{\parindent}{1em}%
    \hspace*{1em}\ignorespaces#2%
  \end{tcolorbox}\par\smallskip
}

% --- 제목 없는 버전 ---
\newcommand{\kf@passagenote@notitle}[1]{%
  \par\ifkfPassageNewPage\ifdim\pagetotal>0.4\textheight\clearpage\fi\fi\smallskip\needspace{6\baselineskip}%
  \begin{tcolorbox}[
    enhanced, breakable,
    width=\dimexpr\linewidth-4cm\relax,
    colback=kfPassageBg, colframe=kfPrimary!70,
    boxrule=0.6pt, arc=3pt,
    left=\kfBoxPad, right=\kfBoxPad, top=\dimexpr\kfBoxPad+8pt\relax, bottom=\kfBoxPad,
    parbox=false,
    overlay unbroken and first={%
      \begin{scope}
        \draw[kfRule, line width=0.4pt, fill=kfNoteBg, rounded corners=2pt]
          ([xshift=\kf@memoGap]frame.north east) rectangle
          ([xshift=\kf@memoGap+\kf@memoW]frame.south east);
        \node[anchor=north east, font=\footnotesize\bfseries\color{kfMute}, inner sep=3pt]
          at ([xshift=\kf@memoGap+\kf@memoW, yshift=-2pt]frame.north east) {메모};
        \node[anchor=north east, fill=kfPrimary!15, text=kfPrimary,
              font=\scriptsize\bfseries, rounded corners=2pt,
              inner xsep=4pt, inner ysep=2pt]
          at ([xshift=\kf@memoGap+\kf@memoW-3pt, yshift=-19pt]frame.north east) {\kf@memocat};
        \begin{scope}
          \path[clip]
            ([xshift=\kf@memoGap]frame.north east) rectangle
            ([xshift=\kf@memoGap+\kf@memoW]frame.south east);
          \foreach \i in {1,...,40}{%
            \draw[kfRule, line width=0.25pt]
              ([xshift=\kf@memoGap+2mm, yshift=-7mm - \i*5mm]frame.north east) --
              ([xshift=\kf@memoGap+\kf@memoW-2mm, yshift=-7mm - \i*5mm]frame.north east);%
          }
        \end{scope}
      \end{scope}
    }
  ]
    \setlength{\parindent}{1em}%
    \hspace*{1em}\ignorespaces#1%
  \end{tcolorbox}\par\smallskip
}

\newcommand{\kfPassageNote}[2]{%
  \def\kf@tmptitle{#1}%
  \ifx\kf@tmptitle\@empty
    \kf@passagenote@notitle{#2}%
  \else
    \kf@passagenote@with{#1}{#2}%
  \fi
}
\newcommand{\kfPassageNoteNT}[1]{\kf@passagenote@notitle{#1}}
\makeatother

% ============================================================
% 풀폭 지문 박스 (메모란 없음) — 활동 내 보기 박스
% PAGE_BREAK_POLICY.md §3
% ============================================================
\makeatletter
\newcommand{\kf@passagefull@with}[2]{%
  \par\ifkfPassageNewPage\ifdim\pagetotal>0.4\textheight\clearpage\fi\fi\smallskip\needspace{6\baselineskip}%
  \begin{tcolorbox}[
    enhanced, breakable,
    colback=kfPassageBg, colframe=kfPrimary!70,
    boxrule=0.6pt, arc=3pt,
    left=\kfBoxPad, right=\kfBoxPad, top=\dimexpr\kfBoxPad+8pt\relax, bottom=\kfBoxPad,
    title={\small\bfseries\color{kfPrimary} #1},
    coltitle=kfPrimary,
    colbacktitle=kfPrimaryLight!60,
    attach boxed title to top left={yshift=-2.4mm, xshift=8pt},
    boxed title style={colframe=kfPrimary!50, boxrule=0.5pt, arc=2pt},
    parbox=false
  ]
    \setlength{\parindent}{1em}%
    \hspace*{1em}\ignorespaces#2%
  \end{tcolorbox}\par\smallskip
}
\newcommand{\kf@passagefull@notitle}[1]{%
  \par\ifkfPassageNewPage\ifdim\pagetotal>0.4\textheight\clearpage\fi\fi\smallskip\needspace{6\baselineskip}%
  \begin{tcolorbox}[
    enhanced, breakable,
    colback=kfPassageBg, colframe=kfPrimary!70,
    boxrule=0.6pt, arc=3pt,
    left=\kfBoxPad, right=\kfBoxPad, top=\dimexpr\kfBoxPad+8pt\relax, bottom=\kfBoxPad,
    parbox=false
  ]
    \setlength{\parindent}{1em}%
    \hspace*{1em}\ignorespaces#1%
  \end{tcolorbox}\par\smallskip
}
\newcommand{\kfPassageFull}[2]{%
  \def\kf@tmptitle{#1}%
  \ifx\kf@tmptitle\@empty
    \kf@passagefull@notitle{#2}%
  \else
    \kf@passagefull@with{#1}{#2}%
  \fi
}
\newcommand{\kfPassageFullNT}[1]{\kf@passagefull@notitle{#1}}
\makeatother

% ============================================================
% 작은 문장 박스 (문장 독해용) — PAGE_BREAK_POLICY.md §2.5
% ============================================================
\newcommand{\kfSentenceBox}[2]{%
  \par\smallskip\needspace{3\baselineskip}%
  \noindent\begin{tcolorbox}[
    enhanced, breakable=false,
    colback=kfPassageBg, colframe=kfMute,
    boxrule=0.4pt, arc=2pt,
    left=8pt, right=6pt, top=4pt, bottom=4pt,
    title={\footnotesize\bfseries\color{kfPrimary} #1},
    coltitle=kfPrimary, colbacktitle=kfPaper,
    attach boxed title to top left={yshift=-1.6mm, xshift=4pt},
    boxed title style={colframe=kfMute, boxrule=0.3pt, arc=1pt}
  ]%
    \setstretch{1.3}#2%
  \end{tcolorbox}\par\smallskip%
}

% ============================================================
% 지문 본문 아래 안내/정리 박스 (v16 신규)
% 사용: \kfPassageHint{한 줄 정리}{본문 안내 텍스트}
% ============================================================
\newcommand{\kfPassageHint}[2]{%
  \par\smallskip\needspace{4\baselineskip}%
  \begin{tcolorbox}[
    enhanced, breakable=false,
    colback=kfPrimary!5, colframe=kfPrimary!40,
    boxrule=0.4pt, arc=2pt,
    left=8pt, right=6pt, top=4pt, bottom=4pt,
    title={\footnotesize\bfseries\color{kfPrimary} #1},
    coltitle=kfPrimary, colbacktitle=kfPaper,
    attach boxed title to top left={yshift=-1.5mm, xshift=4pt},
    boxed title style={colframe=kfPrimary!40, boxrule=0.3pt, arc=1pt}
  ]%
    \small\setstretch{1.3}#2%
  \end{tcolorbox}\par\smallskip%
}

% ============================================================
% 환경 버전 (호환성) — 메모란 포함
% ============================================================
\makeatletter
\newenvironment{kfPassageNoteEnv}[1]%
  {\par\needspace{6\baselineskip}%
   \def\kf@psrc{#1}%
   \begin{tcolorbox}[
     enhanced, breakable,
     width=\dimexpr\linewidth-4cm\relax,
     colback=kfPassageBg, colframe=kfPrimary!70,
     boxrule=0.6pt, arc=3pt,
     left=\kfBoxPad, right=\kfBoxPad, top=\dimexpr\kfBoxPad+8pt\relax, bottom=\kfBoxPad,
     title={\small\bfseries\color{kfPrimary} \kf@psrc},
     coltitle=kfPrimary, colbacktitle=kfPrimaryLight!60,
     attach boxed title to top left={yshift=-2.4mm, xshift=8pt},
     boxed title style={colframe=kfPrimary!50, boxrule=0.5pt, arc=2pt},
     parbox=false,
     overlay unbroken and first={%
       \begin{scope}
         \draw[kfRule, line width=0.4pt, fill=kfNoteBg, rounded corners=2pt]
           ([xshift=\kf@memoGap]frame.north east) rectangle
           ([xshift=\kf@memoGap+\kf@memoW]frame.south east);
         \node[anchor=north east, font=\footnotesize\bfseries\color{kfMute}, inner sep=3pt]
           at ([xshift=\kf@memoGap+\kf@memoW, yshift=-2pt]frame.north east) {메모};
         \begin{scope}
           \path[clip]
             ([xshift=\kf@memoGap]frame.north east) rectangle
             ([xshift=\kf@memoGap+\kf@memoW]frame.south east);
           \foreach \i in {1,...,40}{%
             \draw[kfRule, line width=0.25pt]
               ([xshift=\kf@memoGap+2mm, yshift=-7mm - \i*5mm]frame.north east) --
               ([xshift=\kf@memoGap+\kf@memoW-2mm, yshift=-7mm - \i*5mm]frame.north east);%
           }
         \end{scope}
       \end{scope}
     }
   ]%
   \setlength{\parindent}{1em}%
  }%
  {\end{tcolorbox}\par\smallskip}
\makeatother

% ============================================================
% condition-box.tex
% <조건> 박스 컴포넌트 (tcolorbox)
% 사용:
%   \begin{kfCondition}
%     \item 첫째 조건
%     \item 둘째 조건
%   \end{kfCondition}
% ============================================================

\newenvironment{kfCondition}%
  {\par\smallskip\needspace{4\baselineskip}%
   \begin{tcolorbox}[
     enhanced, breakable=false,
     colback=kfPrimaryLight!40, colframe=kfPrimary,
     boxrule=0.5pt, arc=2pt,
     left=\kfBoxPad, right=\kfBoxPad, top=\kfBoxPad, bottom=\kfBoxPad,
     title={\small\bfseries\color{kfPrimary} \,조건\,},
     coltitle=kfPrimary, colbacktitle=kfPaper,
     attach boxed title to top left={yshift=-2mm, xshift=6pt},
     boxed title style={colframe=kfPrimary, boxrule=0.4pt, arc=1pt}
   ]%
   \begin{enumerate}[leftmargin=16pt, itemsep=1pt, topsep=1pt,
                    label=\textbf{\arabic*.}]}%
  {\end{enumerate}\end{tcolorbox}\par\smallskip}

% ============================================================
% evidence-box.tex
% <보기> 박스 컴포넌트
% 사용:
%   \begin{kfEvidence}
%     보기 본문 ...
%   \end{kfEvidence}
% ============================================================

\newenvironment{kfEvidence}%
  {\par\smallskip\needspace{4\baselineskip}%
   \begin{tcolorbox}[
     enhanced, breakable=false,
     colback=kfPaper, colframe=kfMute,
     boxrule=0.5pt, arc=2pt,
     left=\kfBoxPad, right=\kfBoxPad, top=\kfBoxPad, bottom=\kfBoxPad,
     title={\small\bfseries\color{kfInk} \,보기\,},
     coltitle=kfInk, colbacktitle=kfPrimaryLight!50,
     attach boxed title to top left={yshift=-2mm, xshift=6pt},
     boxed title style={colframe=kfMute, boxrule=0.4pt, arc=1pt}
   ]%
   \leavevmode\mbox{}\ignorespaces%
  }%
  {\end{tcolorbox}\par\smallskip}

% ============================================================
% choice-list.tex
% 5선지 (① ② ③ ④ ⑤) 자동 들여쓰기 컴포넌트
% 사용:
%   \begin{kfChoices}
%     \kfChoice 첫째 선택지
%     \kfChoice 둘째 선택지
%     \kfChoice 셋째 선택지
%     \kfChoice 넷째 선택지
%     \kfChoice 다섯째 선택지
%   \end{kfChoices}
% ============================================================

\newcounter{kfChoiceCnt}
\newcommand{\kfChoiceMark}[1]{%
  \ifcase#1\relax%
  \or ①\or ②\or ③\or ④\or ⑤\or ⑥\or ⑦\or ⑧\or ⑨%
  \fi
}

% 환경: 자동 번호
\newenvironment{kfChoices}%
  {\setcounter{kfChoiceCnt}{0}%
   \par\smallskip%
   \begin{list}{}{%
     \setlength{\leftmargin}{20pt}%
     \setlength{\itemindent}{0pt}%
     \setlength{\labelwidth}{18pt}%
     \setlength{\labelsep}{6pt}%
     \setlength{\itemsep}{2pt}%
     \setlength{\topsep}{1pt}%
     \setlength{\parsep}{0pt}%
     \renewcommand{\makelabel}[1]{##1\hss}%
   }}%
  {\end{list}\par\smallskip}

\newcommand{\kfChoice}{%
  \stepcounter{kfChoiceCnt}%
  \item[\textbf{\kfChoiceMark{\value{kfChoiceCnt}}}]\ignorespaces%
}

% --- 한 줄 5선지 (짧은 선택지용) ---
\newcommand{\kfChoicesInline}[5]{%
  \par\smallskip%
  \noindent\textbf{①}~#1\quad\textbf{②}~#2\quad\textbf{③}~#3\quad\textbf{④}~#4\quad\textbf{⑤}~#5%
  \par\smallskip%
}

% ============================================================
% answer-note.tex
% 답안 작성란 (노트 스타일, 줄친 영역)
% 사용:
%   \kfAnswerLines{5}        % 5줄 답안란
%   \kfAnswerNote{서술형 답안란}{8} % 라벨+8줄
% ============================================================

% 단일 줄 (필요 시 인라인)
\newcommand{\kfAnswerLine}{%
  \par\noindent\textcolor{kfRule}{\rule{\linewidth}{0.4pt}}\par
}

% N줄 답안란 (간단)
\newcommand{\kfAnswerLines}[1]{%
  \par\smallskip%
  \begin{minipage}{\linewidth}%
  \foreach \i in {1,...,#1}{%
    \textcolor{kfRule}{\rule{\linewidth}{0.4pt}}\\[6pt]%
  }%
  \end{minipage}%
  \par\smallskip%
}

% 라벨이 있는 노트형 답안란
\newcommand{\kfAnswerNote}[2]{%
  \par\smallskip\needspace{#2\baselineskip}%
  \begin{tcolorbox}[
    enhanced, breakable=false,
    colback=kfPaper, colframe=kfRule,
    boxrule=0.4pt, arc=2pt,
    left=\kfBoxPad, right=\kfBoxPad, top=\kfBoxPad, bottom=\kfBoxPad,
    title={\footnotesize\bfseries\color{kfMute} #1},
    coltitle=kfMute, colbacktitle=kfPaper,
    attach boxed title to top left={yshift=-1.5mm, xshift=6pt},
    boxed title style={colframe=kfRule, boxrule=0.3pt, arc=1pt}
  ]%
  \foreach \i in {1,...,#2}{%
    \textcolor{kfRule}{\rule{\linewidth}{0.3pt}}\\[6pt]%
  }%
  \end{tcolorbox}\par\smallskip%
}

% 짧은 빈칸 (단답형)
\newcommand{\kfBlank}[1][3cm]{\underline{\hspace{#1}}}
% 넓은 빈칸 (서술 한 줄)
\newcommand{\kfBlankWide}[1][6cm]{\underline{\hspace{#1}}}
% 괄호 빈칸
\newcommand{\kfBlankParen}{(\,\hspace{1cm}\,)}
% O/X 빈칸
\newcommand{\kfBlankOX}{(\,\hspace{1.5cm}\,)}

% ============================================================
% question-block.tex
% 문제 블록 (번호 배지 + 발문 + 보기/조건 + 선택지 + 답안란)
% 모든 구성을 하나의 minipage로 묶어 단/페이지 단절 금지
%
% 사용:
%   \begin{kfQBlock}{1}{객관식}
%     \kfQStem{다음 글에 대한 설명으로 가장 적절한 것은?}
%     \begin{kfEvidence} ... \end{kfEvidence}    % 선택
%     \begin{kfCondition} ... \end{kfCondition}  % 선택
%     \begin{kfChoices}
%       \kfChoice ...
%     \end{kfChoices}
%     \kfAnswerLines{2}                            % 선택
%   \end{kfQBlock}
% ============================================================

% 무단절 minipage로 묶고, 발문/보기/조건/선택지/답안란을 자유롭게 배치
% 사용자 피드백 #4: [객관식]/[서술형] 등 텍스트 라벨 제거 — 번호 배지만 표시
% #2(유형 라벨)는 호환성 인자로만 받고 출력하지 않음.
\newenvironment{kfQBlock}[2]%
  {\par\smallskip\needspace{6\baselineskip}%
   \noindent\begin{minipage}{\linewidth}%
   \samepage%
   \par\noindent
   \kfQNum{#1}\hspace{4pt}%
   \begin{minipage}[t]{\dimexpr\linewidth-30pt\relax}%
  }%
  {\end{minipage}\end{minipage}\par\smallskip}

% 문제 발문
\newcommand{\kfQStem}[1]{%
  \par\noindent#1\par\vspace{2pt}%
}

% 짧은 소문항 라벨 (1), (2), (3)
\newcounter{kfSubQ}
\newcommand{\kfSubQReset}{\setcounter{kfSubQ}{0}}
\newcommand{\kfSubQ}{%
  \stepcounter{kfSubQ}%
  \par\noindent\textbf{(\arabic{kfSubQ})}~\ignorespaces%
}

% ============================================================
% activity-table.tex
% 활동: 정리표 / 내용 확인표 (마크다운 표 → tabularx 변환)
% 사용:
%   % 일반 tabular (직접 컬럼 명시)
%   \begin{kfActTable}{|l|X|X|}
%     \kfTHead{항목} & \kfTHead{설명} & \kfTHead{학생 작성} \\\hline
%     ① & ... & \kfTBlank \\\hline
%   \end{kfActTable}
%
%   % adjustbox 자동 축소 (정해진 폭으로 강제)
%   \kfActTableWrap{
%     ... (\begin{tabular}...\end{tabular} 코드) ...
%   }
% ============================================================

% 사용 시 본문 마지막 행은 \\\hline 으로 끝나야 함.
% v17.1: 흰색 mask로 Canva 배경 본문 침범 차단
\newenvironment{kfActTable}[1]%
  {\par\smallskip\needspace{4\baselineskip}%
   \renewcommand{\arraystretch}{1.3}%
   \arrayrulecolor{kfRule}%
   \begin{tcolorbox}[blanker, colback=white, left=2pt, right=2pt, top=2pt, bottom=2pt]%
   \noindent\begin{adjustbox}{max width=\linewidth}%
   \begin{tabular}{#1}%
   \hline}%
  {\end{tabular}%
   \end{adjustbox}%
   \end{tcolorbox}%
   \arrayrulecolor{black}%
   \par\smallskip}

% 헤더 행 헬퍼
\newcommand{\kfTHead}[1]{\textbf{\color{kfPrimary} #1}}
\newcommand{\kfTHeadRow}[1]{\rowcolor{kfPrimaryLight!50}\textbf{\color{kfPrimary} #1}}

% 빈 셀 (학생 작성용)
\newcommand{\kfTBlank}{\rule[-1.6em]{0pt}{3.4em}}
\newcommand{\kfTBlankSm}{\rule[-1.0em]{0pt}{2.4em}}

% adjustbox 강제 축소 래퍼 — Python에서 변환된 raw tabular 블록을 감쌀 때 사용
\newcommand{\kfActTableWrap}[1]{%
  \par\smallskip\needspace{4\baselineskip}%
  \begin{tcolorbox}[blanker, colback=white, left=2pt, right=2pt, top=2pt, bottom=2pt]%
  \noindent\begin{adjustbox}{max width=\linewidth, center}%
  #1%
  \end{adjustbox}%
  \end{tcolorbox}\par\smallskip%
}

% ============================================================
% activity-vocab.tex
% 어휘 학습 표 (3열: 번호 - 뜻 - 빈칸)
% 사용:
%   \begin{kfVocab}
%     \kfVocabItem{1}{눈으로 보고 마음으로 깨달음}
%     \kfVocabItem{2}{사물의 이치를 따져 헤아림}
%   \end{kfVocab}
% ============================================================

% adjustbox로 폭 강제 + 일반 tabular + p{} 열 사용
\newenvironment{kfVocab}%
  {\par\smallskip\needspace{4\baselineskip}%
   \renewcommand{\arraystretch}{1.35}%
   \arrayrulecolor{kfRule}%
   \noindent\begin{adjustbox}{max width=\linewidth, center}%
   \begin{tabular}{|>{\centering\arraybackslash}m{1.0cm}|m{8.5cm}|>{\centering\arraybackslash}m{4.0cm}|}%
   \hline
   \rowcolor{kfPrimaryLight!50}%
   \multicolumn{1}{|c|}{\textbf{\color{kfPrimary}번호}} &
   \multicolumn{1}{c|}{\textbf{\color{kfPrimary}뜻풀이}} &
   \multicolumn{1}{c|}{\textbf{\color{kfPrimary}어휘 (정답 작성)}} \\\hline
   \ignorespaces}%
  {\end{tabular}%
   \end{adjustbox}%
   \arrayrulecolor{black}%
   \par\smallskip}

\newcommand{\kfVocabItem}[2]{%
  \textbf{#1} & #2 & \rule[-1.0em]{0pt}{2.4em}\\\hline%
}

% ============================================================
% activity-blank.tex
% 빈칸 채우기 활동
% 사용:
%   \begin{kfBlankList}
%     \kfBlankItem 첫 번째 문장에서 (\kfBlank)은(는) ...
%     \kfBlankItem 둘째 문장에서 ...
%   \end{kfBlankList}
% ============================================================

\newenvironment{kfBlankList}%
  {\par\smallskip%
   \begin{enumerate}[leftmargin=20pt, itemsep=4pt, topsep=2pt,
                    label=\textbf{\arabic*.}, labelsep=6pt]}%
  {\end{enumerate}\par\smallskip}

\newcommand{\kfBlankItem}{\item}

% 텍스트 안에 박스형 빈칸 (단어 1개 채우기)
\newcommand{\kfBlankBox}[1][2.4cm]{%
  \tikz[baseline=-0.5ex]{%
    \draw[kfRule, line width=0.4pt] (0,0) -- ++(#1,0);%
  }%
}

% ============================================================
% activity-ox.tex (v15)
% O/X 표기 활동 — 그래픽 디자인
% 사용:
%   \begin{kfOXList}
%     \kfOXItem 글쓴이는 자연을 예찬하고 있다.\kfOX
%     \kfOXItem 1연과 4연은 수미상관 구조이다.\kfOX
%   \end{kfOXList}
%
% \kfOX = 우측에 [O] [X] 두 박스 (학생이 하나 동그라미)
% ============================================================

\newenvironment{kfOXList}%
  {\par\smallskip%
   \begin{enumerate}[leftmargin=20pt, itemsep=5pt, topsep=2pt,
                    label=\textbf{\arabic*.}, labelsep=6pt]}%
  {\end{enumerate}\par\smallskip}

\newcommand{\kfOXItem}{\item}

% v16: O/X 그래픽 박스 키움 (18×14 → 24×20pt) + 영역색 강조
\newcommand{\kfOX}{%
  \hfill\nolinebreak
  \tikz[baseline=-0.9ex]{%
    \node[draw=kfPrimary, fill=kfPrimary!5, line width=0.6pt, rounded corners=3pt,
          minimum width=24pt, minimum height=20pt, inner sep=0pt,
          font=\normalsize\bfseries\color{kfPrimary}] (ob) {O};%
    \node[draw=kfMute, fill=kfMute!5, line width=0.6pt, rounded corners=3pt,
          minimum width=24pt, minimum height=20pt, inner sep=0pt,
          font=\normalsize\bfseries\color{kfMute},
          right=4pt of ob] (xb) {X};%
  }%
}

% ============================================================
% activity-connect.tex
% 좌우 선 긋기 활동 (2단 양식 + 중앙 여백)
% 사용:
%   \begin{kfConnect}
%     \kfPair{사르트르}{실존은 본질에 앞선다}
%     \kfPair{니체}{초인 사상}
%     \kfPair{하이데거}{현존재}
%   \end{kfConnect}
% ============================================================

\newenvironment{kfConnect}%
  {\par\smallskip%
   \renewcommand{\arraystretch}{1.6}%
   \begin{tabular}{@{}>{\raggedleft\arraybackslash}m{4.5cm}
                     @{\hspace{2.5cm}}
                     >{\raggedright\arraybackslash}m{6.5cm}@{}}}%
  {\end{tabular}\par\smallskip}

\newcommand{\kfPair}[2]{%
  \tikz[baseline]{\node[draw=kfRule, fill=kfPrimaryLight!30, rounded corners=2pt,
        inner sep=4pt, font=\small]{#1};} \tikz[baseline]{\node[circle, draw=kfMute, fill=kfPaper, inner sep=1.5pt]{};}
  &
  \tikz[baseline]{\node[circle, draw=kfMute, fill=kfPaper, inner sep=1.5pt]{};} \tikz[baseline]{\node[draw=kfRule, fill=kfNoteBg, rounded corners=2pt,
        inner sep=4pt, font=\small]{#2};} \\
}

% ============================================================
% activity-writing.tex
% 글쓰기 활동 (큰 노트 영역)
% 사용:
%   \kfWriting{독후감}{12}     % 라벨 + 12줄 큰 노트
%   \kfWritingPrompt{프롬프트 텍스트}
% ============================================================

\newcommand{\kfWritingPrompt}[1]{%
  \par\smallskip%
  \begin{tcolorbox}[
    enhanced, breakable=false,
    colback=kfPrimaryLight!30, colframe=kfPrimary,
    boxrule=0.4pt, arc=3pt,
    left=\kfBoxPad, right=\kfBoxPad, top=\kfBoxPad, bottom=\kfBoxPad,
  ]%
  \small\textbf{\color{kfPrimary}글쓰기 안내}\\[2pt]%
  #1
  \end{tcolorbox}\par\smallskip%
}

\newcommand{\kfWriting}[2]{%
  \par\smallskip\needspace{\numexpr#2+2\relax\baselineskip}%
  \begin{tcolorbox}[
    enhanced, breakable=false,
    colback=kfPaper, colframe=kfRule,
    boxrule=0.4pt, arc=2pt,
    left=\kfBoxPad, right=\kfBoxPad, top=\kfBoxPad, bottom=\kfBoxPad,
    title={\footnotesize\bfseries\color{kfMute} #1},
    coltitle=kfMute, colbacktitle=kfPaper,
    attach boxed title to top left={yshift=-1.5mm, xshift=6pt},
    boxed title style={colframe=kfRule, boxrule=0.3pt, arc=1pt}
  ]%
  \foreach \i in {1,...,#2}{%
    \textcolor{kfRule}{\rule{\linewidth}{0.3pt}}\\[7pt]%
  }%
  \end{tcolorbox}\par\smallskip%
}

% ============================================================
% activity-structure.tex
% 구조도 (트리 + 빈칸)
% 사용:
%   \begin{kfStructure}
%     \kfNode{0}{서론}{문제 제기}
%     \kfNode{1}{본론}{(\,\quad\,)}
%     \kfNode{1}{본론}{근거 2}
%     \kfNode{0}{결론}{주장 강화}
%   \end{kfStructure}
% ============================================================

\newenvironment{kfStructure}%
  {\par\smallskip\needspace{6\baselineskip}%
   \begin{tcolorbox}[
     enhanced, breakable=false,
     colback=kfPaper, colframe=kfRule,
     boxrule=0.4pt, arc=2pt,
     left=\kfBoxPad, right=\kfBoxPad, top=\kfBoxPad, bottom=\kfBoxPad,
     title={\footnotesize\bfseries\color{kfMute} 글의 구조},
     coltitle=kfMute, colbacktitle=kfPrimaryLight!40,
     attach boxed title to top left={yshift=-1.5mm, xshift=6pt},
     boxed title style={colframe=kfRule, boxrule=0.3pt, arc=1pt}
   ]%
   \begin{tabbing}
     \hspace{1.4cm}\=\hspace{1.4cm}\=\hspace{8cm}\kill}%
  {\end{tabbing}\end{tcolorbox}\par\smallskip}

% 노드: #1=들여쓰기 단계 (0~3), #2=라벨, #3=내용
\newcommand{\kfNode}[3]{%
  \ifcase#1\relax%
    \>\>\textbf{\color{kfPrimary} ▶ #2}\quad #3\\
  \or
    \>\>\quad\textbf{\color{kfMute} ◦ #2}\quad #3\\
  \or
    \>\>\quad\quad\textbf{\color{kfMute} - #2}\quad #3\\
  \or
    \>\>\quad\quad\quad\textbf{\color{kfMute} \textperiodcentered\ #2}\quad #3\\
  \fi
}

% 빈칸 노드
\newcommand{\kfNodeBlank}[2]{\kfNode{#1}{#2}{(\hspace{2.5cm})}}

% ============================================================
% activity-sentence-reading.tex
% 문장 독해 활동 — 문장 박스 1개 + 그에 묶인 문제 N개를 한 minipage로
% 사용:
%   \begin{kfSentenceUnit}{1}{"바람은 내 귀에 속삭이며..."}
%     \kfSentQ{(1)}{서술어 '속삭이며'의 주어는?}
%     \begin{kfChoices}
%       \kfChoice 바람은
%       ...
%     \end{kfChoices}
%   \end{kfSentenceUnit}
% ============================================================

% kfSentenceBox 매크로는 passage-note.tex에 정의됨

\newenvironment{kfSentenceUnit}[2]%
  {\par\smallskip\needspace{6\baselineskip}%
   \noindent\begin{minipage}{\linewidth}%
   \samepage%
   \kfSentenceBox{문장 #1}{#2}%
   \par\smallskip%
  }%
  {\end{minipage}\par\smallskip}

% 같은 문장에 묶인 소문항 (문장 안내 + 발문)
\newcommand{\kfSentQ}[2]{%
  \par\noindent\textbf{#1}~#2\par\smallskip%
}

% 헤더 없이 문장만 있는 묶음 (sentence_ref가 없는 일반 문항)
\newenvironment{kfSentencelessUnit}%
  {\par\smallskip\noindent\begin{minipage}{\linewidth}\samepage}%
  {\end{minipage}\par\smallskip}

% ============================================================
% chapter-cover.tex (v6)
% 챕터 진입 페이지 (표지) — 하단에 영역 목차 추가
% 사용:
%   \kfChapterCover{1}{근대성과 자아}{Chapter 1}{이번 챕터에서는 ...}
%   \begin{kfChapterAreaList}
%     \kfChapterAreaItem{어휘}{kfAreaVocab}{핵심 어휘 12개}
%     ...
%   \end{kfChapterAreaList}
%
%   영역 목차가 없을 때는 \kfChapterCoverEnd 호출
% ============================================================

\newcommand{\kfChapterCover}[4]{%
  \clearpage%
  \thispagestyle{kfBlank}%
  \kfMarkChapter{Chapter #1. #2}%
  \kfChapterCoverBg%
  % v18.5: 챕터 표지 우상단에 로고
  \begin{tikzpicture}[remember picture, overlay]
    \node[anchor=north east, inner sep=0pt] at ($(current page.north east)+(-18mm,-18mm)$) {\kfLogoMd};
  \end{tikzpicture}%
  \vspace*{20mm}%
  \noindent
  \begin{tikzpicture}[remember picture, overlay]
    % 좌측 액센트 바
    \fill[kfPrimary] (current page.west) ++(12mm,25mm) rectangle ++(5mm, -50mm);
    % 챕터 번호 배지 (이중 원, 약간 작게)
    \node[anchor=north west, inner sep=0pt] at ($(current page.north west)+(24mm,-45mm)$) {%
      \begin{tikzpicture}
        \draw[kfPrimary, line width=1pt] (0,0) circle (10mm);
        \draw[kfPrimary, line width=0.4pt] (0,0) circle (8mm);
        \node[font=\normalsize\bfseries\color{kfPrimary}] at (0,2mm) {CHAPTER};
        \node[font=\LARGE\bfseries\color{kfPrimary}] at (0,-3mm) {#1};
      \end{tikzpicture}%
    };
  \end{tikzpicture}%
  \vspace{42mm}%
  \par\noindent\hspace{32mm}%
  \begin{minipage}{0.65\linewidth}%
    {\footnotesize\color{kfMute} #3}\\[4pt]
    {\LARGE\bfseries\color{kfPrimary} #2}\\[8pt]
    {\color{kfRule}\rule{50mm}{0.5pt}}\\[6pt]
    {\small\color{kfInk} #4}
  \end{minipage}%
  \par\vspace{14mm}%
}

% v6 / v18.5 / v18.6: 챕터 표지의 영역 목차 — 흰색 반투명 박스
% v18.6: 배경 가로선/도형과 안 겹치도록 TikZ overlay로 우하단 절대 위치
\makeatletter
% 영역 항목들을 모아 둘 박스
\newsavebox{\kf@arealistbox}
\newenvironment{kfChapterAreaList}{%
  \begin{lrbox}{\kf@arealistbox}%
  \begin{minipage}{0.62\linewidth}%
    \begin{tcolorbox}[%
      enhanced, colback=white, colframe=white,
      opacityback=0.92, opacityframe=0,
      boxrule=0pt, arc=4pt,
      width=\linewidth,
      top=8pt, bottom=8pt, left=10pt, right=10pt,
    ]%
      {\footnotesize\bfseries\color{kfMute} 이 챕터의 영역}\par\vspace{4pt}%
      {\color{kfRule}\hrule height 0.3pt}\par\vspace{4pt}%
      \begin{itemize}[leftmargin=14pt, itemsep=2pt, topsep=0pt, label={\color{kfPrimary!70}\textbullet}]%
}{%
      \end{itemize}%
    \end{tcolorbox}%
  \end{minipage}%
  \end{lrbox}%
  % v18.7: 배경 상단 가로선 ~ 하단 가로선 사이의 중간 영역에 배치
  % 가로선이 내용을 가리지 않도록 박스 중심이 두 선 사이 중점에 위치
  \begin{tikzpicture}[remember picture, overlay]%
    \node[anchor=center, inner sep=0pt]
      at ($(current page.south)+(0mm,90mm)$)
      {\usebox{\kf@arealistbox}};%
  \end{tikzpicture}%
  \vfill%
  \noindent\hfill\kfSeriesBadge\hspace{12mm}%
  \clearpage%
}
\makeatother

% 영역 항목: \kfChapterAreaItem{영역명}{kfArea색}{부제}
\makeatletter
\newcommand{\kfChapterAreaItem}[3]{%
  \item {\small\bfseries\color{#2} #1}%
  \def\kf@cci@sub{#3}%
  \ifx\kf@cci@sub\@empty\else
    {\small\color{kfMute}\, \textemdash\, #3}%
  \fi
}
\makeatother

% 영역 목록 없이 cover만 (legacy)
\newcommand{\kfChapterCoverEnd}{%
  \vfill%
  \noindent\hfill\kfSeriesBadge\hspace{12mm}%
  \clearpage%
}

% ============================================================
% area-header.tex (v16)
% 영역 시작 표제 — 시리즈별 차별화 라벨 + 큰 영역명 + 부제
% PAGE_BREAK_POLICY.md §1.4
%
% v16 (디자인 고도화 Phase 1):
%   - 시리즈별 라벨 박스 추가:
%     소쉬르 = AREA START (영역색 채움 박스)
%     프레게 = SECTION OPEN (영역색 테두리)
%     러셀   = RUSSELL (영역색 라이트 박스)
%     비트   = WITTGENSTEIN (회색 라이트 박스)
%   - 라벨 위치: 영역명 위 좌측
% ============================================================

\makeatletter

% 시리즈 라벨 텍스트
\newcommand{\kf@serieslabel}{%
  \ifcase\kf@series\relax
    AREA START%
  \or
    AREA START%
  \or
    SECTION OPEN%
  \or
    RUSSELL%
  \or
    WITTGENSTEIN%
  \fi
}

% 시리즈 라벨 박스 스타일 (#1 = 영역색)
\newcommand{\kf@labelbox}[1]{%
  \ifcase\kf@series\relax
    % 0/1 = 소쉬르: 영역색 채움
    \tikz[baseline=-2pt]{\node[fill=#1, text=white,
      rounded corners=3pt, inner xsep=6pt, inner ysep=2pt,
      font=\scriptsize\bfseries]{\kf@serieslabel};}%
  \or
    % 1 = 소쉬르: 영역색 채움
    \tikz[baseline=-2pt]{\node[fill=#1, text=white,
      rounded corners=3pt, inner xsep=6pt, inner ysep=2pt,
      font=\scriptsize\bfseries]{\kf@serieslabel};}%
  \or
    % 2 = 프레게: 영역색 테두리
    \tikz[baseline=-2pt]{\node[draw=#1, line width=0.6pt, text=#1,
      rounded corners=3pt, inner xsep=6pt, inner ysep=2pt,
      font=\scriptsize\bfseries]{\kf@serieslabel};}%
  \or
    % 3 = 러셀: 영역색 라이트 박스
    \tikz[baseline=-2pt]{\node[fill=#1!12, text=#1,
      rounded corners=2pt, inner xsep=5pt, inner ysep=1.5pt,
      font=\scriptsize\bfseries]{\kf@serieslabel};}%
  \or
    % 4 = 비트: 회색 미니 박스
    \tikz[baseline=-2pt]{\node[fill=kfMute!10, text=kfMute,
      rounded corners=2pt, inner xsep=5pt, inner ysep=1.5pt,
      font=\scriptsize\bfseries]{\kf@serieslabel};}%
  \fi
}

% v17: 시리즈+영역별 Canva 배경 자동 매핑
% \kf@hasbg=1로 set되면 \kfAreaStart에서 라벨 박스 출력 생략
\def\kf@areaVocab{어휘}\def\kf@areaGrammar{문법}\def\kf@areaConcept{개념}
\def\kf@areaLit{문학}\def\kf@areaNonlit{비문학}
\newif\ifkf@bgon
% 파일 있으면 배경 적용 + boolean true. 없으면 nothing.
\newcommand{\kfBgSet}[1]{%
  \IfFileExists{#1.pdf}{\kfPageBg{#1}\kf@bgontrue}{}%
}
\newcommand{\kf@trybg}[1]{%
  % v18.4: 영역 배경 PDF 비활성화 — 큰 도형이 본문 영역을 침범해
  % "영역 헤더 후 빈 페이지" 문제를 일으킴. 라벨 박스로만 영역 표시.
  \kf@bgonfalse%
}

% Note: 러셀(rus_*)/비트(wit_*) 배경 PDF 미존재 시 \kfPageBg가 에러 → \IfFileExists로 가드
% (현재는 파일이 모두 존재한다고 가정. 부분 제공 시 cls 수정 필요)

\newcommand{\kfAreaStart}[3]{%
  \clearpage%
  \thispagestyle{fancy}%
  \kfMarkArea{#1}%
  \kf@trybg{#1}%
  \vspace*{1mm}%
  % 배경 없으면 라텍스 라벨 박스 출력 (배경 있으면 일러스트가 라벨 역할)
  \ifkf@bgon
    \vspace*{4mm}%
  \else
    \noindent\kf@labelbox{#2}\par\vspace{4pt}%
  \fi
  % 영역명 + 부제
  \noindent
  \begin{tikzpicture}
    \node[anchor=west, inner sep=0pt] (bar) at (0,0) {\color{#2}\rule{3pt}{22pt}};
    \node[anchor=base west, inner sep=0pt, xshift=8pt, yshift=2pt] (lbl) at (bar.east |- 0,0) {%
      {\LARGE\bfseries\color{#2} #1}%
    };
    \def\kf@areasub{#3}%
    \ifx\kf@areasub\@empty\else
      \node[anchor=base west, inner sep=0pt, xshift=8pt, font=\normalsize\color{kfMute}]
        at (lbl.base east) {\,\textperiodcentered\, #3};%
    \fi
  \end{tikzpicture}%
  \par\vspace{2pt}
  {\color{#2!70}\hrule height 1pt}%
  \par\vspace{1pt}%
  {\color{kfRule}\hrule height 0.3pt}%
  \par\vspace{4pt}%
  \needspace{6\baselineskip}\nopagebreak[4]%
  \global\kf@areaJustOpenedtrue% v18.5: 첫 컨텐츠 needspace 무시
}
\makeatother

% 시리즈/영역 단축 매크로
\newcommand{\kfStartVocab}[1]{\kfAreaStart{어휘}{kfAreaVocab}{#1}}
\newcommand{\kfStartGrammar}[1]{\kfAreaStart{문법}{kfAreaGrammar}{#1}}
\newcommand{\kfStartConcept}[1]{\kfAreaStart{개념}{kfAreaConcept}{#1}}
\newcommand{\kfStartLit}[1]{\kfAreaStart{문학}{kfAreaLiterature}{#1}}
\newcommand{\kfStartNonlit}[1]{\kfAreaStart{비문학}{kfAreaNonlit}{#1}}
\newcommand{\kfStartWeekly}[1]{\kfAreaStart{실력 확인}{kfAreaWeekly}{#1}}

% ============================================================
% section-header.tex
% 섹션 시작 라벨 헤더 — [지문] / [활동] / [문제] 라벨
% 좌측 액센트 바 + 라벨 + 부제 (영역 액센트 색)
%
% 사용:
%   \kfSecHeader{kfAreaLiterature}{지문}{빼앗긴 들에도 봄은 오는가 / 이상화}
%   \kfSecHeader{kfAreaConcept}{활동}{내용 확인}
%   \kfSecHeader{kfAreaNonlit}{문제}{}
%
% 헤더 직후 페이지 분할 금지: \needspace{4\baselineskip} + \nopagebreak
% ============================================================

\providecommand{\kfSecHeader}[3]{%
  \par\needspace{10\baselineskip}\filbreak%
  \vspace{3pt}%
  \noindent
  \begin{tikzpicture}
    \node[anchor=west, inner sep=0pt] (bar) at (0,0) {\color{#1}\rule{3pt}{14pt}};
    \node[anchor=west, inner sep=0pt, xshift=6pt, font=\normalsize\bfseries\color{#1}]
       (lbl) at (bar.east) {[\,#2\,]};
    \node[anchor=west, inner sep=0pt, xshift=4pt, font=\small\color{kfMute}]
       at (lbl.east) {#3};
  \end{tikzpicture}%
  \par\vspace{1pt}%
  {\color{#1!50}\hrule height 0.4pt}%
  \par\vspace{2pt}\nopagebreak%
}

% 영역별 단축 매크로
\newcommand{\kfSecPassageVocab}[1]{\kfSecHeader{kfAreaVocab}{지문}{#1}}
\newcommand{\kfSecPassageGrammar}[1]{\kfSecHeader{kfAreaGrammar}{지문}{#1}}
\newcommand{\kfSecPassageConcept}[1]{\kfSecHeader{kfAreaConcept}{지문}{#1}}
\newcommand{\kfSecPassageLit}[1]{\kfSecHeader{kfAreaLiterature}{지문}{#1}}
\newcommand{\kfSecPassageNonlit}[1]{\kfSecHeader{kfAreaNonlit}{지문}{#1}}
\newcommand{\kfSecPassageWeekly}[1]{\kfSecHeader{kfAreaWeekly}{지문}{#1}}

\newcommand{\kfSecActVocab}[1]{\kfSecHeader{kfAreaVocab}{활동}{#1}}
\newcommand{\kfSecActGrammar}[1]{\kfSecHeader{kfAreaGrammar}{활동}{#1}}
\newcommand{\kfSecActConcept}[1]{\kfSecHeader{kfAreaConcept}{활동}{#1}}
\newcommand{\kfSecActLit}[1]{\kfSecHeader{kfAreaLiterature}{활동}{#1}}
\newcommand{\kfSecActNonlit}[1]{\kfSecHeader{kfAreaNonlit}{활동}{#1}}
\newcommand{\kfSecActWeekly}[1]{\kfSecHeader{kfAreaWeekly}{활동}{#1}}

\newcommand{\kfSecQVocab}[1]{\kfSecHeader{kfAreaVocab}{문제}{#1}}
\newcommand{\kfSecQGrammar}[1]{\kfSecHeader{kfAreaGrammar}{문제}{#1}}
\newcommand{\kfSecQConcept}[1]{\kfSecHeader{kfAreaConcept}{문제}{#1}}
\newcommand{\kfSecQLit}[1]{\kfSecHeader{kfAreaLiterature}{문제}{#1}}
\newcommand{\kfSecQNonlit}[1]{\kfSecHeader{kfAreaNonlit}{문제}{#1}}
\newcommand{\kfSecQWeekly}[1]{\kfSecHeader{kfAreaWeekly}{문제}{#1}}


\begin{document}

\thispagestyle{empty}
\vspace*{40mm}
\begin{center}
{\Huge\bfseries\color{kfPrimary} 프레게2}\\[10pt]
{\Large\color{kfMute} 학생용 교재 — 제 1 권}\\[20pt]
{\small\color{kfMute} (챕터 1 \textendash{} 4)}
\end{center}
\vfill
\hfill\kfSeriesBadge\hspace{15mm}
\clearpage
\kfChapterCover{1}{프레게2 (챕터1)}{Chapter 1}{이번 챕터의 학습 내용을 시작합니다.}

\kfStartVocab{어휘 영역 — 사랑과 사회적 제약 관련 어휘}


\kfSecHeader{kfAreaVocab}{문제}{}

\begin{kfQBlock}{1}{}
\kfQStem{신념}
\kfAnswerNote{답안}{4}
\end{kfQBlock}
\begin{kfQBlock}{2}{}
\kfQStem{희생}
\kfAnswerNote{답안}{4}
\end{kfQBlock}

\kfSecHeader{kfAreaVocab}{활동}{사랑과 사회적 제약 관련 어휘}
\noindent\textit{\small [1-5] <보기>에서 알맞은 단어를 골라 문장을 완성하세요.}\par\medskip
\begin{kfBlankList}
\kfBlankItem 춘향은 죽음 앞에서도 도련님을 향한 사랑이 [　　　　　]에 가도 변하지 않을 것이라 믿었다.
\begin{kfEvidence}보\end{kfEvidence}
\kfBlankItem 조선 시대와 같은 [　　　　　] 사회에서는 신분에 따라 결혼이 자유롭지 못했다.
\begin{kfEvidence}기\end{kfEvidence}
\kfBlankItem 이 시의 주제는 현실의 [　　　　　]을 극복하고 얻어낸 자유와 사랑이다.
\begin{kfEvidence}>\end{kfEvidence}
\kfBlankItem 모든 사람은 차별과 억압으로부터 [　　　　　]되어야 한다.
\begin{kfEvidence} \end{kfEvidence}
\kfBlankItem 두 주인공은 신분 차이 때문에 심각한 [　　　　　]을 겪었다.
\begin{kfEvidence}　\end{kfEvidence}
\end{kfBlankList}

\kfSecHeader{kfAreaVocab}{활동}{사랑과 사회적 제약 관련 어휘}
\noindent\textit{\small [1-5] 단어의 정확한 의미를 찾아 선으로 연결해 보세요.}\par\medskip
\begin{kfConnect}
\kfPair{갈등}{㉠ 굳게 믿어 의심하지 않는 마음}
\kfPair{신념}{㉡ 자신의 목숨, 재산 등을 다른 사람이나 목적을 위해 바침}
\kfPair{희생}{㉢ 조건을 붙여 내용을 제한하거나 자유를 막음}
\kfPair{제약}{㉣ 일이나 사정이 서로 복잡하게 뒤얽혀 화합하지 못하는 상태}
\kfPair{절개}{㉤ 신념을 굽히지 않고 굳게 지키는 꿋꿋한 태도}
\end{kfConnect}

\kfSecHeader{kfAreaVocab}{활동}{사랑과 사회적 제약 관련 어휘}
\noindent\textit{\small [1-5] 뜻풀이에 알맞은 단어를 초성 힌트를 참고해 쓰세요.}\par\medskip
\begin{enumerate}[leftmargin=22pt, itemsep=6pt, topsep=4pt, label=\textbf{\arabic*.}]
\item 굳게 믿어 의심하지 않는 마음 [  ㅅ  ㄴ  ]
\item 작품을 통해 글쓴이가 말하고자 하는 중심 생각 [  ㅈ  ㅈ  ]
\item 신분 제도나 옛 관습을 중요하게 여겼던 전근대적인 사회 제도 [  ㅂ  ㄱ  ]
\item 다른 사람이나 어떤 목적을 위하여 자신의 목숨, 재산, 명예 등을 바침 [  ㅎ  ㅅ  ]
\item 여자가 굳은 마음으로 지키는 깨끗한 절개 [  ㅈ  ㅈ  ]
\end{enumerate}

\kfSecHeader{kfAreaVocab}{문제}{}



\kfStartConcept{개념 영역 — 화자, 청자}


\kfSecHeader{kfAreaConcept}{개념}{화자, 청자}

시를 읽을 때는 “누가 누구에게 말하고 있는가?”를 파악하는 것이 중요합니다.

1. 화자 (말하는 이)   　• 뜻: 시 속에서 이야기를 들려주는 사람입니다. 시인이 자신의 생각과 느낌을 효과적으로 전달하기 위해 내세운 사람입니다.   　• 찾는 법: 시 속에서 “나”, “우리”라는 말이 나오면 화자 자신을 가리킵니다.   　• 특징: 화자가 남자인지 여자인지, 어른인지 아이인지에 따라 시의 말투와 느낌이 달라집니다. (예: 어린아이 화자 → 순수하고 귀여운 말투)  

2. 청자 (듣는 이)   　• 뜻: 화자의 이야기를 듣는 상대방입니다.   　• 찾는 법: 화자가 “당신”, “너”, “그대”라고 부르는 대상입니다.   　• 특징: 사람일 수도 있고, 동물이나 자연(달, 별), 사물일 수도 있습니다. 때로는 듣는 사람이 따로 없이 혼잣말을 하기도 합니다.
\par\medskip\begin{itemize}[leftmargin=18pt, itemsep=2pt, topsep=2pt, label=\textbullet]
\item 화자(나·우리) — '나는 오늘도 길을 걷는다'에서 '나'
\item 청자(너·당신) — '엄마야 누나야 강변 살자'의 '엄마·누나'
\item 독백(청자 없음) — '오늘은 왜 이리 가슴이 떨리는가' 같은 혼잣말
\end{itemize}

\kfSecHeader{kfAreaConcept}{문제}{}

\begin{kfQBlock}{1}{}
\kfQStem{화자에 대한 설명으로 알맞지 \uline{않은} 것은 무엇인가요?}
\begin{kfChoices}
\kfChoice 시 속에서 말하는 사람을 뜻한다.
\kfChoice 화자는 반드시 작가 자신과 똑같은 사람이다.
\kfChoice 화자의 성별이나 나이에 따라 시의 분위기가 달라진다.
\kfChoice 시 속에 '나'라는 말이 없어도 화자는 존재한다.
\end{kfChoices}
\end{kfQBlock}
\begin{kfQBlock}{1}{}
\kfQStem{화자: 내일 소풍을 가야 하는 '초등학생'

청자: 하늘}
\kfAnswerNote{답안}{4}
\end{kfQBlock}
\begin{kfQBlock}{2}{}
\kfQStem{화자: 며칠 동안 물을 못 마셔 목이 말랐던 '나무'

청자: 비(소나기)}
\kfAnswerNote{답안}{4}
\end{kfQBlock}
\begin{kfQBlock}{3}{}
\kfQStem{화자: 비를 피할 곳이 없는 '길고양이'

청자: 지나가는 사람들}
\kfAnswerNote{답안}{4}
\end{kfQBlock}
\par\noindent\textit{\small \{'1': '하고 싶은 말: \_\_\_\_\_\_\_\_\_\_\_\_\_\_\_\_\_\_\_\_\_\_\_\_\_\_\_\_\_\_\_\_\_\_\_\_\_\_\_\_\_\_', '2': '하고 싶은 말: \_\_\_\_\_\_\_\_\_\_\_\_\_\_\_\_\_\_\_\_\_\_\_\_\_\_\_\_\_\_\_\_\_\_\_\_\_\_\_\_\_\_', '3': '하고 싶은 말: \_\_\_\_\_\_\_\_\_\_\_\_\_\_\_\_\_\_\_\_\_\_\_\_\_\_\_\_\_\_\_\_\_\_\_\_\_\_\_\_\_\_'\}}

\kfSecHeader{kfAreaConcept}{활동}{화자, 청자}
\noindent\textit{\small [1-3] 다음 시의 구절을 읽고 화자와 청자를 찾아보세요.}\par\medskip
\begin{kfBlankList}
\kfBlankItem 김소월, 「엄마야 누나야」 “엄마야 누나야 강변 살자. 뜰에는 반짝이는 금모래 빛...”
\kfBlankItem 김소월, 「진달래꽃」 “나 보기가 역겨워 가실 때에는 말없이 고이 보내 드리우리다.”
\kfBlankItem 서정주 「춘향의 마지막 편지」, “안녕히 계세요. / 도련님. / ... / 춘향의 사랑보단 오히려 더 먼 / 딴 나라는 아마 아닐 것입니다.”
\end{kfBlankList}

\kfSecHeader{kfAreaConcept}{활동}{화자, 청자}
\noindent\textit{\small 다음 진술이 맞으면 ○, 틀리면 × 표시하시오.}\par\medskip


\kfStartLit{문학 영역 — 배경 지식 돋보기}


\kfSecHeader{kfAreaLiterature}{지문}{배경 지식 돋보기}
\kfPassageNote{배경 지식 돋보기}{%
이 시를 더 깊이 이해하기 위해 ‘춘향전’의 이야기와 시인의 상상력을 알아봅시다.

1. 춘향이와 몽룡이의 사랑 이야기 고전 소설 ‘춘향전’을 알고 있나요? 성춘향과 이몽룡은 신분을 뛰어넘어 서로 사랑했습니다. 하지만 몽룡이가 서울로 떠나고, 춘향이는 감옥에 갇히는 시련을 겪지요. 이 시는 춘향이가 죽음을 앞두고 몽룡이에게 마지막 편지를 쓴다는 상상으로 지어진 시입니다.

2. 죽어도 사라지지 \uline{않는} 사랑 보통 죽으면 모든 것이 끝난다고 생각합니다. 하지만 이 시에서 춘향이는 죽어서 사라지는 것이 아니라, ‘물’이나 ‘구름’, ‘소나기’처럼 모습만 바꾸어 몽룡이 곁에 영원히 머물겠다고 말합니다.
}


\kfSecHeader{kfAreaLiterature}{지문}{작품}
\kfPassageNote{작품}{%
안녕히 계세요. \\[4pt]
도련님. \\[14pt]
지난 오월 단옷날, 처음 만나던 날 \\[4pt]
우리 둘이서 그늘 밑에 서 있던 \\[4pt]
그 크고 푸르던 나무처럼 \\[4pt]
늘 안녕히 안녕히 계세요. \\[14pt]
저승이 어딘지는 똑똑히 모르지만, \\[4pt]
춘향의 사랑보단 오히려 더 먼 \\[4pt]
딴 나라는 아마 아닐 것입니다. \\[14pt]
아주 깊은 땅속을 검은 물이 되어 흐르거나, \\[4pt]
하늘 높이 구름이 되어 날아다니더라도 \\[4pt]
\uline{㉠그건} 결국 도련님 곁을 떠나지 않는 것입니다. \\[14pt]
더구나 그 구름이 소나기가 되어 퍼부을 때 \\[4pt]
춘향은 틀림없이 \uline{㉡거기} 있을 것입니다.
}


\kfSecHeader{kfAreaLiterature}{활동}{차근차근 내용 확인}
\noindent\textit{\small [1-4] 글의 내용과 일치하면 O, 일치하지 않으면 X를 쓰세요.}\par\medskip
\begin{kfOXList}
\kfOXItem 춘향이는 도련님과 처음 만났던 날을 기억하며 편지를 쓰고 있다. \kfOX
\kfOXItem 춘향이는 저승을 도련님과 아주 멀리 떨어진 '딴 나라'라고 생각한다. \kfOX
\kfOXItem 춘향이는 죽어서라도 도련님 곁을 떠나지 않겠다고 다짐하고 있다. \kfOX
\kfOXItem 이 시는 춘향이가 도련님을 원망하며 쓴 편지이다. \kfOX
\end{kfOXList}

\kfSecHeader{kfAreaLiterature}{활동}{나만의 생각 펼치기}
\noindent\textit{\small 다음 조건에 맞게 글을 써 보세요.}\par\medskip
\kfWritingPrompt{춘향이는 물, 구름, 소나기가 되어 도련님 곁에 있겠다고 했습니다. 만약 여러분이 사랑하는 가족이나 친구 곁에 영원히 머물고 싶다면, 어떤 자연물이 되고 싶은가요?

되고 싶은 것: ( 예: 따뜻한 햇살, 시원한 바람, 튼튼한 의자 등 ) 이유: \_\_\_\_\_\_\_\_\_\_\_\_\_\_\_\_\_\_\_\_\_\_\_\_\_\_\_\_\_\_\_\_\_\_\_\_\_\_\_\_\_\_\_\_\_\_\_\_\_\_\_\_\_\_\_\_\_\_\_\_\_\_\_\_\_\_}
\kfWriting{답안 작성}{8}
\kfWritingPrompt{춘향이의 편지를 받은 도련님(이몽룡)은 어떤 마음이 들었을까요? 춘향이에게 보내는 짧은 답장을 써 보세요.}
\kfWriting{답안 작성}{8}

\kfSecHeader{kfAreaLiterature}{문제}{}

\begin{kfQBlock}{1}{}
\kfQStem{이 시에서 '도련님'이라고 부르는 사람은 누구이며, 이 편지를 쓰고 있는 사람은 누구인가요?}
\begin{kfChoices}
\kfChoice 화자: 이몽룡 / 대상: 춘향
\kfChoice 화자: 변 사또 / 대상: 이몽룡
\kfChoice 화자: 춘향 / 대상: 이몽룡
\kfChoice 화자: 춘향 어머니 / 대상: 춘향
\end{kfChoices}
\end{kfQBlock}
\begin{kfQBlock}{2}{}
\kfQStem{㉠'그건'이 가리키는 구체적인 내용으로 알맞은 것은 무엇인가요?}
\begin{kfChoices}
\kfChoice 춘향이 죽어서 가는 저 세상
\kfChoice 하늘의 구름이나 땅속의 검은 물
\kfChoice 도련님이 춘향이를 잊고 떠나버리는 상황
\kfChoice 지난 오월 단오에 처음 만났던 크고 푸른 나무
\end{kfChoices}
\end{kfQBlock}
\begin{kfQBlock}{3}{}
\kfQStem{㉡'거기'가 가리키는 구체적인 장소나 상황은 무엇인가요?}
\begin{kfChoices}
\kfChoice 깊은 땅속
\kfChoice 아주 먼 딴 나라
\kfChoice 도련님이 계신 곳
\kfChoice 구름 위 하늘나라
\end{kfChoices}
\end{kfQBlock}
\begin{kfQBlock}{1}{}
\kfQStem{춘향이가 도련님에게 “안녕히 계세요”라고 말한 의미로 적절하지 \uline{않은} 것은 무엇인가요?}
\begin{kfChoices}
\kfChoice 이제 사랑이 식었으니 그만 헤어지자는 뜻이다.
\kfChoice 다시는 만나지 말자고 차갑게 이별을 말하는 것이다.
\kfChoice 내가 없더라도 나무처럼 변함없이 잘 지내시라는 뜻이다.
\kfChoice 서울로 가는 길이 위험하니 조심해서 가라는 뜻이다.
\end{kfChoices}
\end{kfQBlock}
\begin{kfQBlock}{2}{}
\kfQStem{이 시에서 '나무'가 의미하는 바로 알맞지 \uline{않은} 것은 무엇인가요?}
\begin{kfChoices}
\kfChoice 도련님과 처음 만났던 행복한 추억이 깃든 존재
\kfChoice 춘향이가 도련님을 기다리다 지쳐서 변해버린 모습
\kfChoice 도련님이 춘향이에게 선물로 주었던 아주 비싼 물건
\kfChoice 비가 오면 피할 수 있는 곳으로 도련님의 보호를 뜻함
\end{kfChoices}
\end{kfQBlock}
\begin{kfQBlock}{3}{}
\kfQStem{춘향이가 '검은 물'이나 '구름'이 되겠다고 한 이유는 무엇인가요?}
\begin{kfChoices}
\kfChoice 도련님의 눈을 피해 아주 멀리 도망가기 위해서
\kfChoice 자유롭게 세상을 떠돌아다니며 여행을 하기 위해서
\kfChoice 모양을 바꾸어서라도 도련님 곁에 머물기 위해서
\kfChoice 도련님이 자신을 알아보지 못하게 변장하기 위해서
\end{kfChoices}
\end{kfQBlock}
\begin{kfQBlock}{4}{}
\kfQStem{“그 구름이 소나기가 되어 퍼부을 때 / 춘향은 틀림없이 거기 있을 것입니다”에 담긴 춘향의 마음은?}
\begin{kfChoices}
\kfChoice 비가 오는 날에는 위험하니 밖으로 나가지 말라는 당부
\kfChoice 어떤 모습으로 변해서라도 도련님과 함께하겠다는 의지
\kfChoice 나중에 비가 되어 도련님의 옷을 적시겠다는 복수심
\kfChoice 날씨가 안 좋을 때는 춘향이를 생각해 달라는 부탁
\end{kfChoices}
\end{kfQBlock}
\begin{kfQBlock}{5}{}
\kfQStem{이 시의 주제로 알맞지 \uline{않은} 것은 무엇인가요?}
\begin{kfChoices}
\kfChoice 죽음조차 갈라놓을 수 없는 영원한 사랑
\kfChoice 헤어진 연인에 대한 그리움과 원망
\kfChoice 아름다운 자연 속에서 느끼는 즐거움
\kfChoice 자신의 운명을 거부하고 싶은 슬픈 마음
\end{kfChoices}
\end{kfQBlock}
\begin{kfQBlock}{1}{}
\kfQStem{시간의 흐름과 춘향이의 태도 변화를 중심으로 시의 내용을 정리해 봅시다. 빈칸에 들어갈 말을 <보기>에서 찾아 쓰세요.

지난 5월 단옷날, 처음 만나 그늘 밑에 서 있던 크고 푸른 ( ① 　　　 )처럼 도련님이 늘 평안하시기를 빕니다. 비록 몸은 떠나지만, 도련님께 ( ② 　　　 ) 계시라고 마지막 인사를 올립니다. 저는 죽어서 사라지는 것이 아닙니다. 물이 되고 구름이 되어, 결국 ( ③ 　　　 )가 되어 도련님 곁에 다시 돌아올 것입니다.}
\begin{kfEvidence}나무\end{kfEvidence}
\kfAnswerLines{1}
\end{kfQBlock}


\kfStartNonlit{비문학 영역 — 신분과 능력이 부딪치다, 조선의 신분 제도}


\kfSecHeader{kfAreaNonlit}{지문}{신분과 능력이 부딪치다, 조선의 신분 제도}
\kfPassageNote{신분과 능력이 부딪치다, 조선의 신분 제도}{%
　여러분은 태어날 때부터 직업이 정해져 있다면 어떨까요? 조선 시대에는 태어나면서부터 신분이 정해져 있었고, 평생 그 신분에 따라 살아야 했습니다. \uline{㉠이것}을 '신분 제도'라고 합니다. 조선 시대 사람들은 양반, 중인, 상민, 천민의 네 신분으로 나뉘었고, 이에 따라 할 수 있는 일과 권리가 달랐습니다.

　가장 높은 신분인 '양반'은 나랏일을 하는 관리가 되어 권력과 재산을 가졌습니다. 양반의 자녀는 글 공부를 하여 과거 시험을 볼 수 있었습니다. 그 아래 '중인'은 의원이나 통역관 같은 전문 기술직이었습니다. \uline{㉡이들}은 양반보다 지위는 낮았지만, 기술 덕분에 비교적 넉넉하게 살았습니다.

　백성의 대부분은 '상민'이었습니다. 농사를 짓거나 장사를 하며 나라에 세금을 냈습니다. 상민은 법적으로 자유로웠지만, 양반처럼 높은 벼슬을 하기는 어려웠습니다. 가장 낮은 신분은 '천민'이었습니다. 천민에 속하는 노비는 물건처럼 사고팔 수 있었고, 마음대로 직업을 바꿀 수도 없었습니다.

　조선 시대 신분 제도에는 몇 가지 특징이 있었습니다. 첫째, 신분은 부모에게서 그대로 물려받았습니다. 둘째, 다른 신분끼리는 결혼하기가 매우 어려웠습니다. 셋째, 신분 상승이 힘들었습니다. 상민이 과거 시험에 합격하면 양반이 될 수도 있었지만, \uline{㉢이것}은 극히 드문 일이었습니다.

　이처럼 조선 시대는 능력보다 신분이 중요한 사회였습니다. 오늘날 우리가 누리는 자유와 평등은 과거의 불합리한 차별을 없애고 모든 사람이 존중받는 세상을 만들기 위해 노력해 온 결과입니다.
}


\kfSecHeader{kfAreaNonlit}{활동}{어휘력 쑥쑥! 단어 익히기}
\noindent\textit{\small [1-4] 다음 단어의 뜻을 찾아 선으로 연결해 보세요.}\par\medskip
\begin{kfConnect}
\kfPair{신분}{㉠ 남을 지배하거나 복종시킬 수 있는 힘}
\kfPair{권력}{㉡ 이치에 맞지 않거나 도리에 어긋남}
\kfPair{벼슬}{㉢ 개인의 사회적인 위치나 계급}
\kfPair{불합리}{㉣ 나랏일을 담당하는 직위나 직무}
\end{kfConnect}

\kfSecHeader{kfAreaNonlit}{활동}{차근차근 내용 확인}
\noindent\textit{\small [1-4] 글의 내용과 일치하면 O, 일치하지 않으면 X를 쓰세요.}\par\medskip
\begin{kfOXList}
\kfOXItem 조선 시대 사람들은 자신이 원하면 언제든지 직업을 바꿀 수 있었다. \kfOX
\kfOXItem 상민은 인구의 대부분을 차지했으며 나라에 세금을 냈다. \kfOX
\kfOXItem 천민인 노비는 주인의 재산으로 취급되어 사고팔 수 있었다. \kfOX
\kfOXItem 조선 시대는 능력보다 신분이 더 중요한 사회였다. \kfOX
\end{kfOXList}

\kfSecHeader{kfAreaNonlit}{활동}{나만의 생각 펼치기}
\noindent\textit{\small 다음 조건에 맞게 글을 써 보세요.}\par\medskip
\kfWritingPrompt{여러분이 조선 시대에 살고 있는 깨어 있는 지식인이라고 상상해 보세요. 임금님께 “신분 제도를 없애고 능력에 따라 사람을 뽑아야 합니다”라고 주장하는 짧은 건의문을 써 보세요.}
\kfWriting{답안 작성}{8}
\kfWritingPrompt{조선 시대 양반 선비가 타임머신을 타고 2025년의 대한민국에 왔습니다. “아니, 어떻게 감히 어린아이가 어른과 겸상을 하고, 누구나 사장님이 될 수 있단 말이오?”라고 놀라는 선비에게, 오늘날의 '평등'에 대해 설명해 주세요.}
\kfWriting{답안 작성}{8}

\kfSecHeader{kfAreaNonlit}{문제}{}

\begin{kfQBlock}{1}{}
\kfQStem{㉠ '이것'이 가리키는 내용으로 가장 적절한 것은 무엇인가요?}
\begin{kfChoices}
\kfChoice 직업을 마음대로 바꿀 수 있는 자유
\kfChoice 태어날 때부터 정해진 사회적 위치와 규칙
\kfChoice 누구나 노력하면 왕이 될 수 있는 기회
\kfChoice 돈을 많이 벌면 양반이 되는 제도
\end{kfChoices}
\end{kfQBlock}
\begin{kfQBlock}{2}{}
\kfQStem{㉡ '이들'은 어떤 신분의 사람들을 말하나요?}
\begin{kfChoices}
\kfChoice 양반
\kfChoice 중인
\kfChoice 상민
\kfChoice 천민
\end{kfChoices}
\end{kfQBlock}
\begin{kfQBlock}{3}{}
\kfQStem{㉢ '이것'이 가리키는 상황은 무엇인가요?}
\begin{kfChoices}
\kfChoice 상민이 농사를 짓거나 장사를 하는 것
\kfChoice 상민이 세금을 내지 않고 도망가는 것
\kfChoice 상민이 신분을 뛰어넘어 양반이 되는 것
\kfChoice 상민이 천민과 결혼하여 사는 것
\end{kfChoices}
\end{kfQBlock}
\begin{kfQBlock}{1}{}
\kfQStem{조선 시대 '중인'에 대한 설명으로 알맞은 것은 무엇인가요?}
\begin{kfChoices}
\kfChoice 가장 높은 신분으로 나랏일을 결정했다.
\kfChoice 농사를 짓는 백성의 대부분을 차지했다.
\kfChoice 전문 기술을 가지고 있어 생활이 비교적 넉넉했다.
\kfChoice 물건처럼 사고팔릴 수 있는 가장 낮은 신분이었다.
\end{kfChoices}
\end{kfQBlock}
\begin{kfQBlock}{2}{}
\kfQStem{이 글을 통해 알 수 있는 조선 시대 신분 제도의 특징이 \uline{아닌} 것은?}
\begin{kfChoices}
\kfChoice 부모님의 신분이 자녀에게 그대로 이어진다.
\kfChoice 신분이 다른 사람끼리 자유롭게 결혼할 수 있었다.
\kfChoice 태어날 때부터 신분에 따라 할 수 있는 일이 정해졌다.
\kfChoice 상민이 양반이 되는 것은 불가능하진 않지만 매우 어려웠다.
\end{kfChoices}
\end{kfQBlock}
\begin{kfQBlock}{3}{}
\kfQStem{'상민'에 대한 설명으로 옳은 것은 무엇인가요?}
\begin{kfChoices}
\kfChoice 의원이나 통역관 같은 전문직으로 일했다.
\kfChoice 신분이 가장 낮아 주인의 재산으로 취급받았다.
\kfChoice 과거 시험을 통해 관리가 되는 일이 매우 흔했다.
\kfChoice 법적으로는 자유로웠지만 높은 벼슬을 하기는 어려웠다.
\end{kfChoices}
\end{kfQBlock}
\begin{kfQBlock}{4}{}
\kfQStem{조선 시대에 '양반'이 누렸던 특권으로 알맞지 \uline{않은} 것은 무엇인가요?}
\begin{kfChoices}
\kfChoice 나라에 세금을 가장 많이 내야 했다.
\kfChoice 나랏일을 하는 관리가 되어 권력을 가졌다.
\kfChoice 다른 사람의 물건을 마음대로 빼앗을 수 있었다.
\kfChoice 힘든 농사일을 도맡아 하며 백성을 섬겼다.
\end{kfChoices}
\end{kfQBlock}
\begin{kfQBlock}{5}{}
\kfQStem{오늘날의 사회가 조선 시대와 다른점으로 가장 적절한 것은 무엇인가요?}
\begin{kfChoices}
\kfChoice 태어날 때부터 평생 직업이 정해져 있다.
\kfChoice 부모님의 직업을 무조건 물려받아야 한다.
\kfChoice 신분 차별이 없고 모든 사람이 법적으로 평등하다.
\kfChoice 능력이 없어도 신분이 높으면 존경받는다.
\end{kfChoices}
\end{kfQBlock}
\begin{kfQBlock}{1}{}
\kfQStem{조선 시대의 신분 계급과 그 특징을 표로 정리해 봅시다. 빈칸에 들어갈 알맞은 말을 쓰세요.

조선의 네 가지 신분

| 신분 이름 | 하는 일 및 특징 | | :--- | :--- | | ( ① 　　　 ) | 나랏일을 하는 관리, 권력과 재산을 가짐. | | 중인 | 의원, 통역관 같은 ( ② 　　　 ) 기술직. | | ( ③ 　　　 ) | 백성의 대부분. 농사나 장사를 하며 세금을 냄. | | 천민(노비) | 가장 낮은 신분. 물건처럼 사고팔리기도 함. |

신분 제도의 특징

| 특징 | 내용 | | :--- | :--- | | 세습 | 신분은 ( ④ 　　　 )에게서 그대로 물려받는다. | | 결혼 | 다른 신분끼리는 결혼하기가 매우 어렵다. | | 변화 | 한번 정해진 신분은 바꾸기가 ( ⑤ 쉽다 / 어렵다 ). |}
\kfAnswerLines{1}
\end{kfQBlock}


\kfStartWeekly{실력확인 영역 — 주간 실력 확인}


\noindent\textit{\small (제한시간: 30분 / 총 20문항)}\par\medskip
\par\noindent\textbf{[6-10] 다음 시를 읽고, 물음에 답하세요.}\par\medskip
\kfPassageNote{지문 6}{%
안녕히 계세요. \\[4pt]
도련님. \\[14pt]
지난 오월 단옷날, 처음 만나던 날 \\[4pt]
우리 둘이서 그늘 밑에 서 있던 \\[4pt]
\uline{㉠그 크고 푸르던 나무}처럼 \\[4pt]
늘 안녕히 안녕히 계세요. \\[14pt]
저승이 어딘지는 똑똑히 모르지만, \\[4pt]
\uline{㉡춘향의 사랑}보단 오히려 더 먼 \\[4pt]
딴 나라는 아마 아닐 것입니다. \\[14pt]
아주 깊은 땅속을 \uline{㉢검은 물}이 되어 흐르거나, \\[4pt]
하늘 높이 구름이 되어 날아다니더라도 \\[4pt]
그건 결국 도련님 곁을 떠나지 않는 것입니다. \\[14pt]
더구나 \uline{㉣그 구름이 소나기가 되어 퍼부을 때} \\[4pt]
춘향은 틀림없이 거기 있을 것입니다.
}
\par\noindent\textbf{[11-15] 다음 글을 읽고, 물음에 답하세요.}\par\medskip
\kfPassageNote{지문 11}{%
여러분은 태어날 때부터 직업이 정해져 있다면 어떨까요? 조선 시대에는 태어나면서부터 신분이 정해져 있었고, 평생 그 신분에 따라 살아야 했습니다. 이것을 '신분 제도'라고 합니다. 조선 시대 사람들은 양반, 중인, 상민, 천민의 네 신분으로 나뉘었고, 이에 따라 할 수 있는 일과 권리가 달랐습니다. 가장 높은 신분인 '양반'은 나랏일을 하는 관리가 되어 권력과 재산을 가졌습니다. 양반의 자녀는 글 공부를 하여 과거 시험을 볼 수 있었습니다. 그 아래 '중인'은 의원이나 통역관 같은 전문 기술직이었습니다. 이들은 양반보다 지위는 낮았지만, 기술 덕분에 비교적 넉넉하게 살았습니다. 백성의 대부분은 '상민'이었습니다. 농사를 짓거나 장사를 하며 나라에 세금을 냈습니다. 상민은 법적으로 자유로웠지만, 양반처럼 높은 벼슬을 하기는 어려웠습니다. 가장 낮은 신분은 '천민'이었습니다. 천민에 속하는 노비는 물건처럼 사고팔 수 있었고, 마음대로 직업을 바꿀 수도 없었습니다. 조선 시대 신분 제도에는 몇 가지 특징이 있었습니다. 첫째, 신분은 부모에게서 그대로 물려받았습니다. 둘째, 다른 신분끼리는 결혼하기가 매우 어려웠습니다. 셋째, 신분 상승이 힘들었습니다. 상민이 과거 시험에 합격하면 양반이 될 수도 있었지만, 이것은 극히 드문 일이었습니다. 이처럼 조선 시대는 능력보다 신분이 중요한 사회였습니다. 오늘날 우리가 누리는 자유와 평등은 과거의 불합리한 차별을 없애고 모든 사람이 존중받는 세상을 만들기 위해 노력해 온 결과입니다.
}
\begin{kfQBlock}{1}{}
\kfQStem{다음 <보기>의 상황에 가장 어울리는 단어는 무엇인가요?}
\begin{kfEvidence}<보기> 　독립운동가들은 나라를 되찾기 위해 자신의 재산과 목숨까지 기꺼이 바쳤습니다. 그들의 숭고한 [　　　　　] 덕분에 우리는 자유로운 세상에서 살 수 있게 되었습니다.\end{kfEvidence}
\begin{kfChoices}
\kfChoice 신념
\kfChoice 희생
\kfChoice 절개
\kfChoice 해방
\end{kfChoices}
\end{kfQBlock}

\begin{kfQBlock}{2}{}
\kfQStem{다음 문장의 빈칸에 들어갈 낱말로 가장 적절한 것은 무엇인가요?}
\begin{kfEvidence}“작가들은 소설이나 영화를 통해 독자들에게 전하고 싶은 중심 생각, 즉 [　　　　　]을/를 이야기 속에 담아냅니다.”\end{kfEvidence}
\begin{kfChoices}
\kfChoice 신념
\kfChoice 절개
\kfChoice 주제
\kfChoice 제약
\end{kfChoices}
\end{kfQBlock}

\begin{kfQBlock}{3}{}
\kfQStem{'봉건'이라는 단어와 가장 관련 깊은 시대적 상황은 무엇인가요?}
\begin{kfChoices}
\kfChoice 누구나 자유롭게 직업을 선택할 수 있는 평등한 사회
\kfChoice 신분 제도나 옛날의 관습을 중요하게 여기는 옛날 사회
\kfChoice 과학 기술이 발달하여 기계가 사람을 대신하는 미래 사회
\kfChoice 다른 나라의 문화를 받아들여 새롭게 변화하는 개방 사회
\end{kfChoices}
\end{kfQBlock}

\begin{kfQBlock}{4}{}
\kfQStem{춘향이가 지키려고 했던 '절개'의 의미를 현대적으로 가장 잘 해석한 것은?}
\begin{kfChoices}
\kfChoice 자신의 뜻을 굽혀 강한 사람의 의견을 따르는 유연한 태도
\kfChoice 한번 맺은 약속이나 믿음을 끝까지 지키는 마음
\kfChoice 시대의 흐름에 맞춰 약속을 고쳐 가는 융통성 있는 자세
\kfChoice 자신의 이익을 위해 신념을 잠시 잠시 바꿀 줄 아는 지혜
\end{kfChoices}
\end{kfQBlock}

\begin{kfQBlock}{5}{}
\kfQStem{다음 중 단어와 그 뜻풀이가 바르게 연결된 것은 무엇인가요?}
\begin{kfChoices}
\kfChoice 해방 - 어떤 일에 얽매여 꼼짝 못 하게 됨
\kfChoice 제약 - 조건을 붙여 자유를 막거나 제한함
\kfChoice 저승 - 사람들이 살고 있는 현실의 세상
\kfChoice 갈등 - 서로 마음이 잘 맞아 평화로운 상태
\end{kfChoices}
\end{kfQBlock}

\begin{kfQBlock}{6}{}
\kfQStem{이 시의 표현상 특징으로 가장 알맞은 것은 무엇인가요?}
\begin{kfChoices}
\kfChoice 소리를 흉내 내는 말을 사용하여 생동감을 주었다.
\kfChoice 말하는 이가 듣는 이에게 편지를 쓰듯이 공손한 말투를 사용했다.
\kfChoice 묻고 답하는 형식을 사용하여 두 사람이 대화하는 것처럼 표현했다.
\kfChoice 사투리를 사용하여 시골의 정겨운 느낌을 살려 표현했다.
\end{kfChoices}
\end{kfQBlock}

\begin{kfQBlock}{7}{}
\kfQStem{이 시의 화자인 춘향이 '도련님'에게 편지를 쓰는 상황으로 가장 적절한 것은?}
\begin{kfChoices}
\kfChoice 도련님과 함께 서울로 떠나게 되어 기쁜 마음으로 쓰는 편지
\kfChoice 감옥에 갇혀 죽음을 앞두고 도련님께 마지막으로 쓰는 편지
\kfChoice 변 사또와의 결혼 소식을 알리며 미안한 마음으로 쓰는 편지
\kfChoice 도련님이 과거 시험에 합격한 것을 축하하기 위해 쓰는 편지
\end{kfChoices}
\end{kfQBlock}

\begin{kfQBlock}{8}{}
\kfQStem{㉠\textasciitilde{}㉣에 대한 설명으로 적절하지 \uline{않은} 것은 무엇인가요?}
\begin{kfChoices}
\kfChoice ㉠: 도련님과 처음 만났던 행복한 추억을 떠올리며, 도련님이 나무처럼 변함없이 건강하게 지내기를 바라는 마음을 담고 있다.
\kfChoice ㉡: 춘향의 사랑이 아주 깊고 커서, 죽어서 가는 저승조차도 도련님을 향한 사랑의 거리보다는 가깝게 느껴진다는 의미이다.
\kfChoice ㉢: 죽은 뒤에는 도련님에게 얽매이지 않고, 물과 구름처럼 자유롭게 세상을 떠돌아다니며 여행하고 싶다는 뜻이다.
\kfChoice ㉣: 비록 몸은 사라지고 모습은 변했지만, 비가 되어 내리는 순간만큼은 도련님과 가장 가까이에서 만나고 싶다는 간절한 마음이다.
\end{kfChoices}
\end{kfQBlock}

\begin{kfQBlock}{9}{}
\kfQStem{춘향이가 '검은 물'이나 '구름'이 되겠다고 한 이유는 무엇인가요?}
\begin{kfChoices}
\kfChoice 도련님의 눈을 피해 아주 멀리 도망가기 위해서
\kfChoice 자유롭게 세상을 떠돌아다니며 여행을 하기 위해서
\kfChoice 도련님이 자신을 알아보지 못하게 변장하기 위해서
\kfChoice 죽어서 모양을 바꾸어서라도 도련님 곁에 머물기 위해서
\end{kfChoices}
\end{kfQBlock}

\begin{kfQBlock}{10}{}
\kfQStem{<보기>를 읽고, 이 시에서 말하는 '사랑'의 특징을 가장 잘 설명한 것은 무엇인가요?}
\begin{kfEvidence}<보기> 　보통 사람들은 죽음을 모든 것이 끝나는 '이별'이라고 생각합니다. 하지만 문학 작품에서는 죽음을 넘어선 영원한 사랑을 노래하기도 합니다. 몸은 사라져도 마음은 자연물(비, 바람, 나무 등)로 변하여 사랑하는 사람 곁에 남는다는 믿음을 보여줍니다.\end{kfEvidence}
\begin{kfChoices}
\kfChoice 죽음은 피할 수 없으므로, 사랑하는 사람과 헤어지는 슬픔을 받아들이고 포기하고 있다.
\kfChoice 죽어서 저승으로 가더라도 도련님을 잊지 않기 위해, 처음 만난 날의 추억을 계속 떠올리고 있다.
\kfChoice 죽음은 끝이 아니라 새로운 시작이며, 모습만 바뀔 뿐 도련님 곁을 떠나지 않겠다는 마음을 보여준다.
\kfChoice 자연물로 변해서라도 곁에 있겠다는 것은 거짓말이며, 사실은 더 살고 싶은 마음을 표현한 것이다.
\end{kfChoices}
\end{kfQBlock}

\begin{kfQBlock}{11}{}
\kfQStem{조선 시대 신분 제도의 특징으로 알맞지 \uline{않은} 것은 무엇인가요?}
\begin{kfChoices}
\kfChoice 태어날 때부터 신분이 정해져 있어 평생 따라다녔다.
\kfChoice 부모님의 신분이 자녀에게 그대로 이어지는 사회였다.
\kfChoice 능력이 뛰어나면 누구나 쉽게 양반이 되어 벼슬을 할 수 있었다.
\kfChoice 신분이 다른 사람끼리는 결혼하는 것이 매우 어려웠다.
\end{kfChoices}
\end{kfQBlock}

\begin{kfQBlock}{12}{}
\kfQStem{다음 중 '중인' 신분에 대한 설명으로 옳은 것은 무엇인가요?}
\begin{kfChoices}
\kfChoice 가장 낮은 신분으로 주인의 재산처럼 사고팔렸다.
\kfChoice 농사나 장사를 하며 나라에 세금을 내는 평민이었다.
\kfChoice 의원이나 통역관 같은 전문 기술을 가진 사람들이었다.
\kfChoice 나랏일을 결정하는 높은 벼슬아치들이 대부분이었다.
\end{kfChoices}
\end{kfQBlock}

\begin{kfQBlock}{13}{}
\kfQStem{다음 <보기>와 같은 상황을 겪은 조선 시대 '중인'의 심정으로 가장 적절한 것은?}
\begin{kfEvidence}<보기> 　“나는 의술 실력이 뛰어나서 죽어가는 사람도 살려냈다. 하지만 신분이 낮다는 이유로 나라의 중요한 정책을 결정하는 높은 자리에 오를 수는 없었다.”\end{kfEvidence}
\begin{kfChoices}
\kfChoice “내 실력이 부족하니 더 열심히 공부해야겠어.”
\kfChoice “신분이 낮아도 돈을 많이 벌 수 있으니 만족해.”
\kfChoice “능력이 있어도 신분 때문에 뜻을 펼칠 수 없어 억울해.”
\kfChoice “양반님들이 나랏일을 잘하시니 나는 내 일만 하면 돼.”
\end{kfChoices}
\end{kfQBlock}

\begin{kfQBlock}{14}{}
\kfQStem{다음 <보기>의 상황에서 조선 시대 양반이 했을 법한 생각으로 가장 적절한 것은?}
\begin{kfEvidence}<보기> 　“상민 출신인 김 씨가 장사를 해서 큰돈을 벌어 기와집을 짓고 잘 산다는 소문이 들린다.”\end{kfEvidence}
\begin{kfChoices}
\kfChoice 김 씨를 우리 마을의 사또로 추천해야겠다.
\kfChoice 나도 장사를 배워서 김 씨처럼 부자가 되어야겠다.
\kfChoice 능력이 있으면 신분에 상관없이 성공하는 것이 당연하지.
\kfChoice 감히 상민이 양반 흉내를 내다니, 신분 질서가 무너지는구나.
\end{kfChoices}
\end{kfQBlock}

\begin{kfQBlock}{15}{}
\kfQStem{오늘날 우리 사회가 조선 시대와 다른 점으로 가장 적절한 것은 무엇인가요?}
\begin{kfChoices}
\kfChoice 모든 사람은 법 앞에 평등하며 신분에 따른 차별이 없다.
\kfChoice 양반 자녀만 학교에서 글공부를 할 수 있는 자격을 받는다.
\kfChoice 의원이나 통역관 같은 전문직은 양반의 추천을 받아야 한다.
\kfChoice 결혼은 같은 신분의 사람끼리만 할 수 있도록 정해져 있다.
\end{kfChoices}
\end{kfQBlock}

\begin{kfQBlock}{16}{}
\kfQStem{시 속에서 말하는 사람인 '화자'에 대한 설명으로 옳은 것은 무엇인가요?}
\begin{kfChoices}
\kfChoice 화자는 반드시 시를 쓴 작가 자신과 똑같은 사람이다.
\kfChoice 시인은 자신의 생각을 효과적으로 전달하기 위해 화자를 설정한다.
\kfChoice 화자는 항상 시의 겉면에 드러나서 “나”라고 말해야만 한다.
\kfChoice 화자의 성별이나 나이는 시의 분위기와 아무런 관련이 없다.
\end{kfChoices}
\end{kfQBlock}

\begin{kfQBlock}{17}{}
\kfQStem{다음 시 구절에서 화자가 누구인지를 짐작할 수 있는 단어는 무엇인가요?}
\begin{kfEvidence}“오빠 생각 / 뜸북 뜸북 뜸북새 논에서 울고 / 뻐꾹 뻐꾹 뻐꾹새 숲에서 울 제 / 우리 오빠 말 타고 서울 가시면...”\end{kfEvidence}
\begin{kfChoices}
\kfChoice 서울
\kfChoice 오빠
\kfChoice 뜸북새
\kfChoice 뻐꾹새
\end{kfChoices}
\end{kfQBlock}

\begin{kfQBlock}{18}{}
\kfQStem{다음 시 구절에서 화자와 청자가 바르게 짝지어진 것은 무엇인가요?}
\begin{kfEvidence}“엄마야 누나야 강변 살자. 뜰에는 반짝이는 금모래 빛...”\end{kfEvidence}
\begin{kfChoices}
\kfChoice 화자: 엄마 / 청자: 누나
\kfChoice 화자: 강변 / 청자: 금모래
\kfChoice 화자: 시인 / 청자: 독자
\kfChoice 화자: 남자아이 / 청자: 엄마, 누나
\end{kfChoices}
\end{kfQBlock}

\begin{kfQBlock}{19}{}
\kfQStem{화자가 누구냐에 따라 시의 느낌이 달라집니다. 다음 중 짝이 \uline{어색한} 것은?}
\begin{kfChoices}
\kfChoice 화자: 어린아이 - 순수하고 부드러운 말투
\kfChoice 화자: 독립운동가 - 강하고 의지적인 말투
\kfChoice 화자: 사랑에 빠진 사람 - 달콤하고 부드러운 말투
\kfChoice 화자: 장난꾸러기 소년 - 점잖고 예의 바른 말투
\end{kfChoices}
\end{kfQBlock}

\begin{kfQBlock}{20}{}
\kfQStem{시 속의 화자를 '어린아이'로 설정했을 때 얻을 수 있는 효과는 무엇인가요?}
\begin{kfChoices}
\kfChoice 순수하고 솔직한 마음을 독자에게 더 효과적으로 전달할 수 있다.
\kfChoice 어려운 한자어를 사용하여 지식을 뽐내기에 아주 유리하다.
\kfChoice 어른들은 시를 읽지 못하게 하고 오직 어린이만 읽게 할 수 있다.
\kfChoice 화자가 누군지 숨겨서 독자가 추리하는 재미를 느낄 수 있게 한다.
\end{kfChoices}
\end{kfQBlock}


\kfSecHeader{kfAreaWeekly}{문제}{}


\kfSecHeader{kfAreaWeekly}{활동}{실력 점검: 주간 종합 평가}
\noindent\textit{\small 다음 빈칸에 알맞은 말을 골라 채우시오.}\par\medskip



\clearpage

\kfChapterCover{2}{프레게2 (챕터2)}{Chapter 2}{이번 챕터의 학습 내용을 시작합니다.}

\kfStartVocab{어휘 영역 — 민주주의와 정치 관련 어휘}


\kfSecHeader{kfAreaVocab}{문제}{}

\begin{kfQBlock}{1}{}
\kfQStem{선거}
\kfAnswerNote{답안}{4}
\end{kfQBlock}
\begin{kfQBlock}{2}{}
\kfQStem{정권}
\kfAnswerNote{답안}{4}
\end{kfQBlock}

\kfSecHeader{kfAreaVocab}{활동}{민주주의와 정치 관련 어휘}
\noindent\textit{\small [1-5] <보기>에서 알맞은 단어를 골라 문장을 완성하세요.}\par\medskip
\begin{kfBlankList}
\kfBlankItem 대한민국 헌법 제1조에 따르면 대한민국의 [　　　　　]은 국민에게 있다.
\begin{kfEvidence}보\end{kfEvidence}
\kfBlankItem 우리 학교 전교 회장 [　　　　　]에 출마한 후보들이 열띤 토론을 벌였다.
\begin{kfEvidence}기\end{kfEvidence}
\kfBlankItem 도서관이나 공원은 여러 사람이 함께 이용하는 [　　　　　] 장소이므로 예절을 지켜야 한다.
\begin{kfEvidence}>\end{kfEvidence}
\kfBlankItem 4·19 혁명 당시 [　　　　　] 들은 힘을 합쳐 독재 정권을 무너뜨렸다.
\begin{kfEvidence} \end{kfEvidence}
\kfBlankItem 독재 [　　　　　] 이/가 힘을 잃자 국민들은 자유를 되찾기 위해 거리로 나왔다.
\begin{kfEvidence}　\end{kfEvidence}
\end{kfBlankList}

\kfSecHeader{kfAreaVocab}{활동}{민주주의와 정치 관련 어휘}
\noindent\textit{\small [1-5] 단어의 정확한 의미를 찾아 선으로 연결해 보세요.}\par\medskip
\begin{kfConnect}
\kfPair{혁명}{㉠ 남을 따르게 하거나 지배할 수 있는 힘}
\kfPair{대의}{㉡ 기존의 사회 체제를 새롭게 바꾸기 위해 일어나는 급격한 변화}
\kfPair{권력}{㉢ 대표를 뽑아 그들이 정치를 의논하거나 결정하게 함}
\kfPair{독재}{㉣ 특정한 개인이나 집단이 권력을 독차지하여 마음대로 정치를 하는 것}
\kfPair{시민}{㉤ 민주 사회의 구성원으로서 권리와 의무를 가진 사람}
\end{kfConnect}

\kfSecHeader{kfAreaVocab}{활동}{민주주의와 정치 관련 어휘}
\noindent\textit{\small [1-5] 뜻풀이에 알맞은 단어를 초성 힌트를 참고해 쓰세요.}\par\medskip
\begin{enumerate}[leftmargin=22pt, itemsep=6pt, topsep=4pt, label=\textbf{\arabic*.}]
\item 국가의 주권이 국민에게 있고, 국민을 위해 정치를 하는 제도: [  ㅁ  ㅈ  ㅈ  ㅇ  ]
\item 기존의 사회 체제를 새롭게 바꾸기 위해 일어나는 급격한 변화: [  ㅎ  ㅁ  ]
\item 특정한 개인이나 집단이 권력을 독차지하여 마음대로 정치를 하는 것: [  ㄷ  ㅈ  ]
\item 국민이 직접 결정하는 대신, 대표를 뽑아 그들이 정치를 의논하거나 결정하게 함: [  ㄷ  ㅇ  ]
\item 개인의 이익이 아닌, 국가나 사회의 구성원에게 두루 관계되는 것: [  ㄱ  ㄱ  ]
\end{enumerate}

\kfSecHeader{kfAreaVocab}{문제}{}



\kfStartConcept{개념 영역 — 비유하기}


\kfSecHeader{kfAreaConcept}{개념}{비유하기}

비유가 무엇인지, 그리고 비유를 이루는 요소에는 무엇이 있는지 꼼꼼히 읽어 보세요.

1. 비유란 무엇일까요? 　어떤 현상이나 사물을 직접 설명하지 않고, 비슷한 다른 사물에 빗대어 표현하는 방법을 말해요.

2. 비유를 이루는 두 가지 요소: 원관념과 보조 관념 　비유가 성립하려면 두 가지가 꼭 필요해요! 　• 원관념 (A): 내가 원래 표현하고자 하는 대상 　• 보조 관념 (B): 원관념을 더 잘 표현하기 위해 빗대어 가져온 대상

[예시] “사과 같은 내 얼굴” 　• 원관념: 내 얼굴 (내가 진짜 말하고 싶은 것) 　• 보조 관념: 사과 (얼굴을 더 잘 표현하기 위해 가져온 것)

3. 비유의 세 가지 종류 　• 직유법: '\textasciitilde{}같이', '\textasciitilde{}처럼', '\textasciitilde{}듯이' 등을 사용하여 원관념과 보조 관념을 직접 연결해요. (예: 호수 같은 내 마음) 　• 은유법: 연결어 없이 'A는 B다'의 형태로 원관념과 보조 관념을 은근하게 연결해요. (예: 내 마음은 호수요) 　• 의인법: 사람이 \uline{아닌} 것(원관념)을 사람(보조 관념)처럼 생각하고 행동하게 표현해요. (예: 꽃이 웃는다)
\par\medskip\begin{itemize}[leftmargin=18pt, itemsep=2pt, topsep=2pt, label=\textbullet]
\item 직유법 — '쟁반같이 둥근 달'(달=원관념, 쟁반=보조 관념, 연결어 '같이')
\item 은유법 — '인생은 마라톤이다'(원관념=인생, 보조 관념=마라톤)
\item 의인법 — '바람이 속삭인다'(바람을 사람처럼 표현)
\end{itemize}

\kfSecHeader{kfAreaConcept}{문제}{}

\begin{kfQBlock}{1}{}
\kfQStem{다음 중 '은유법'이 사용되지 \uline{않은} 문장은 무엇인가요?}
\begin{kfChoices}
\kfChoice 구름은 솜사탕처럼 하얗다.
\kfChoice 책은 마음의 양식이다.
\kfChoice 파도가 바위에 부딪혀 소리친다.
\kfChoice 내 목소리는 꾀꼬리 같다.
\end{kfChoices}
\end{kfQBlock}
\begin{kfQBlock}{2}{}
\kfQStem{밑줄 친 부분에 사용된 비유법은 무엇인가요?}
\begin{kfEvidence}“가로수들이 \uline{나란히 서서 손을 흔들며 반겨준다}.”\end{kfEvidence}
\begin{kfChoices}
\kfChoice 직유법
\kfChoice 은유법
\kfChoice 의인법
\kfChoice 반복법
\end{kfChoices}
\end{kfQBlock}
\begin{kfQBlock}{3}{}
\kfQStem{다음 문장에 사용된 비유법과 원관념이 바르게 짝지어지지 \uline{않은} 것은 무엇인가요?}
\begin{kfEvidence}“아버지의 손은 거친 나무껍질 같았다.”\end{kfEvidence}
\begin{kfChoices}
\kfChoice 비유법: 직유법 / 원관념: 아버지의 손
\kfChoice 비유법: 은유법 / 원관념: 아버지의 손
\kfChoice 비유법: 직유법 / 원관념: 나무껍질
\kfChoice 비유법: 은유법 / 원관념: 나무껍질
\end{kfChoices}
\end{kfQBlock}
\begin{kfQBlock}{1}{}
\kfQStem{원관념: 어머니의 사랑 (직유법 사용)}
\kfAnswerNote{답안}{4}
\end{kfQBlock}
\begin{kfQBlock}{2}{}
\kfQStem{원관념: 도서관 (은유법 사용)}
\kfAnswerNote{답안}{4}
\end{kfQBlock}
\begin{kfQBlock}{3}{}
\kfQStem{원관념: 참새 (의인법 사용)}
\kfAnswerNote{답안}{4}
\end{kfQBlock}

\kfSecHeader{kfAreaConcept}{활동}{비유하기}
\noindent\textit{\small [1-3] 다음 비유적 표현에서 '원관념'과 '보조 관념'을 찾아보세요.}\par\medskip
\begin{kfBlankList}
\kfBlankItem “시간은 금이다.”
\kfBlankItem “쟁반같이 둥근 달”
\kfBlankItem “내 동생은 청개구리야.”
\end{kfBlankList}

\kfSecHeader{kfAreaConcept}{활동}{비유하기}
\noindent\textit{\small 다음 진술이 맞으면 ○, 틀리면 × 표시하시오.}\par\medskip

\kfSecHeader{kfAreaConcept}{문제}{}



\kfStartLit{문학 영역 — 배경 지식 돋보기}


\kfSecHeader{kfAreaLiterature}{지문}{배경 지식 돋보기}
\kfPassageNote{배경 지식 돋보기}{%
시를 더 깊이 이해하기 위해, 시가 쓰인 시대 상황을 읽어 봅시다.

“풀은 그냥 풀이 아니에요” 이 시를 쓴 김수영 시인이 살았던 1960년대는 국민들이 자신의 생각을 자유롭게 말하기 \uline{어려운} 시절이었어요. 힘을 가진 권력자들이 국민들을 억누르고 괴롭히기도 했지요. 시인은 무서운 힘 앞에서 기죽지 않고 끝까지 버티는 국민들의 모습이 마치 '풀'과 닮았다고 생각했어요. 풀은 바람이 불면 금방 눕고 울지만, 바람이 지나가면 누구보다 먼저 다시 일어나는 강한 생명력을 가졌으니까요. 그래서 이 시에 나오는 '비바람'은 국민들을 괴롭히는 독재 권력이나 시련을, '풀'은 끈질기게 저항하고 이겨내는 평범한 시민들을 상징한답니다.
}


\kfSecHeader{kfAreaLiterature}{지문}{작품}
\kfPassageNote{작품}{%
풀이 눕는다. \\[4pt]
비를 몰고 오는 비바람에 흔들려 \\[4pt]
풀은 눕고 \\[4pt]
드디어 울었다. \\[4pt]
날이 흐려서 더 울다가 \\[4pt]
다시 울었다. \\[14pt]
풀이 눕는다. \\[4pt]
바람보다도 더 빨리 눕는다. \\[4pt]
바람보다도 더 빨리 울고 \\[4pt]
바람보다 먼저 일어난다. \\[14pt]
날이 흐리고 풀이 눕는다. \\[4pt]
발목까지 \\[4pt]
발밑까지 눕는다. \\[4pt]
바람보다 늦게 누워도 \\[4pt]
바람보다 먼저 일어나고 \\[4pt]
바람보다 늦게 울어도 \\[4pt]
바람보다 먼저 웃는다. \\[4pt]
날이 흐리고 풀뿌리가 눕는다.
}


\kfSecHeader{kfAreaLiterature}{활동}{배경 지식 돋보기}
\noindent\textit{\small [1-2] 배경 지식을 바탕으로 볼 때, 다음 시어들이 실제로 의미하는 것은 무엇일까요? 선으로 연결해 보세요.}\par\medskip
\begin{kfConnect}
\kfPair{풀}{㉠ 국민들을 억누르는 독재 권력이나 힘든 시련}
\kfPair{비바람}{㉡ 억압에도 굴하지 않고 다시 일어나는 강한 시민들}
\end{kfConnect}

\kfSecHeader{kfAreaLiterature}{활동}{차근차근 내용 확인}
\noindent\textit{\small [1-5] 글의 내용과 일치하면 O, 일치하지 않으면 X를 쓰세요.}\par\medskip
\begin{kfOXList}
\kfOXItem 이 시에서 풀은 비바람에 흔들려 결국 울고 만다. \kfOX
\kfOXItem 풀은 바람보다 늦게 눕고, 바람보다 늦게 일어난다. \kfOX
\kfOXItem 시의 1연에서는 풀뿌리가 눕는 모습이 구체적으로 제시된다. \kfOX
\kfOXItem 풀은 바람보다 늦게 울고, 바람보다 먼저 웃는다. \kfOX
\kfOXItem 이 시에서 날이 흐린 것은 풀이 눕고 우는 상황과 관련이 있다. \kfOX
\end{kfOXList}

\kfSecHeader{kfAreaLiterature}{활동}{나만의 생각 펼치기}
\noindent\textit{\small 다음 조건에 맞게 글을 써 보세요.}\par\medskip
\kfWritingPrompt{여러분도 '풀'처럼 힘들고 슬픈 일(비바람)이 있었지만, 포기하지 않고 다시 일어난 경험이 있나요? - 나에게 불어온 비바람(힘든 일):  - 내가 다시 일어난 방법(극복):}
\kfWriting{답안 작성}{8}

\kfSecHeader{kfAreaLiterature}{문제}{}

\begin{kfQBlock}{1}{}
\kfQStem{시인이 국민들을 '풀'에 빗대어 표현한 이유로 가장 알맞은 것은 무엇인가요?}
\begin{kfChoices}
\kfChoice 풀은 사람들에게 밟히면 시들어 흔적도 없이 사라지기 때문에
\kfChoice 풀은 바람이 불어오면 이리저리 흔들리며 춤을 추기 때문에
\kfChoice 풀은 눕고 울더라도, 다시 일어나는 강한 생명력이 있기 때문에
\kfChoice 풀은 햇볕을 받고 자라서 나중에는 커다란 나무가 되기 때문에
\end{kfChoices}
\end{kfQBlock}
\begin{kfQBlock}{2}{}
\kfQStem{풀이 “바람보다 먼저 일어난다”라는 것은 시민들의 어떤 태도를 말하는 걸까요?}
\begin{kfChoices}
\kfChoice 시련을 이겨내려는 강한 의지
\kfChoice 괴로워서 도망치고 싶은 마음
\kfChoice 힘센 권력자를 무조건 따르는 태도
\kfChoice 아무 생각 없이 시간에 맡기는 태도
\end{kfChoices}
\end{kfQBlock}
\begin{kfQBlock}{3}{}
\kfQStem{마지막 연에서 풀은 울다가 결국 '웃는다'고 했습니다. 이때 시민들의 마음 상태로 가장 적절한 것은 무엇인가요?}
\begin{kfChoices}
\kfChoice 너무 슬퍼서 헛웃음만 나왔을 것이다.
\kfChoice 아무 생각 없이 그냥 웃었을 것이다.
\kfChoice 이겨냈다는 기쁨과 자신감이 생겼을 것이다.
\kfChoice 바람이 다시 불어올까 봐 매우 걱정했을 것이다.
\end{kfChoices}
\end{kfQBlock}
\begin{kfQBlock}{1}{}
\kfQStem{이 시에 사용된 비유적 표현 중, 원관념과 보조 관념의 연결이 바르게 짝지어진 것은 무엇인가요?}
\begin{kfChoices}
\kfChoice 원관념: 시련, 억압 / 보조 관념: 풀뿌리
\kfChoice 원관념: 독재 권력 / 보조 관념: 비바람
\kfChoice 원관념: 억압받는 민중 / 보조 관념: 울었다
\kfChoice 원관념: 풀의 강한 생명력 / 보조 관념: 눕는다
\end{kfChoices}
\end{kfQBlock}
\begin{kfQBlock}{1}{}
\kfQStem{이 시에서 '풀'을 괴롭히며 눕게 만드는 존재는 무엇인가요?}
\begin{kfChoices}
\kfChoice 햇살
\kfChoice 친구
\kfChoice 비바람
\kfChoice 돌멩이
\end{kfChoices}
\end{kfQBlock}
\begin{kfQBlock}{2}{}
\kfQStem{시의 내용에 따르면 '풀'의 행동으로 알맞지 \uline{않은} 것은 무엇인가요?}
\begin{kfChoices}
\kfChoice 바람보다 먼저 웃는다.
\kfChoice 바람이 무서워 도망친다.
\kfChoice 바람보다 먼저 일어난다.
\kfChoice 바람보다 더 빨리 눕는다.
\end{kfChoices}
\end{kfQBlock}
\begin{kfQBlock}{3}{}
\kfQStem{3연의 '바람보다 늦게 누워도 / 바람보다 먼저 일어나고 / 바람보다 늦게 울어도 / 바람보다 먼저 웃는다.'는 풀의 어떤 특성을 가장 잘 드러내고 있나요?}
\begin{kfChoices}
\kfChoice 풀은 바람이 불어오면 두려움을 느끼고 피하려고만 한다.
\kfChoice 풀은 바람의 힘에 저항하지 못하고 순순히 복종하기만 한다.
\kfChoice 풀은 쓰러져도 다시 일어나는 강한 생명력을 가지고 있다.
\kfChoice 풀은 주변 환경이 어떻게 변하더라도 전혀 영향을 받지 않는다.
\end{kfChoices}
\end{kfQBlock}
\begin{kfQBlock}{1}{}
\kfQStem{시상 전개 구조: 빈칸을 채워 각 연의 상황과 풀의 태도를 정리해 봅시다. 1연: 비바람에 풀이 흔들려 ( ① 　　　 ) 2연: 바람보다 빨리 눕지만, 바람보다 ( ② 　　　 ) 3연: 발밑까지 눕지만, 바람보다 늦게 울고 ( ③ 　　　 )}
\kfAnswerLines{1}
\end{kfQBlock}
\begin{kfQBlock}{2}{}
\kfQStem{시의 상징적 의미: '풀'과 '바람'의 대조 구조를 완성해 봅시다. 약자(풀): ( ④ 　　　 ) 강자(바람): ( ⑤ 　　　 )}
\kfAnswerLines{1}
\end{kfQBlock}


\kfStartNonlit{비문학 영역 — 국민이 주인인 나라, 민주주의와 깨어 있는 시민}


\kfSecHeader{kfAreaNonlit}{지문}{국민이 주인인 나라, 민주주의와 깨어 있는 시민}
\kfPassageNote{국민이 주인인 나라, 민주주의와 깨어 있는 시민}{%
　대한민국은 '민주 공화국'입니다. \uline{㉠이는} 국가의 주인이 국민이며, 국민을 위한 정치가 이루어지는 나라라는 뜻입니다. 하지만 우리 역사를 살펴보면 민주주의가 처음부터 자리 잡은 것은 아니었습니다. 수많은 시민들의 희생과 참여가 있었기에 오늘날의 민주주의가 가능했습니다.

　1960년 4·19 혁명은 부정 선거에 맞서 일어난 최초의 민주주의 혁명입니다. 학생과 시민들은 “국민의 뜻에 따르라”고 외치며 독재 정부를 무너뜨렸습니다. 1980년 5·18 민주화 운동에서는 광주 시민들이 군사 정권에 맞서 민주주의를 지키기 위해 싸웠습니다. 또한, 1987년 6월 민주 항쟁은 대통령을 국민의 손으로 직접 뽑는 '대통령 직선제'를 만들어 냈습니다. \uline{㉡이처럼} 대한민국의 민주주의는 시민들의 피와 땀으로 발전해 왔습니다.

　오늘날의 민주주의는 선거를 통한 참여뿐만 아니라 다양한 방식으로 이루어집니다. 시민들은 신문과 방송, 시민 단체 활동을 통해 자신의 의견을 표현하고, 토론회나 주민 회의에 참석하여 지역 사회의 문제를 해결합니다. 최근에는 인터넷과 SNS를 통해 정치적 생각을 밝히고 정책을 제안하는 일도 활발해졌습니다.

　민주주의 사회에서 시민의 역할은 매우 중요합니다. 시민이 정치에 무관심하면 권력은 부패하기 쉽고, 국민의 권리는 빼앗길 수 있습니다. 따라서 우리는 사회 문제에 관심을 가지고, 비판적인 시각으로 뉴스를 보며, 선거에 빠짐없이 참여해야 합니다. 깨어 있는 시민들의 참여가 건강한 민주주의를 만듭니다.
}


\kfSecHeader{kfAreaNonlit}{활동}{어휘력 쑥쑥! 단어 익히기}
\noindent\textit{\small [1-3] 다음 단어와 그 뜻이 올바르게 연결되도록 선을 그으세요.}\par\medskip
\begin{kfConnect}
\kfPair{직선제}{㉠ 정치나 사회 제도의 잘못된 점을 따져 보고 평가하는 태도}
\kfPair{비판적}{㉡ 국민이 직접 투표하여 대통령 등의 대표를 뽑는 제도}
\kfPair{부패}{㉢ 정치나 사회가 타락하여 올바르지 못한 상태로 변하는 것}
\end{kfConnect}

\kfSecHeader{kfAreaNonlit}{활동}{차근차근 내용 확인}
\noindent\textit{\small [1-5] 글의 내용과 일치하면 O, 일치하지 않으면 X를 쓰세요.}\par\medskip
\begin{kfOXList}
\kfOXItem 대한민국이 '민주 공화국'인 것은 수많은 시민들의 희생과 참여가 없었다면 불가능했을 것이다. \kfOX
\kfOXItem 6월 민주 항쟁을 통해 독재 정부가 무너지고, 국민이 직접 대통령을 뽑게 되었다. \kfOX
\kfOXItem 현대 사회에서 시민들이 정치에 참여하는 방법은 신문, 방송 등 전통적인 매체에만 한정된다. \kfOX
\kfOXItem 시민들이 정치에 무관심할 경우 권력이 부패할 위험성이 커진다고 글쓴이는 지적한다. \kfOX
\kfOXItem 깨어 있는 시민들의 참여는 글의 마지막 문단에서 언급된 '건강한 민주주의'의 필수 조건이다. \kfOX
\end{kfOXList}

\kfSecHeader{kfAreaNonlit}{활동}{나만의 생각 펼치기}
\noindent\textit{\small 다음 조건에 맞게 글을 써 보세요.}\par\medskip
\kfWritingPrompt{내가 우리 학교나 마을을 위해 실천할 수 있는 '참여' 방법 한 가지를 구체적으로 써 보세요.}
\kfWriting{답안 작성}{8}
\kfWritingPrompt{이 글의 내용을 바탕으로 '건강한 민주주의'를 만들기 위해 가장 필요한 것은 무엇인지 자신의 생각을 2\textasciitilde{}3문장으로 정리해 보세요.}
\kfWriting{답안 작성}{8}

\kfSecHeader{kfAreaNonlit}{문제}{}

\begin{kfQBlock}{1}{}
\kfQStem{밑줄 친 단어와 바꿔 쓰기에 가장 적절한 말은 무엇인가요?}
\begin{kfEvidence}“권력은 \uline{부패}하기 쉽고, 국민의 권리는 빼앗길 수 있습니다.”\end{kfEvidence}
\begin{kfChoices}
\kfChoice 발전
\kfChoice 썩기
\kfChoice 성장
\kfChoice 성공
\end{kfChoices}
\end{kfQBlock}
\begin{kfQBlock}{1}{}
\kfQStem{㉠ '이는'이 가리키는 내용으로 가장 적절한 것은 무엇인가요?}
\begin{kfChoices}
\kfChoice 우리나라의 지난 역사를 차근차근 살펴보는 것
\kfChoice 대한민국이 '민주 공화국'이라는 사실
\kfChoice 민주주의가 우리나라에 처음부터 자리 잡은 것
\kfChoice 시민들의 희생과 참여가 있었던 것
\end{kfChoices}
\end{kfQBlock}
\begin{kfQBlock}{2}{}
\kfQStem{㉡ '이처럼'이 가리키는 내용을 지문에서 찾아 정리할 때, 포함되지 않는 사건은 무엇인가요?}
\begin{kfChoices}
\kfChoice 4·19 혁명
\kfChoice 5·18 민주화 운동
\kfChoice 6월 민주 항쟁
\kfChoice 주민 회의 참석
\end{kfChoices}
\end{kfQBlock}
\begin{kfQBlock}{3}{}
\kfQStem{다음 문장을 읽고 '원인(조건)'과 그에 따른 '결과'를 바르게 연결한 것은 무엇인가요? “시민이 정치에 무관심하면 권력은 부패하기 쉽고, 국민의 권리는 빼앗길 수 있습니다.”}
\begin{kfChoices}
\kfChoice 원인: 국민의 권리가 빼앗김 / 결과: 시민이 정치에 무관심해짐
\kfChoice 원인: 권력이 부패함 / 결과: 시민이 정치에 무관심해짐
\kfChoice 원인: 시민이 정치에 무관심함 / 결과: 권력이 부패하고 권리를 빼앗김
\kfChoice 원인: 시민이 정치에 무관심함 / 결과: 건강한 민주주의가 만들어짐
\end{kfChoices}
\end{kfQBlock}
\begin{kfQBlock}{1}{}
\kfQStem{이 글을 읽고 알 수 있는 내용으로 적절하지 \uline{않은} 것은 무엇인가요?}
\begin{kfChoices}
\kfChoice 대한민국은 국가의 주인이 국민인 민주 공화국이다.
\kfChoice 대한민국의 민주주의는 시민들의 희생과 참여로 발전했다.
\kfChoice 오늘날에는 선거 외에는 정치에 참여할 방법이 없다.
\kfChoice 시민이 정치에 관심을 가져야 국민의 권리를 지킬 수 있다.
\end{kfChoices}
\end{kfQBlock}
\begin{kfQBlock}{2}{}
\kfQStem{이 글의 중심 내용으로 알맞지 \uline{않은} 것은 무엇인가요?}
\begin{kfChoices}
\kfChoice 대한민국 헌법의 역사와 종류
\kfChoice 독재 정권이 생겨나는 원인과 과정
\kfChoice 인터넷과 SNS를 활용한 선거 운동 방법
\kfChoice 민주주의의 발전 과정과 시민 참여의 중요성
\end{kfChoices}
\end{kfQBlock}
\begin{kfQBlock}{3}{}
\kfQStem{글쓴이의 주장을 바탕으로 추론할 때, 민주주의 사회에서 국민의 권리가 빼앗기는 근본적인 이유로 가장 적절한 것은 무엇인가요?}
\begin{kfChoices}
\kfChoice 시민들이 인터넷과 SNS를 너무 많이 사용하기 때문에
\kfChoice 권력자들이 너무 많은 활동을 하기 때문에
\kfChoice 시민이 사회 문제에 관심을 두지 않고 정치에 무관심하기 때문에
\kfChoice 선거를 통한 참여 방법이 너무 복잡하고 어렵기 때문에
\end{kfChoices}
\end{kfQBlock}
\begin{kfQBlock}{4}{}
\kfQStem{<보기>의 '지수'의 행동이 글에서 제시된 '깨어 있는 시민'의 역할과 일치하는 부분은 무엇인가요?}
\begin{kfEvidence}<보기>   5학년 지수는 학교 앞 횡단보도가 위험하다고 생각했습니다. 그래서 시청 홈페이지에 “어린이 보호 구역 안전장치를 설치해 주세요”라는 글을 올렸습니다.\end{kfEvidence}
\begin{kfChoices}
\kfChoice 신문을 통해 자신의 의견을 표현했다.
\kfChoice 군사 정권에 맞서 독재에 저항했다.
\kfChoice 인터넷을 통해 지역 사회 문제 해결에 참여했다.
\kfChoice 선거에 출마하여 대표가 되려고 노력하고 있다.
\end{kfChoices}
\end{kfQBlock}
\begin{kfQBlock}{5}{}
\kfQStem{다음 중 '깨어 있는 시민'의 태도로 보기 \uline{어려운} 것은 무엇인가요?}
\begin{kfChoices}
\kfChoice 사회 문제에 대해 “나랑 상관없는 일이야”라며 신경 쓰지 않는다.
\kfChoice 뉴스를 볼 때 사실인지 아닌지 비판적인 시각으로 바라본다.
\kfChoice 선거 날 귀찮더라도 투표소에 가서 소중한 한 표를 행사한다.
\kfChoice 우리 동네의 불편한 점을 해결하기 위해 주민 회의에 참석한다.
\end{kfChoices}
\end{kfQBlock}
\begin{kfQBlock}{1}{}
\kfQStem{다음은 대한민국의 민주화 과정을 정리한 표입니다. (가)와 (나)에 들어갈 알맞은 말을 쓰세요. 1960년: 4·19 혁명 (부정 선거에 맞서 독재 정부를 무너뜨림) 1980년: ( 가 ) (광주 시민들이 군사 정권에 맞서 싸움) 1987년: 6월 민주 항쟁 ( ( 나 )를 이끌어 냄 )}
\kfAnswerLines{1}
\end{kfQBlock}
\begin{kfQBlock}{1}{}
\kfQStem{다음 문장에서 '누가(주어)'와 '무엇을 했는지(서술어)'를 찾아 쓰세요. “1987년 6월 민주 항쟁은 대통령을 국민의 손으로 직접 뽑는 대통령 직선제를 만들어 냈습니다.” 누가(무엇이): \_\_\_\_\_\_\_\_\_\_\_\_\_\_\_\_\_\_\_\_\_\_\_\_\_\_\_\_\_\_\_\_\_\_\_\_\_\_\_\_ 무엇을 했나: \_\_\_\_\_\_\_\_\_\_\_\_\_\_\_\_\_\_\_\_\_\_\_\_\_\_\_\_\_\_\_\_\_\_\_\_\_\_\_\_}
\kfAnswerNote{답안}{4}
\end{kfQBlock}


\kfStartWeekly{실력확인 영역 — 주간 실력 확인}


\noindent\textit{\small (제한시간: 30분 / 총 20문항)}\par\medskip
\par\noindent\textbf{[6-10] 다음 시를 읽고, 물음에 답하세요.}\par\medskip
\kfPassageNote{지문 6}{%
풀이 눕는다. \\[4pt]
비를 몰고 오는 \uline{㉠비바람}에 흔들려 \\[4pt]
풀은 눕고 \\[4pt]
드디어 울었다. \\[4pt]
날이 흐려서 더 울다가 \\[4pt]
다시 울었다. \\[14pt]
풀이 눕는다. \\[4pt]
바람보다도 더 빨리 눕는다. \\[4pt]
\uline{㉡바람보다도 더 빨리 울고} \\[4pt]
\uline{㉢바람보다 먼저 일어난다.} \\[14pt]
날이 흐리고 풀이 눕는다. \\[4pt]
발목까지 \\[4pt]
발밑까지 눕는다. \\[4pt]
바람보다 늦게 누워도 \\[4pt]
바람보다 먼저 일어나고 \\[4pt]
바람보다 늦게 울어도 \\[4pt]
\uline{㉣바람보다 먼저 웃는다.} \\[4pt]
날이 흐리고 풀뿌리가 눕는다.
}
\par\noindent\textbf{[11-15] 다음 글을 읽고, 물음에 답하세요.}\par\medskip
\kfPassageNote{지문 11}{%
대한민국은 '민주 공화국'입니다. 이는 국가의 주인이 국민이며, 국민을 위한 정치가 이루어지는 나라라는 뜻입니다. 하지만 우리 역사를 살펴보면 민주주의가 처음부터 자리 잡은 것은 아니었습니다. 수많은 시민들의 희생과 참여가 있었기에 오늘날의 민주주의가 가능했습니다. 1960년 4·19 혁명은 부정 선거에 맞서 일어난 최초의 민주주의 혁명입니다. 학생과 시민들은 “국민의 뜻에 따르라”고 외치며 독재 정부를 무너뜨렸습니다. 1980년 5·18 민주화 운동에서는 광주 시민들이 군사 정권에 맞서 민주주의를 지키기 위해 싸웠습니다. 또한, 1987년 6월 민주 항쟁은 대통령을 국민의 손으로 직접 뽑는 '대통령 직선제'를 만들어 냈습니다. 이처럼 대한민국의 민주주의는 시민들의 피와 땀으로 발전해 왔습니다. 오늘날의 민주주의는 선거를 통한 참여뿐만 아니라 다양한 방식으로 이루어집니다. 시민들은 신문과 방송, 시민 단체 활동을 통해 자신의 의견을 표현하고, 토론회나 주민 회의에 참석하여 지역 사회의 문제를 해결합니다. 최근에는 인터넷과 SNS를 통해 정치적 생각을 밝히고 정책을 제안하는 일도 활발해졌습니다. 민주주의 사회에서 시민의 역할은 매우 중요합니다. 시민이 정치에 무관심하면 권력은 부패하기 쉽고, 국민의 권리는 빼앗길 수 있습니다. 따라서 우리는 사회 문제에 관심을 가지고, 비판적인 시각으로 뉴스를 보며, 선거에 빠짐없이 참여해야 합니다. 깨어 있는 시민들의 참여가 건강한 민주주의를 만듭니다.
}
\begin{kfQBlock}{1}{}
\kfQStem{다음 단어와 그 뜻풀이가 바르게 연결되지 \uline{않은} 것은 무엇인가요?}
\begin{kfChoices}
\kfChoice 공공 - 개인적인 이익만을 위하는 것
\kfChoice 정권 - 정치를 맡아 나라를 다스리는 힘
\kfChoice 권력 - 남을 따르게 하거나 지배할 수 있는 힘
\kfChoice 독재 - 특정한 개인이나 집단이 권력을 독차지하여 마음대로 정치를 하는 것
\end{kfChoices}
\end{kfQBlock}

\begin{kfQBlock}{2}{}
\kfQStem{다음 중 '혁명'이라는 단어를 사용하기에 가장 적절한 상황은 무엇인가요?}
\begin{kfChoices}
\kfChoice 친구와 약속을 정할 때
\kfChoice 잃어버린 물건을 찾았을 때
\kfChoice 반장 선거에 후보로 나갈 때
\kfChoice 사회 제도를 완전히 바꿀 때
\end{kfChoices}
\end{kfQBlock}

\begin{kfQBlock}{3}{}
\kfQStem{다음 문장의 빈칸에 들어갈 알맞은 말은 무엇인가요?}
\begin{kfEvidence}“민주주의 국가에서 모든 (　 　 　)은 국민으로부터 나온다.”\end{kfEvidence}
\begin{kfChoices}
\kfChoice 왕
\kfChoice 독재
\kfChoice 권력
\kfChoice 부패
\end{kfChoices}
\end{kfQBlock}

\begin{kfQBlock}{4}{}
\kfQStem{밑줄 친 단어와 바꿔 쓰기에 가장 적절한 말은 무엇인가요?}
\begin{kfEvidence}“우리는 이번 주말에 열리는 마을 청소 봉사에 \uline{참여}하기로 했다.”\end{kfEvidence}
\begin{kfChoices}
\kfChoice 구경
\kfChoice 참가
\kfChoice 반대
\kfChoice 후회
\end{kfChoices}
\end{kfQBlock}

\begin{kfQBlock}{5}{}
\kfQStem{다음 단어들의 관계가 나머지와 \uline{다른} 것은 무엇인가요?}
\begin{kfChoices}
\kfChoice 권력 - 힘
\kfChoice 공공 - 공동
\kfChoice 정권 - 정부
\kfChoice 독재 - 민주주의
\end{kfChoices}
\end{kfQBlock}

\begin{kfQBlock}{6}{}
\kfQStem{이 시에서 '비바람'과 바꾸어 쓸 수 있는 시어는 무엇인가요?}
\begin{kfChoices}
\kfChoice 평화
\kfChoice 행복
\kfChoice 억압
\kfChoice 자유
\end{kfChoices}
\end{kfQBlock}

\begin{kfQBlock}{7}{}
\kfQStem{시의 내용에 대한 설명으로 알맞지 \uline{않은} 것은 무엇인가요?}
\begin{kfChoices}
\kfChoice 풀은 비바람 때문에 눕고 운다.
\kfChoice 풀은 바람보다 더 빨리 눕는다.
\kfChoice 풀은 바람보다 먼저 일어난다.
\kfChoice 풀은 바람이 무서워서 도망친다.
\end{kfChoices}
\end{kfQBlock}

\begin{kfQBlock}{8}{}
\kfQStem{“바람보다 더 빨리 눕는다”와 “바람보다 먼저 일어난다”는 표현에 담긴 '풀'의 행동 특징을 가장 잘 설명한 것은 무엇인가요?}
\begin{kfChoices}
\kfChoice 바람이 너무 무서워서 정신없이 움직이고 있다.
\kfChoice 바람이 불기도 전에 미리 알아채고 피하고 있다.
\kfChoice 아무런 생각 없이 바람이 부는 대로 흔들리고 있다.
\kfChoice 시련에 민감하게 반응하며 능동적으로 이겨내고 있다.
\end{kfChoices}
\end{kfQBlock}

\begin{kfQBlock}{9}{}
\kfQStem{이 시를 읽고 난 후의 감상으로 가장 적절한 것은 무엇인가요?}
\begin{kfChoices}
\kfChoice 비가 와서 풀이 자라는 모습이 신기해.
\kfChoice 바람이 불면 풀이 눕는 건 당연한 일이야.
\kfChoice 약해 보여도 포기하지 않는 마음이 중요해.
\kfChoice 농사를 지을 때는 날씨가 정말 중요한 것 같아.
\end{kfChoices}
\end{kfQBlock}

\begin{kfQBlock}{10}{}
\kfQStem{이 시의 주제로 가장 적절한 것은 무엇인가요?}
\begin{kfChoices}
\kfChoice 비바람을 막아 주는 자연의 다정한 보호
\kfChoice 풀이 바람보다 먼저 눕고 늦게 일어나는 차분한 순응의 자세
\kfChoice 비바람 속에서 풀과 사람이 함께 외로움을 견디는 슬픈 정서
\kfChoice 비바람(시련)에 짓눌려도 다시 일어나 웃는 민중의 끈질긴 생명력
\end{kfChoices}
\end{kfQBlock}

\begin{kfQBlock}{11}{}
\kfQStem{4·19 혁명이나 6월 민주 항쟁 같은 사건이 일어날 수밖에 없었던 이유는 무엇인가요?}
\begin{kfChoices}
\kfChoice 나라의 경제를 더 빠르게 발전시키기 위해서
\kfChoice 다른 나라보다 군사력이 강함을 보여 주기 위해서
\kfChoice 대통령이 왕처럼 나라를 마음대로 다스리기 위해서
\kfChoice 국민의 권리가 제대로 지켜지지 않았기 때문에
\end{kfChoices}
\end{kfQBlock}

\begin{kfQBlock}{12}{}
\kfQStem{1987년 6월 민주 항쟁의 결과로 얻어낸 제도는 무엇인가요?}
\begin{kfChoices}
\kfChoice 대통령 직선제
\kfChoice 국회의원 선거
\kfChoice 반장 선거
\kfChoice 주민 투표제
\end{kfChoices}
\end{kfQBlock}

\begin{kfQBlock}{13}{}
\kfQStem{오늘날 시민들이 정치에 참여하는 방법으로 알맞지 \uline{않은} 것은 무엇인가요?}
\begin{kfChoices}
\kfChoice 선거에 참여하여 투표권을 행사한다.
\kfChoice 인터넷이나 SNS를 통해 의견을 표현한다.
\kfChoice 주민 회의에 참석하여 문제를 해결한다.
\kfChoice 자신의 이익을 위해 법을 어긴다.
\end{kfChoices}
\end{kfQBlock}

\begin{kfQBlock}{14}{}
\kfQStem{다음 <보기>의 상황에 대해 글쓴이가 할 수 있는 조언으로 가장 적절한 것은 무엇인가요?}
\begin{kfEvidence}<보기> 　A나라의 국민들은 정치에 관심이 없습니다. 선거일에는 투표하러 가는 대신 놀러 가고, 뉴스는 재미없다며 보지 않습니다.\end{kfEvidence}
\begin{kfChoices}
\kfChoice 정치에 관심이 없으니 마음 편하게 살 수 있겠군요.
\kfChoice 권력이 부패하여 국민의 권리를 빼앗길 수 있으니 주의하세요.
\kfChoice 정치인들이 알아서 잘할 테니 국민들은 생업에만 집중하세요.
\kfChoice 강력한 독재자가 나타나서 나라를 다스리는 것이 더 효율적입니다.
\end{kfChoices}
\end{kfQBlock}

\begin{kfQBlock}{15}{}
\kfQStem{위 글의 내용을 바탕으로 추론한 내용 중 적절하지 \uline{않은} 것은 무엇인가요?}
\begin{kfChoices}
\kfChoice 시민의 참여가 늘어날수록 권력은 더 부패하기 쉬워진다.
\kfChoice 과거에는 대통령을 국민이 직접 뽑지 못했던 시절도 있었다.
\kfChoice 오늘날에는 직접 시위에 나가지 않아도 정치에 참여할 방법이 있다.
\kfChoice 민주주의는 저절로 얻어진 것이 아니라 시민들이 싸워서 얻어낸 것이다.
\end{kfChoices}
\end{kfQBlock}

\begin{kfQBlock}{16}{}
\kfQStem{다음 문장에서 '원관념'과 '보조 관념'을 바르게 짝지은 것은 무엇인가요?}
\begin{kfEvidence}“운동장은 아이들이 뛰노는 거대한 개미집이다.”\end{kfEvidence}
\begin{kfChoices}
\kfChoice 원관념: 개미집 / 보조 관념: 운동장
\kfChoice 원관념: 운동장 / 보조 관념: 개미집
\kfChoice 원관념: 아이들 / 보조 관념: 개미
\kfChoice 원관념: 운동장 / 보조 관념: 아이들
\end{kfChoices}
\end{kfQBlock}

\begin{kfQBlock}{17}{}
\kfQStem{다음 중 <보기>의 설명에 해당하는 문장은 무엇인가요?}
\begin{kfEvidence}<보기> 　이 비유법은 원관념과 보조 관념을 'A는 B다'와 같은 형식으로 은근하게 연결합니다.\end{kfEvidence}
\begin{kfChoices}
\kfChoice 시간은 쏜살같이 흘러간다.
\kfChoice 빗방울이 창문을 두드린다.
\kfChoice 선생님은 우리의 등대이시다.
\kfChoice 아기 피부가 비단처럼 곱다.
\end{kfChoices}
\end{kfQBlock}

\begin{kfQBlock}{18}{}
\kfQStem{다음 중 '직유법'이 사용된 문장은 무엇인가요?}
\begin{kfChoices}
\kfChoice 내 마음은 호수요.
\kfChoice 쟁반같이 둥근 달.
\kfChoice 학교는 우리들의 집이다.
\kfChoice 파도가 바위에 입 맞춘다.
\end{kfChoices}
\end{kfQBlock}

\begin{kfQBlock}{19}{}
\kfQStem{다음 중 '의인법'이 사용된 문장은 무엇인가요?}
\begin{kfChoices}
\kfChoice 사과 같은 내 얼굴
\kfChoice 호수같이 맑은 눈동자
\kfChoice 내 동생은 호랑이 선생님이다
\kfChoice 비가 오니 우산이 활짝 웃는다
\end{kfChoices}
\end{kfQBlock}

\begin{kfQBlock}{20}{}
\kfQStem{다음 중 비유의 종류와 그 예시가 알맞게 짝지어진 것은 무엇인가요?}
\begin{kfChoices}
\kfChoice 의인법 - “내 동생은 청개구리다.”
\kfChoice 은유법 - “책은 마음의 양식이다.”
\kfChoice 직유법 - “꽃이 나를 보고 웃는다.”
\kfChoice 은유법 - “아버지 손은 거친 나무껍질 같다.”
\end{kfChoices}
\end{kfQBlock}




\clearpage

\kfChapterCover{3}{프레게2 (챕터3)}{Chapter 3}{이번 챕터의 학습 내용을 시작합니다.}

\kfStartVocab{어휘 영역 — 건강과 과학 관련 어휘}


\kfSecHeader{kfAreaVocab}{문제}{}

\begin{kfQBlock}{1}{}
\kfQStem{면역}
\kfAnswerNote{답안}{4}
\end{kfQBlock}
\begin{kfQBlock}{2}{}
\kfQStem{백신}
\kfAnswerNote{답안}{4}
\end{kfQBlock}

\kfSecHeader{kfAreaVocab}{활동}{건강과 과학 관련 어휘}
\noindent\textit{\small [1-5] <보기>에서 알맞은 단어를 골라 문장을 완성하세요.}\par\medskip
\begin{kfBlankList}
\kfBlankItem 독감 예방을 위해 병원에 가서 [　　　　　]을/를 맞았다.
\begin{kfEvidence}보\end{kfEvidence}
\kfBlankItem [　　　　　]에 감염되어 아플 때는 의사 선생님의 처방에 따라 항생제를 복용해야 한다.
\begin{kfEvidence}기\end{kfEvidence}
\kfBlankItem 평소 건강 관리를 잘하여 [　　　　　]력을 높이는 것이 중요하다.
\begin{kfEvidence}>\end{kfEvidence}
\kfBlankItem 감기는 주로 [　　　　　] 때문에 생기며, 특별한 치료 없이도 나을 수 있다.
\begin{kfEvidence} \end{kfEvidence}
\kfBlankItem 예방 접종을 하면 몸 안에 병균과 싸울 수 있는 [　　　　　]가/이 생긴다.
\begin{kfEvidence}　\end{kfEvidence}
\end{kfBlankList}

\kfSecHeader{kfAreaVocab}{활동}{건강과 과학 관련 어휘}
\noindent\textit{\small [1-5] 단어의 정확한 의미를 찾아 선으로 연결해 보세요.}\par\medskip
\begin{kfConnect}
\kfPair{면역}{㉠ 우리 몸에 침입하여 면역 반응을 일으키는 원인 물질}
\kfPair{항체}{㉡ 병균이 우리 몸에 들어왔을 때 이겨 내는 힘}
\kfPair{감염}{㉢ 생물을 이루는 가장 작은 단위}
\kfPair{항원}{㉣ 병균이 몸 안에 들어가 병을 일으키는 것}
\kfPair{세포}{㉤ 병균과 싸우기 위해 몸 안에서 만들어지는 물질}
\end{kfConnect}

\kfSecHeader{kfAreaVocab}{활동}{건강과 과학 관련 어휘}
\noindent\textit{\small [1-5] 뜻풀이에 알맞은 단어를 초성 힌트를 참고해 쓰세요.}\par\medskip
\begin{enumerate}[leftmargin=22pt, itemsep=6pt, topsep=4pt, label=\textbf{\arabic*.}]
\item 전염병을 예방하기 위해 미리 맞는 주사나 약: [  ㅂ  ㅅ  ]
\item 몸이나 마음에 생기는 병: [  ㅈ  ㅂ  ]
\item 세균을 죽이거나 세균이 자라지 못하게 하는 약: [  ㅎ  ㅅ  ㅈ  ]
\item 우리 몸에 침입하여 면역 반응을 일으키는 원인 물질: [  ㅎ  ㅇ  ]
\item 생물을 이루는 가장 작은 단위: [  ㅅ  ㅍ  ]
\end{enumerate}

\kfSecHeader{kfAreaVocab}{문제}{}



\kfStartConcept{개념 영역 — 유의어, 반의어}


\kfSecHeader{kfAreaConcept}{개념}{유의어, 반의어}

우리가 쓰는 단어 중에는 뜻이 서로 비슷한 짝도 있고, 뜻이 정반대인 짝도 있습니다. 이 관계를 잘 알면 글을 훨씬 더 풍부하고 정확하게 쓸 수 있답니다.

　• 유의어: 소리는 다르지만 뜻이 서로 비슷한 단어입니다.   예시: 친구 - 벗, 기쁨 - 즐거움, 걱정 - 근심, 견디다 - 참다 　• 반의어: 뜻이 서로 정반대되는 단어입니다.   예시: 크다 ↔ 작다, 빠르다 ↔ 느리다, 성공 ↔ 실패, 밀다 ↔ 당기다
\par\medskip\begin{itemize}[leftmargin=18pt, itemsep=2pt, topsep=2pt, label=\textbullet]
\item 유의어 — 가르치다·교육하다·지도하다(같은 뜻, 다른 분위기)
\item 반의어(정도) — 뜨겁다↔차갑다 / 가볍다↔무겁다
\item 반의어(관계) — 부모↔자식, 사다↔팔다(쌍을 이루는 관계)
\end{itemize}

\kfSecHeader{kfAreaConcept}{문제}{}

\begin{kfQBlock}{1}{}
\kfQStem{장점 ↔}
\begin{kfEvidence}보\end{kfEvidence}
\kfAnswerLines{1}
\end{kfQBlock}
\begin{kfQBlock}{2}{}
\kfQStem{절약 ↔}
\begin{kfEvidence}기\end{kfEvidence}
\kfAnswerLines{1}
\end{kfQBlock}
\begin{kfQBlock}{3}{}
\kfQStem{구매 ↔}
\begin{kfEvidence}>\end{kfEvidence}
\kfAnswerLines{1}
\end{kfQBlock}
\begin{kfQBlock}{4}{}
\kfQStem{만남 ↔}
\kfAnswerLines{1}
\end{kfQBlock}
\begin{kfQBlock}{5}{}
\kfQStem{입원 ↔}
\kfAnswerLines{1}
\end{kfQBlock}
\begin{kfQBlock}{1}{}
\kfQStem{다음 문장의 빈칸에 '입구'와 반대되는 뜻을 가진 낱말을 쓰시오.}
\begin{kfEvidence}이 건물의 입구는 정면에 있지만, 나가는 [　　　　　]는 건물 뒤쪽에 있다.\end{kfEvidence}
\kfAnswerNote{답안}{4}
\end{kfQBlock}
\begin{kfQBlock}{2}{}
\kfQStem{다음 문장의 흐름을 고려하여 '확대'와 반대되는 의미를 지닌 낱말을 쓰시오.}
\begin{kfEvidence}스마트폰으로 사진을 확대해서 보려고 했는데, 실수로 더 [　　　　　]해 버렸다.\end{kfEvidence}
\kfAnswerNote{답안}{4}
\end{kfQBlock}
\begin{kfQBlock}{3}{}
\kfQStem{다음 문장의 빈칸에 '단점'의 반의어를 활용하여 문장을 완성하시오.}
\begin{kfEvidence}나의 단점을 부끄러워하기보다는, 나만의 [　　　　　]을 살려서 자신감을 갖는 것이 좋다.\end{kfEvidence}
\kfAnswerNote{답안}{4}
\end{kfQBlock}

\kfSecHeader{kfAreaConcept}{활동}{유의어, 반의어}
\noindent\textit{\small [1-5] 뜻이 서로 비슷한 단어끼리 선으로 연결해 보세요.}\par\medskip
\begin{kfConnect}
\kfPair{증가}{㉠ 늘어남}
\kfPair{이유}{㉡ 근심}
\kfPair{희망}{㉢ 가난}
\kfPair{빈곤}{㉣ 까닭}
\kfPair{걱정}{㉤ 소망}
\end{kfConnect}

\kfSecHeader{kfAreaConcept}{활동}{유의어, 반의어}
\noindent\textit{\small 다음 진술이 맞으면 ○, 틀리면 × 표시하시오.}\par\medskip

\kfSecHeader{kfAreaConcept}{문제}{}



\kfStartLit{문학 영역 — 배경 지식 돋보기}


\kfSecHeader{kfAreaLiterature}{지문}{배경 지식 돋보기}
\kfPassageNote{배경 지식 돋보기}{%
이 소설은 우리나라가 아주 빠르게 변화하던 1960\textasciitilde{}70년대를 배경으로 하고 있어요. 당시 사회에 어떤 일이 있었는지 알아봅시다.

1. 무명천과 카시미론의 대결 - 무명천: 황씨 아저씨가 만들던 천연 섬유. 땀 흡수가 잘 되지만, 잘 구겨지고 투박한 느낌이 들어요. - 카시미론: 공장에서 기계로 뽑아낸 합성 섬유. 가볍고 부드러운 데다 알록달록 예쁜 색깔을 낼 수 있어 당시 사람들에게 엄청난 인기를 끌었어요. - 결과: 사람들은 촌스러운 무명천 대신 화려하고 편리한 카시미론을 선택했습니다. 이것은 '전통적인 것'이 '새로운 유행'에 밀려나는 당시의 시대 상황을 보여줍니다.

2. 산업화라는 거대한 바람 큰 공장이 들어서고 기계가 물건을 대량으로 찍어내기 시작한 것을 '산업화'라고 해요. 산업화 덕분에 나라는 발전했지만, 황씨 아저씨처럼 전통적인 방식으로 살아가던 사람들은 일자리를 잃고 가난해지기도 했어요.

3. 아저씨에게 암소란? 옛날 농촌에서 소는 단순한 가축이 아니었어요. 농사일을 돕는 든든한 일꾼이자, 집안의 가장 큰 재산이었죠. 황씨 아저씨에게 암소는 빚을 갚고 다시 일어설 수 있는 '유일한 희망'이었어요.
}


\kfSecHeader{kfAreaLiterature}{지문}{작품}
\kfPassageNote{작품}{%
　황씨 아저씨는 원래 옷감을 만드는 작은 공장 사장님이었습니다. 아저씨가 만든 하얀 무명천은 아주 인기가 좋아서, 장날마다 직접 짊어지고 나가면 금방 다 팔리곤 했습니다. 돈도 꽤 벌었지요. 아저씨는 친구 선출이에게 빌린 돈도 곧 갚을 수 있을 거라고 믿었습니다.

　그런데 세상이 갑자기 변하기 시작했습니다. 마을 근처에 커다란 공장 단지가 들어선 것입니다. \uline{㉠그곳에서는} 기계들이 ‘윙윙’ 돌아가며, 무명천보다 훨씬 부드럽고 알록달록한 ‘카시미론’이라는 새로운 옷감을 쏟아냈습니다.

　설상가상으로 황씨 아저씨네 공장에서 일하던 일꾼들도 모두 그만두겠다고 했습니다.

　“사장님, 새로 생긴 큰 공장에서는 월급을 두 배나 준대요. 우리는 \uline{㉡그리}로 갈래요.”

　황씨 아저씨는 월급을 올려줄 돈이 없었습니다. 사람들은 이제 촌스러운 무명천 대신 화려한 새 옷감을 찾았고, 아저씨의 기계들은 거미줄만 쳐진 채 멈춰 버렸습니다. 결국 유행에 밀려난 아저씨의 공장은 문을 닫아야 했습니다.

　하지만 황씨 아저씨는 포기하지 않았습니다. 친구 선출이에게 진 빚을 갚기 위해, 두 사람은 힘을 합쳐 암소 한 마리를 키우기로 했습니다.

　“이 소만 잘 키워서 송아지를 낳으면 다시 일어설 수 있어!”

　마침내 암소가 송아지를 갖게 되자, 두 사람은 빚을 갚고 다시 일어설 수 있다는 희망에 뛸 듯이 기뻐했습니다.

　어느 날, 황씨 아저씨네 집에서 일이 잘되기를 비는 작은 잔치가 열렸습니다. 그런데 아저씨의 아내가 실수로 술을 만들고 남은 찌꺼기인 술지게미를 소 여물통 근처에 두었습니다.

　배가 고팠던 암소는 고소한 냄새에 이끌려 \uline{㉢그것}을 먹어 치웠습니다. 심지어 창고 문이 열린 틈을 타, 그 안에 있던 막걸리 한 독을 몽땅 마셔버리고 말았습니다.

　“엄마! 소가 이상해! 소가 어디 갔어?”

　딸 양순이의 비명에 사람들이 뛰어나왔습니다. 얌전하던 암소가 술에 잔뜩 취해 마당에서 날뛰고 있었습니다. 소는 비틀거리고 넘어지고, 벌러덩 드러눕기를 반복했습니다.

　“아이고, 내 소! 내 소!”

　황씨 아저씨와 선출이는 너무 놀라 소를 붙잡으려 했지만, 술 취한 소의 힘을 당해낼 수가 없었습니다. 동네 사람들도 “소가 술 취해 죽는 건 처음 본다”며 발을 동동 구를 뿐이었습니다.

　사람들은 소를 살리기 위해 마당에 모닥불을 피우고, 해독에 좋다는 녹두 물을 먹였습니다. 하지만 소는 깊은 잠에 빠져 다시는 일어나지 못했습니다. 추운 겨울밤, 마당 한가운데서 타오르는 모닥불 불티만 하늘로 솟아올랐습니다.

　황씨 아저씨는 죽은 소를 붙잡고 울부짖었습니다.

　“\uline{㉣이 녀석아}, 어떻게 네가 나를 배신하니! 내가 너만 믿고 살았는데!”

　아저씨의 아내는 속으로 ‘죽은 소 배 속에 있는 송아지라도 꺼내서 남편 몸보신이나 시켜야겠다’고 생각하며 불을 더 지폈습니다. 구경하던 마을 사람들은 겉으로는 안타까워하면서도, 속으로는 ‘죽은 소 고기는 싸게 팔 테니, 고기나 실컷 먹을 수 있겠네’라고 생각했습니다.

　그렇게 황씨 아저씨의 마지막 희망은 허무하게 사라지고 말았습니다. 산업화에 밀려난 것도 서러운데, 다시 일어서려던 꿈마저 잃어버린 슬픈 겨울밤이었습니다.
}


\kfSecHeader{kfAreaLiterature}{활동}{차근차근 내용 확인}
\noindent\textit{\small [1-5] 글의 내용과 일치하면 O, 일치하지 않으면 X를 쓰세요.}\par\medskip
\begin{kfOXList}
\kfOXItem 황씨 아저씨는 원래 하얀 무명천을 만드는 공장 사장님이었다. \kfOX
\kfOXItem 암소는 건강하게 송아지를 낳은 뒤에 늙어서 죽었다. \kfOX
\kfOXItem 아저씨의 아내는 일부러 소를 취하게 하려고 술 찌꺼기를 주었다. \kfOX
\kfOXItem 마을 사람들은 소를 살리기 위해 마당에 모닥불을 피우고 녹두 물을 먹였다. \kfOX
\kfOXItem 사람들은 무명천이 촌스럽다고 생각하여 화려한 새 옷감을 더 좋아했다. \kfOX
\end{kfOXList}

\kfSecHeader{kfAreaLiterature}{활동}{나만의 생각 펼치기}
\noindent\textit{\small 다음 조건에 맞게 글을 써 보세요.}\par\medskip
\kfWritingPrompt{황씨 아저씨에게 암소는 다시 일어설 수 있게 해주는 소중한 '희망'이었습니다. 여러분에게도 힘든 순간에 힘이 되어주는 '희망'이 있나요? * 나의 보물: (　　　　) * 이유:}
\kfWriting{답안 작성}{8}
\kfWritingPrompt{모든 것을 잃고 슬퍼하는 황씨 아저씨에게 다시 힘을 낼 수 있도록 희망의 메시지를 담은 짧은 편지를 써 주세요. 황씨 아저씨께:}
\kfWriting{답안 작성}{8}

\kfSecHeader{kfAreaLiterature}{문제}{}

\begin{kfQBlock}{1}{}
\kfQStem{'㉠그곳'이 가리키는 장소는 어디인가요?}
\begin{kfChoices}
\kfChoice 황씨 아저씨네 공장
\kfChoice 마을 근처 큰 공장 단지
\kfChoice 물건을 파는 시골 장터
\kfChoice 선출이 아저씨네 집
\end{kfChoices}
\end{kfQBlock}
\begin{kfQBlock}{2}{}
\kfQStem{일꾼들이 말한 '㉡그리'는 어디를 뜻하나요?}
\begin{kfChoices}
\kfChoice 황씨 아저씨의 공장
\kfChoice 새로 생긴 큰 공장
\kfChoice 일꾼들의 고향 마을
\kfChoice 서울에 있는 백화점
\end{kfChoices}
\end{kfQBlock}
\begin{kfQBlock}{3}{}
\kfQStem{암소가 먹은 '㉢그것'은 무엇인가요?}
\begin{kfChoices}
\kfChoice 광에 있던 막걸리
\kfChoice 술을 만들고 남은 찌꺼기
\kfChoice 소가 평소에 먹던 여물
\kfChoice 아저씨가 사 온 과자
\end{kfChoices}
\end{kfQBlock}
\begin{kfQBlock}{4}{}
\kfQStem{황씨 아저씨가 부르는 '㉣이 녀석'이 가리키는 대상은 누구인가요?}
\begin{kfChoices}
\kfChoice 공장을 떠난 일꾼들
\kfChoice 돈을 빌려준 선출이
\kfChoice 아저씨가 믿었던 소
\kfChoice 술 찌꺼기를 둔 아내
\end{kfChoices}
\end{kfQBlock}
\begin{kfQBlock}{1}{}
\kfQStem{황씨 아저씨의 공장이 문을 닫게 된 진짜 이유는 무엇인가요?}
\begin{kfChoices}
\kfChoice 아저씨가 게을러서 일을 열심히 하지 않았기 때문이다.
\kfChoice 큰 공장에서 만든 화려한 새 옷감이 유행했기 때문이다.
\kfChoice 공장에 큰 불이 나서 기계가 모두 타 버렸기 때문이다.
\kfChoice 친구가 빌려준 돈을 빨리 갚으라고 재촉했기 때문이다.
\end{kfChoices}
\end{kfQBlock}
\begin{kfQBlock}{2}{}
\kfQStem{황씨 아저씨가 빚을 갚고 다시 일어서기 위해 키운 동물은 무엇인가요?}
\begin{kfChoices}
\kfChoice 꿀꿀 돼지
\kfChoice 꼬꼬 닭
\kfChoice 누렁 암소
\kfChoice 바둑 강아지
\end{kfChoices}
\end{kfQBlock}
\begin{kfQBlock}{3}{}
\kfQStem{암소에게 벌어진 슬픈 사건의 직접적인 원인은 무엇인가요?}
\begin{kfChoices}
\kfChoice 밤사이 도둑이 들어서 몰래 소를 훔쳐 갔기 때문이다.
\kfChoice 소가 술 찌꺼기와 술을 먹고 심하게 취해 버렸기 때문이다.
\kfChoice 날씨가 너무 추워서 소가 꽁꽁 얼어 버렸기 때문이다.
\kfChoice 산에서 호랑이가 내려와 소를 무섭게 물어 갔기 때문이다.
\end{kfChoices}
\end{kfQBlock}
\begin{kfQBlock}{1}{}
\kfQStem{1960\textasciitilde{}70년대, 기계를 갖춘 큰 공장이 생겨나고 사회가 빠르게 발전하던 현상을 (　　　　)(이)라고 합니다.}
\kfAnswerLines{1}
\end{kfQBlock}
\begin{kfQBlock}{2}{}
\kfQStem{황씨 아저씨네 무명천 공장이 문을 닫게 된 결정적인 이유는 새롭고 화려한 합성 섬유인 (　　　　　)이/가 유행했기 때문입니다.}
\kfAnswerLines{1}
\end{kfQBlock}
\begin{kfQBlock}{3}{}
\kfQStem{이 소설에서 '암소'는 황씨 아저씨가 빚을 갚고 재기할 수 있는 마지막 (　　　)을/를 상징합니다.}
\kfAnswerLines{1}
\end{kfQBlock}
\begin{kfQBlock}{1}{}
\kfQStem{다음은 소설의 주요 사건을 시간 순서대로 정리한 것입니다. 빈칸에 들어갈 내용을 채워 봅시다.

| 주요 내용 | 의미 | | :---- | :---- | | 황씨 아저씨, ( ① 　　　　 )에 밀려 공장 문을 닫음 | 시대적 시련 시작 | | 친구 선출이와 ( ② 　　　　 )을/를 사서 희망을 가짐 | 다시 일어설 희망 | | 암소가 실수로 ( ③ 　　　　 )을/를 먹고 쓰러짐 | 희망에 닥친 예상치 못한 비극 | | 동네 사람들 노력에도 불구하고 암소가 ( ④ 　　　　 ) | 모든 희망의 사라짐 | | 아저씨의 절망, 주변 인물의 이기적인 속마음 | 비극적 상황 심화 및 주제 암시 |}
\kfAnswerLines{1}
\end{kfQBlock}
\begin{kfQBlock}{1}{}
\kfQStem{소가 죽은 뒤, 마을 사람들의 겉모습과 속마음이 어떻게 달랐는지 쓰세요. 1) 겉으로는: (　　　　) 2) 속으로는:}
\kfAnswerNote{답안}{4}
\end{kfQBlock}


\kfStartNonlit{비문학 영역 — 우리 몸을 지키는 든든한 방패, 면역}


\kfSecHeader{kfAreaNonlit}{지문}{우리 몸을 지키는 든든한 방패, 면역}
\kfPassageNote{우리 몸을 지키는 든든한 방패, 면역}{%
　우리 눈에는 보이지 않지만, 공기 중이나 물건 표면에는 수많은 세균과 바이러스가 살고 있습니다. \uline{㉠이들은} 기회가 생기면 우리 몸속으로 들어와 감염을 일으키고 질병을 만듭니다. 하지만 우리는 매일 아프지 않고 건강하게 지냅니다. 바로 우리 몸 안에 튼튼한 성벽과 군대가 있기 때문입니다. \uline{㉡이것을} '면역'이라고 합니다.

　면역은 크게 두 단계로 우리 몸을 보호합니다. 1단계는 성벽과 같은 역할을 하는 '피부'와 '점막'입니다. 피부는 병균이 들어오지 못하게 막고, 코와 목 안의 끈적끈적한 점막은 먼지와 세균을 붙잡습니다.

　만약 병균이 \uline{㉢이 성벽을} 뚫고 몸속으로 들어오면 2단계 보호가 시작됩니다. 이때 등장하는 것이 바로 '백혈구'라는 군인입니다. 백혈구는 종류에 따라 나쁜 병균을 직접 잡아먹기도 하고, 병균을 무찌를 수 있는 '항체'라는 무기를 만들어내기도 합니다. 한번 만들어진 항체는 우리 몸에 기억되어, 나중에 같은 병균이 또 들어오면 아주 빠르고 강력하게 물리칩니다.

　그렇다면 '백신'은 무엇일까요? 백신은 병균의 힘을 아주 약하게 만들거나 죽인 것을 이용해 만든 약입니다. 백신을 맞으면 우리 몸은 마치 진짜 적이 쳐들어온 줄 알고 미리 항체라는 무기를 만들고 훈련을 합니다. \uline{㉣이것은} 가짜 적을 이용해 미리 하는 '군사 훈련'과 같습니다. 덕분에 나중에 진짜 강력한 바이러스가 들어와도, 우리 몸은 이미 훈련이 되어 있어 쉽게 이겨 낼 수 있습니다. 이처럼 면역과 백신은 우리 건강을 지키는 가장 소중한 방패입니다.
}


\kfSecHeader{kfAreaNonlit}{활동}{어휘력 쑥쑥! 단어 익히기}
\noindent\textit{\small [1-4] 다음은 글에 나온 낱말들입니다. 낱말의 뜻풀이와 알맞게 연결하세요.}\par\medskip
\begin{kfConnect}
\kfPair{면역}{㉠ 몸속에 들어온 세균이나 바이러스를 무찌를 수 있도록 백혈구가 만들어 내는 물질.}
\kfPair{백혈구}{㉡ 병원균에 대한 저항력을 갖게 하기 위해, 약하게 만든 균 따위를 이용하여 미리 면역이 생기도록 만든 의약품.}
\kfPair{항체}{㉢ 병균이 우리 몸에 들어왔을 때 이겨 내는 힘.}
\kfPair{백신}{㉣ 병원균을 잡아먹는 역할을 하며 우리 몸의 면역 작용에 중요한 역할을 하는 혈액 속의 세포.}
\end{kfConnect}

\kfSecHeader{kfAreaNonlit}{활동}{차근차근 내용 확인}
\noindent\textit{\small [1-5] 글의 내용과 일치하면 O, 일치하지 않으면 X를 쓰세요.}\par\medskip
\begin{kfOXList}
\kfOXItem 우리 몸의 1단계 면역 보호는 오직 피부를 통해서만 이루어진다. \kfOX
\kfOXItem 점막은 코와 목 안에서 끈적끈적하게 병균을 붙잡는 역할을 한다. \kfOX
\kfOXItem 백혈구는 병균을 직접 잡아먹거나 항체라는 무기를 만들어낸다. \kfOX
\kfOXItem 백신을 통해 만들어진 항체는 병균을 기억하지 못해 다시 감염된다. \kfOX
\kfOXItem 백신은 힘이 약해진 병균이 아닌, 건강한 병균을 이용해 만든 약이다. \kfOX
\end{kfOXList}

\kfSecHeader{kfAreaNonlit}{활동}{나만의 생각 펼치기}
\noindent\textit{\small [1-1] 다음 활동을 해보세요.}\par\medskip
\kfWritingPrompt{내가 생각하는 '튼튼한 성벽'과 '강한 군대'를 만들기 위한 나만의 방법(생활 습관)을 2가지 이상 써 보세요.}
\kfWriting{답안 작성}{8}

\kfSecHeader{kfAreaNonlit}{문제}{}

\begin{kfQBlock}{1}{}
\kfQStem{'㉠이들'이 가리키는 것은 무엇인가요?}
\begin{kfChoices}
\kfChoice 우리 몸을 지키는 세포
\kfChoice 공기 중에 떠다니는 먼지
\kfChoice 수많은 세균과 바이러스
\kfChoice 물건 표면에 묻은 얼룩
\end{kfChoices}
\end{kfQBlock}
\begin{kfQBlock}{2}{}
\kfQStem{'㉡이것'이 가리키는 상황은 무엇인가요?}
\begin{kfChoices}
\kfChoice 감염을 일으키는 나쁜 질병
\kfChoice 매일 아프지 않고 건강한 것
\kfChoice 우리 몸 안의 성벽과 군대
\kfChoice 공기 중에 떠다니는 바이러스
\end{kfChoices}
\end{kfQBlock}
\begin{kfQBlock}{3}{}
\kfQStem{병균이 뚫고 들어온 '㉢이 성벽'은 구체적으로 무엇을 말하나요?}
\begin{kfChoices}
\kfChoice 우리 몸의 피부와 점막
\kfChoice 혈액 속에 있는 백혈구
\kfChoice 병균을 무찌르는 항체
\kfChoice 병을 예방하는 백신
\end{kfChoices}
\end{kfQBlock}
\begin{kfQBlock}{4}{}
\kfQStem{'㉣이것'이 비유하고 있는 과정은 무엇인가요?}
\begin{kfChoices}
\kfChoice 병균이 우리 몸을 공격해서 아프게 만드는 과정
\kfChoice 백신을 통해 미리 항체를 만들고 훈련하는 과정
\kfChoice 피부와 점막이 병균의 침입을 막아내는 과정
\kfChoice 백혈구가 나쁜 병균을 직접 잡아먹는 과정
\end{kfChoices}
\end{kfQBlock}
\begin{kfQBlock}{1}{}
\kfQStem{윗글의 내용과 일치하지 \uline{않는} 것은 무엇인가요?}
\begin{kfChoices}
\kfChoice 우리 몸의 1단계 면역 보호는 오직 피부를 통해서만 이루어진다.
\kfChoice 점막은 코와 목 안에서 끈적끈적하게 병균을 붙잡는 역할을 한다.
\kfChoice 백신을 통해 만들어진 항체는 병균을 기억하지 못해 다시 감염된다.
\kfChoice 백신은 힘이 약해진 병균이 아닌, 건강한 병균을 이용해 만든 약이다.
\end{kfChoices}
\end{kfQBlock}
\begin{kfQBlock}{2}{}
\kfQStem{이 글에서 사용된 비유적 표현과 그 의미가 바르게 연결되지 \uline{않은} 것은 무엇인가요?}
\begin{kfChoices}
\kfChoice 성벽 - 백혈구
\kfChoice 군인 - 항체
\kfChoice 무기 - 피부와 점막
\kfChoice 가짜 적 - 백신
\end{kfChoices}
\end{kfQBlock}
\begin{kfQBlock}{3}{}
\kfQStem{다음 대화의 빈칸에 들어갈 말이 바르게 짝지어지지 \uline{않은} 것은 무엇인가요?}
\begin{kfEvidence}환자: “선생님, 독감 백신 주사는 왜 맞는 건가요?” 의사: “백신은 우리 몸이 미리 병균과 싸우는 연습을 하도록 도와주는 거예요. 마치 군인들이 전쟁 전에 미리 ( ① )을/를 하는 것과 같답니다.” 환자: “아, 그럼 백신을 맞으면 우리 몸속에 무기가 생기나요?” 의사: “맞아요. 우리 몸은 백신을 보고 병균을 물리칠 수 있는 ( ② )라는 무기를 미리 만들어 놓습니다.” 환자: “그럼 나중에 진짜 독감 병균이 들어오면 어떻게 되나요?” 의사: “우리 몸은 예전에 만든 무기를 ( ③ )하고 있다가, ( ④ ) 독감 바이러스가 들어오면 아주 빠르게 물리칠 수 있게 된답니다.”\end{kfEvidence}
\begin{kfChoices}
\kfChoice 훈련, 항체, 기억, 진짜
\kfChoice 휴식, 세균, 저장, 가짜
\kfChoice 공격, 백신, 삭제, 약한
\kfChoice 방어, 세포, 무시, 죽은
\end{kfChoices}
\end{kfQBlock}
\begin{kfQBlock}{4}{}
\kfQStem{1단계 보호에서 '성벽' 역할을 하며 병균이 들어오지 못하게 막는 것이 \uline{아닌} 것은 무엇인가요?}
\begin{kfChoices}
\kfChoice 백혈구와 항체
\kfChoice 피부와 점막
\kfChoice 백신과 항체
\kfChoice 피부와 백혈구
\end{kfChoices}
\end{kfQBlock}
\begin{kfQBlock}{5}{}
\kfQStem{병균이 몸속으로 들어왔을 때 싸우는 '군인' 역할을 하는 세포가 \uline{아닌} 것은 무엇인가요?}
\begin{kfChoices}
\kfChoice 적혈구
\kfChoice 백혈구
\kfChoice 피부 세포
\kfChoice 신경 세포
\end{kfChoices}
\end{kfQBlock}
\begin{kfQBlock}{6}{}
\kfQStem{'백신'을 만드는 재료로 알맞지 \uline{않은} 것은 무엇인가요?}
\begin{kfChoices}
\kfChoice 건강한 사람의 피
\kfChoice 깨끗한 물과 소금
\kfChoice 몸에 좋은 비타민과 영양제
\kfChoice 힘을 아주 약하게 만들거나 죽인 병균
\end{kfChoices}
\end{kfQBlock}
\begin{kfQBlock}{1}{}
\kfQStem{우리 몸이 병균을 막아내는 과정을 순서대로 정리했습니다. 빈칸에 들어갈 말을 <보기>에서 찾아 쓰세요. * 1단계 보호: ( ① 　　　 )와/과 점막이 병균의 침입을 막거나 붙잡는다. * (뚫리면) * 2단계 보호: ( ② 　　　 )가 출동하여 병균을 잡아먹거나 무기를 만든다. * 결과: 병균을 무찌르는 무기인 ( ③ 　　　 )이/가 만들어지고, 우리 몸에 ( ④ 　　　 )되어 나중에 더 잘 싸우게 된다.}
\begin{kfEvidence}백혈구\end{kfEvidence}
\kfAnswerLines{1}
\end{kfQBlock}


\kfStartWeekly{실력확인 영역 — 주간 실력 확인}


\noindent\textit{\small (제한시간: 30분 / 총 20문항)}\par\medskip
\par\noindent\textbf{[6-10] 다음 글을 읽고, 물음에 답하세요.}\par\medskip
\kfPassageNote{지문 6}{%
황씨 아저씨는 원래 옷감을 만드는 작은 공장 사장님이었습니다. 아저씨가 만든 하얀 무명천은 아주 인기가 좋아서, 장날마다 직접 짊어지고 나가면 금방 다 팔리곤 했습니다. 돈도 꽤 벌었지요. 아저씨는 친구 선출이에게 빌린 돈도 곧 갚을 수 있을 거라고 믿었습니다.

그런데 세상이 갑자기 변하기 시작했습니다. 마을 근처에 커다란 공장 단지가 들어선 것입니다. 그곳에서는 기계들이 ‘윙윙’ 돌아가며, 무명천보다 훨씬 부드럽고 알록달록한 ‘카시미론’이라는 새로운 옷감을 쏟아냈습니다.

설상가상으로 황씨 아저씨네 공장에서 일하던 일꾼들도 모두 그만두겠다고 했습니다.

“사장님, 새로 생긴 큰 공장에서는 월급을 두 배나 준대요. 우리는 그리로 갈래요.”

황씨 아저씨는 월급을 올려줄 돈이 없었습니다. 사람들은 이제 촌스러운 무명천 대신 화려한 새 옷감을 찾았고, 아저씨의 기계들은 거미줄만 쳐진 채 멈춰 버렸습니다. 결국 유행에 밀려난 아저씨의 공장은 문을 닫아야 했습니다.

하지만 황씨 아저씨는 포기하지 않았습니다. 친구 선출이에게 진 빚을 갚기 위해, 두 사람은 힘을 합쳐 암소 한 마리를 키우기로 했습니다.

“이 소만 잘 키워서 송아지를 낳으면 다시 일어설 수 있어!”

마침내 암소가 송아지를 갖게 되자, 두 사람은 빚을 갚고 다시 일어설 수 있다는 희망에 뛸 듯이 기뻐했습니다.

어느 날, 황씨 아저씨네 집에서 일이 잘되기를 비는 작은 잔치가 열렸습니다. 그런데 아저씨의 아내가 실수로 술을 만들고 남은 찌꺼기(술지게미)를 소 여물통 근처에 두었습니다.

배가 고팠던 암소는 고소한 냄새에 이끌려 그것을 먹어 치웠습니다. 심지어 창고 문이 열린 틈을 타, 그 안에 있던 막걸리 한 독을 몽땅 마셔버리고 말았습니다.

“엄마! 소가 이상해! 소가 어디 갔어?”

딸 양순이의 비명에 사람들이 뛰어나왔습니다. 얌전하던 암소가 술에 잔뜩 취해 마당에서 날뛰고 있었습니다. 소는 비틀거리고 넘어지고, 벌러덩 드러눕기를 반복했습니다.

“아이고, 내 소! 내 소!”

황씨 아저씨와 선출이는 너무 놀라 소를 붙잡으려 했지만, 술 취한 소의 힘을 당해낼 수가 없었습니다. 동네 사람들도 “소가 술 취해 죽는 건 처음 본다”며 발을 동동 구를 뿐이었습니다.

사람들은 소를 살리기 위해 마당에 모닥불을 피우고, 해독에 좋다는 녹두 물을 먹였습니다. 하지만 소는 깊은 잠에 빠져 다시는 일어나지 못했습니다. 추운 겨울밤, 마당 한가운데서 타오르는 모닥불 불티만 하늘로 솟아올랐습니다.

황씨 아저씨는 죽은 소를 붙잡고 울부짖었습니다.

“이 녀석아, 어떻게 네가 나를 배신하니! 내가 너만 믿고 살았는데!”

아저씨의 아내는 속으로 ‘죽은 소 배 속에 있는 송아지라도 꺼내서 남편 몸보신이나 시켜야겠다’고 생각하며 불을 더 지폈습니다. 구경하던 마을 사람들은 겉으로는 안타까워하면서도, 속으로는 ‘죽은 소 고기는 싸게 팔 테니, 고기나 실컷 먹을 수 있겠네’라고 생각했습니다.

그렇게 황씨 아저씨의 마지막 희망은 허무하게 사라지고 말았습니다. 산업화에 밀려난 것도 서러운데, 다시 일어서려던 꿈마저 잃어버린 슬픈 겨울밤이었습니다.

- 이문구, 「암소」
}
\par\noindent\textbf{[11-15] 다음 글을 읽고, 물음에 답하세요.}\par\medskip
\kfPassageNote{지문 11}{%
우리 눈에는 보이지 않지만, 공기 중이나 물건 표면에는 수많은 세균과 바이러스가 살고 있습니다. 이들은 기회가 생기면 우리 몸속으로 들어와 감염을 일으키고 질병을 만듭니다. 하지만 우리는 매일 아프지 않고 건강하게 지냅니다. 바로 우리 몸 안에 튼튼한 성벽과 군대가 있기 때문입니다. 이것을 '면역'이라고 합니다.

면역은 크게 두 단계로 우리 몸을 보호합니다. 1단계는 성벽과 같은 역할을 하는 '피부'와 '점막'입니다. 피부는 병균이 들어오지 못하게 막고, 코와 목 안의 끈적끈적한 점막은 먼지와 세균을 붙잡습니다.

만약 병균이 이 성벽을 뚫고 몸속으로 들어오면 2단계 보호가 시작됩니다. 이때 등장하는 것이 바로 '백혈구'라는 군인입니다. 백혈구는 종류에 따라 나쁜 병균을 직접 잡아먹기도 하고, 병균을 무찌를 수 있는 '항체'라는 무기를 만들어내기도 합니다. 한번 만들어진 항체는 우리 몸에 기억되어, 나중에 같은 병균이 또 들어오면 아주 빠르고 강력하게 물리칩니다.

그렇다면 '백신'은 무엇일까요? 백신은 병균의 힘을 아주 약하게 만들거나 죽인 것을 이용해 만든 약입니다. 백신을 맞으면 우리 몸은 마치 진짜 적이 쳐들어온 줄 알고 미리 항체라는 무기를 만들고 훈련을 합니다. 이것은 가짜 적을 이용해 미리 하는 '군사 훈련'과 같습니다. 덕분에 나중에 진짜 강력한 바이러스가 들어와도, 우리 몸은 이미 훈련이 되어 있어 쉽게 이겨 낼 수 있습니다. 이처럼 면역과 백신은 우리 건강을 지키는 가장 소중한 방패입니다.
}
\begin{kfQBlock}{1}{}
\kfQStem{우리 주변 어디에나 있으며 좋은 것도 나쁜 것도 있고, 항생제로 치료하는 아주 작은 생물은 무엇인가요?}
\begin{kfChoices}
\kfChoice 바이러스
\kfChoice 항원
\kfChoice 세균
\kfChoice 항체
\end{kfChoices}
\end{kfQBlock}

\begin{kfQBlock}{2}{}
\kfQStem{다음 중 '몸이나 마음에 생기는 병'을 뜻하는 단어는 무엇인가요?}
\begin{kfChoices}
\kfChoice 감염
\kfChoice 면역
\kfChoice 질병
\kfChoice 항체
\end{kfChoices}
\end{kfQBlock}

\begin{kfQBlock}{3}{}
\kfQStem{다음 단어와 그 뜻풀이가 바르게 연결되지 \uline{않은} 것은 무엇인가요?}
\begin{kfChoices}
\kfChoice 면역 - 병균이 몸에 들어왔을 때 이겨 내는 힘
\kfChoice 세포 - 생물을 이루는 가장 작은 단위
\kfChoice 항체 - 병원균에 대한 저항력을 갖게 하는 약
\kfChoice 항원 - 면역 반응을 일으키는 원인 물질
\end{kfChoices}
\end{kfQBlock}

\begin{kfQBlock}{4}{}
\kfQStem{다음 중 '감염'이라는 단어를 사용하여 문장을 만들 때, 문맥상 가장 \uline{어색한} 것은 무엇인가요?}
\begin{kfChoices}
\kfChoice 독감 바이러스에 감염되지 않도록 마스크를 썼다.
\kfChoice 강아지가 전염병에 감염되어 동물 병원에 입원했다.
\kfChoice 나는 친구의 밝고 긍정적인 웃음에 감염되었다.
\kfChoice 식중독은 상한 음식을 먹었을 때 세균에 감염되어 생긴다.
\end{kfChoices}
\end{kfQBlock}

\begin{kfQBlock}{5}{}
\kfQStem{다음 문장의 빈칸에 들어갈 알맞은 단어를 순서대로 짝지은 것은 무엇인가요?}
\begin{kfEvidence}“전염병을 막기 위해 [　　　　　]을/를 맞아 몸속에 [　　　　　]를/을 미리 만들고, 병에 걸리면 [　　　　　]를/을 사용하여 치료한다.”\end{kfEvidence}
\begin{kfChoices}
\kfChoice 백신 - 항원 - 항생제
\kfChoice 항생제 - 백신 - 항체
\kfChoice 백신 - 항체 - 항생제
\kfChoice 항체 - 항생제 - 백신
\end{kfChoices}
\end{kfQBlock}

\begin{kfQBlock}{6}{}
\kfQStem{이 소설의 시대적 배경인 1960\textasciitilde{}70년대 상황에 대한 설명으로 가장 적절한 것은 무엇인가요?}
\begin{kfChoices}
\kfChoice 전통적인 방법이 새로운 기계 문명에 밀려나던 산업화 시대였다.
\kfChoice 나라 전체가 농사에만 집중하며 소를 가장 중요한 재산으로 여겼다.
\kfChoice 큰 공장이 모두 망하고 황씨 아저씨 같은 작은 공장만 살아남았다.
\kfChoice 국민들이 소비를 줄이고 절약하여 무명천 같은 물건만 찾았다.
\end{kfChoices}
\end{kfQBlock}

\begin{kfQBlock}{7}{}
\kfQStem{황씨 아저씨가 만든 '무명천'과 새로 유행한 '카시미론'의 관계가 상징하는 것은 무엇인가요?}
\begin{kfChoices}
\kfChoice 신제품의 등장과 대중의 환영
\kfChoice 거친 느낌과 부드러운 느낌이 함께 있음
\kfChoice 오래된 가치와 새로운 가치의 대립
\kfChoice 생산자와 소비자의 협력
\end{kfChoices}
\end{kfQBlock}

\begin{kfQBlock}{8}{}
\kfQStem{황씨 아저씨가 “이 녀석아, 어떻게 네가 나를 배신하니!”라고 울부짖은 이유로 가장 알맞은 것은 무엇인가요?}
\begin{kfChoices}
\kfChoice 소가 다른 집으로 도망가버렸기 때문에
\kfChoice 소가 아저씨에게 폭력을 행사했기 때문에
\kfChoice 소가 빚을 갚을 수 있는 유일한 희망이었기 때문에
\kfChoice 소가 아저씨의 아내에게 심한 상처를 입혔기 때문에
\end{kfChoices}
\end{kfQBlock}

\begin{kfQBlock}{9}{}
\kfQStem{다음 중 소설의 내용과 일치하지 \uline{않은} 것은 무엇인가요?}
\begin{kfChoices}
\kfChoice 황씨 아저씨의 아내는 실수로 술 찌꺼기를 소 여물통 근처에 두었다.
\kfChoice 아저씨는 친구 선출이에게 빌린 돈을 갚기 위해 암소를 키우기 시작했다.
\kfChoice 마을 사람들은 겉으로는 아저씨를 걱정했지만, 속으로는 소고기를 기대했다.
\kfChoice 암소가 송아지를 낳은 직후에 술에 취해 죽어 황씨 아저씨가 큰 절망에 빠졌다.
\end{kfChoices}
\end{kfQBlock}

\begin{kfQBlock}{10}{}
\kfQStem{이 글을 읽고 아래 <보기>를 참고하여 황씨 아저씨의 상황을 이해한 것으로 가장 적절하지 \uline{않은} 것은 무엇인가요?}
\begin{kfEvidence}<보기> 　산업화 시대의 '거대한 바람'은 개인의 노력이나 의지와 상관없이 많은 것을 변화시키고 파괴합니다. 이 소설에서 황씨 아저씨는 시대의 흐름에 저항하지 못하고 끝내 좌절하는 인물로 그려집니다.\end{kfEvidence}
\begin{kfChoices}
\kfChoice 공장의 폐업은 아저씨의 게으름이 아닌 시대의 흐름 탓으로 볼 수 있다.
\kfChoice 카시미론의 유행은 개인이 막을 수 없는 거대한 시대적 변화를 보여준다.
\kfChoice 암소를 키운 것은 산업화에 맞서 다시 일어서려 했던 마지막 희망이었다.
\kfChoice 암소의 죽음은 아저씨가 산업화의 시련을 결국 이겨냈음을 의미한다.
\end{kfChoices}
\end{kfQBlock}

\begin{kfQBlock}{11}{}
\kfQStem{윗글의 내용과 일치하지 \uline{않은} 것은 무엇인가요?}
\begin{kfChoices}
\kfChoice 백신은 가짜 적을 이용해 미리 하는 군사 훈련과 같습니다.
\kfChoice 피부와 점막은 병균이 몸속에 들어오지 못하게 성벽처럼 막습니다.
\kfChoice 백혈구는 항체라는 무기를 만들어 내거나 병균을 직접 잡아먹습니다.
\kfChoice 항체는 병균을 기억하지 못하여 같은 병균이 또 들어오면 새로 훈련해야 합니다.
\end{kfChoices}
\end{kfQBlock}

\begin{kfQBlock}{12}{}
\kfQStem{우리 몸의 면역 활동을 설명하기 위해 사용한 비유적 표현과, 그 비유가 가리키는 대상이 바르게 짝지어진 것은 무엇인가요?}
\begin{kfChoices}
\kfChoice 성벽 - 백혈구
\kfChoice 군인 - 항체
\kfChoice 무기 - 항체
\kfChoice 가짜 적 - 진짜 바이러스
\end{kfChoices}
\end{kfQBlock}

\begin{kfQBlock}{13}{}
\kfQStem{우리 몸의 면역 활동 2단계에서 일어나는 일들을 모두 고른 것은 무엇인가요?}
\begin{kfEvidence}가. 피부와 점막이 병균의 침입을 막습니다. 나. 백혈구가 출동하여 병균을 잡아먹습니다. 다. 병균과 싸울 수 있는 항체가 만들어집니다. 라. 코와 목 안의 점막이 세균을 끈적하게 붙잡습니다.\end{kfEvidence}
\begin{kfChoices}
\kfChoice 가, 라
\kfChoice 나, 다
\kfChoice 가, 나, 다
\kfChoice 나, 다, 라
\end{kfChoices}
\end{kfQBlock}

\begin{kfQBlock}{14}{}
\kfQStem{홍역을 앓은 사람이 홍역에 다시 잘 걸리지 않는 근본적인 이유로 가장 적절한 것은 무엇인가요?}
\begin{kfChoices}
\kfChoice 홍역이 약한 병이라 쉽게 이겨낼 수 있기 때문에
\kfChoice 홍역을 앓는 동안 피부와 점막이 두꺼워졌기 때문에
\kfChoice 홍역 항체가 몸속에 기억되어 있다가 빠르게 싸우기 때문에
\kfChoice 홍역 병균이 다른 바이러스와 싸우느라 힘이 약해졌기 때문에
\end{kfChoices}
\end{kfQBlock}

\begin{kfQBlock}{15}{}
\kfQStem{이 글의 내용을 바탕으로 볼 때, 건강한 면역 체계를 유지하는 방법으로 가장 거리가 \uline{먼} 것은 무엇인가요?}
\begin{kfChoices}
\kfChoice 손을 자주 씻어 피부와 점막을 청결하게 유지한다.
\kfChoice 예방 접종을 통해 몸 안에 항체를 미리 만들도록 유도한다.
\kfChoice 질병에 걸렸을 때 즉시 항생제를 사용하여 세균을 완전히 없앤다.
\kfChoice 평소 규칙적인 생활로 백혈구가 제 역할을 할 수 있도록 관리한다.
\end{kfChoices}
\end{kfQBlock}

\begin{kfQBlock}{16}{}
\kfQStem{다음 중 유의어에 대한 설명으로 가장 적절한 것은 무엇인가요?}
\begin{kfChoices}
\kfChoice 소리는 같지만 뜻이 서로 비슷한 단어입니다.
\kfChoice 뜻이 서로 정반대되는 관계의 단어입니다.
\kfChoice 소리는 다르지만 뜻이 서로 비슷한 단어입니다.
\kfChoice 하나의 단어가 여러 가지 뜻을 가지고 있는 것입니다.
\end{kfChoices}
\end{kfQBlock}

\begin{kfQBlock}{17}{}
\kfQStem{다음 중 서로 유의어 관계인 단어끼리 바르게 짝지어진 것은 무엇인가요?}
\begin{kfChoices}
\kfChoice 증가 - 줄어듦
\kfChoice 이유 - 결과
\kfChoice 빈곤 - 가난
\kfChoice 희망 - 절망
\end{kfChoices}
\end{kfQBlock}

\begin{kfQBlock}{18}{}
\kfQStem{다음 중 '성공과 실패'처럼, 두 단어의 관계가 반의어인 것은 무엇인가요?}
\begin{kfChoices}
\kfChoice 친구와 벗
\kfChoice 기쁨과 즐거움
\kfChoice 빠르다와 느리다
\kfChoice 걱정과 근심
\end{kfChoices}
\end{kfQBlock}

\begin{kfQBlock}{19}{}
\kfQStem{빈칸에 들어갈 말이 '단점'의 반대말인 문장은 무엇인가요?}
\begin{kfChoices}
\kfChoice 나가는 [　　　　　]는 건물 뒤쪽에 있다.
\kfChoice 실수로 사진을 더 [　　　　　]해 버렸다.
\kfChoice 나의 단점을 부끄러워하기보다는, 나만의 [　　　　　]을 살려야 한다.
\kfChoice 나는 약속 시간에 늦은 것을 [　　　　　]한다.
\end{kfChoices}
\end{kfQBlock}

\begin{kfQBlock}{20}{}
\kfQStem{다음 중 단어와 그 반의어가 바르게 짝지어지지 \uline{않은} 것은 무엇인가요?}
\begin{kfChoices}
\kfChoice 장점 ↔ 단점
\kfChoice 절약 ↔ 판매
\kfChoice 구매 ↔ 판매
\kfChoice 만남 ↔ 이별
\end{kfChoices}
\end{kfQBlock}




\clearpage

\kfChapterCover{4}{프레게2 (챕터4)}{Chapter 4}{이번 챕터의 학습 내용을 시작합니다.}

\kfStartVocab{어휘 영역 — 논리와 주장 관련 어휘}


\kfSecHeader{kfAreaVocab}{문제}{}

\begin{kfQBlock}{1}{}
\kfQStem{다음 중 '객관적'과 뜻이 비슷하지 않은 단어는 무엇인가요?}
\begin{kfChoices}
\kfChoice 사실적
\kfChoice 주관적
\kfChoice 합리적
\kfChoice 분석적
\end{kfChoices}
\end{kfQBlock}
\begin{kfQBlock}{2}{}
\kfQStem{다음 중 '사실'과 뜻이 같은 부류가 아닌 단어는 무엇인가요?}
\begin{kfChoices}
\kfChoice 진실
\kfChoice 실제
\kfChoice 의견
\kfChoice 진상
\end{kfChoices}
\end{kfQBlock}
\begin{kfQBlock}{1}{}
\kfQStem{근거}
\kfAnswerNote{답안}{4}
\end{kfQBlock}
\begin{kfQBlock}{2}{}
\kfQStem{주장}
\kfAnswerNote{답안}{4}
\end{kfQBlock}

\kfSecHeader{kfAreaVocab}{활동}{논리와 주장 관련 어휘}
\noindent\textit{\small [1-5] <보기>에서 알맞은 단어를 골라 문장을 완성하세요.}\par\medskip
\begin{kfBlankList}
\kfBlankItem 토론할 때는 자신의 주장만 내세우지 말고 알맞은 [　　　　　]을/를 제시해야 한다.
\begin{kfEvidence}보\end{kfEvidence}
\kfBlankItem 재판에서는 피고인의 죄를 [　　　　　]할 수 있는 증거가 필요하다.
\begin{kfEvidence}기\end{kfEvidence}
\kfBlankItem 신문 기사는 글쓴이의 상상보다는 확인된 [　　　　　]을/를 중심으로 써야 한다.
\begin{kfEvidence}>\end{kfEvidence}
\kfBlankItem “이 영화는 정말 재미있다”는 말은 개인의 느낌을 담은 [　　　　　]인 생각이다.
\begin{kfEvidence} \end{kfEvidence}
\kfBlankItem 우리는 [　　　　　]인 시각으로 사건을 바라보고 판단해야 한다.
\begin{kfEvidence}　\end{kfEvidence}
\end{kfBlankList}

\kfSecHeader{kfAreaVocab}{활동}{논리와 주장 관련 어휘}
\noindent\textit{\small [1-5] 단어의 정확한 의미를 찾아 선으로 연결해 보세요.}\par\medskip
\begin{kfConnect}
\kfPair{사실}{㉠ 어떤 증거를 내세워 증명함}
\kfPair{의견}{㉡ 실제로 있었던 일이나 현재에 있는 일}
\kfPair{근거}{㉢ 어떤 대상이나 일에 대해 가지는 생각}
\kfPair{입증}{㉣ 자신의 견해나 관점에 따라 판단하는 것}
\kfPair{주관적}{㉤ 어떤 의견이나 주장의 바탕이 되는 까닭이나 이유}
\end{kfConnect}

\kfSecHeader{kfAreaVocab}{활동}{논리와 주장 관련 어휘}
\noindent\textit{\small [1-5] 뜻풀이에 알맞은 단어를 초성 힌트를 참고해 쓰세요.}\par\medskip
\begin{enumerate}[leftmargin=22pt, itemsep=6pt, topsep=4pt, label=\textbf{\arabic*.}]
\item 말이나 글이 앞뒤가 잘 맞고 이치에 닿도록 생각의 순서와 규칙을 지키는 법: [  ㄴ  ㄹ  ]
\item 자신의 의견이나 주의를 굳게 내세움: [  ㅈ  ㅈ  ]
\item 사물의 이치에 맞아 올바른 성질: [  ㅌ  ㄷ  ㅅ  ]
\item 실제로 있었던 일이나 현재에 있는 일: [  ㅅ  ㅅ  ]
\item 남의 잘못이나 흠을 잡아 나쁘게 말함: [  ㅂ  ㄴ  ]
\end{enumerate}


\kfStartConcept{개념 영역 — 문장 성분 (주어, 서술어)}


\kfSecHeader{kfAreaConcept}{개념}{문장 성분 (주어, 서술어)}

문장을 이루는 가장 중요한 두 가지 재료, '주어'와 '서술어'를 알아봅시다.

1. 주어 　• 문장의 주인입니다. “누가”, “무엇이”에 해당합니다. 　• 보통 단어 뒤에 '이/가', '은/는', '께서' 등이 붙습니다. 　• 예시: 철수가 밥을 먹는다. / 꽃이 피었다.

2. 서술어 　• 주어를 설명하는 말입니다. “어찌하다(동작)”, “어떠하다(상태)”, “무엇이다”에 해당합니다. 　• 보통 문장의 맨 끝에 오며 '\textasciitilde{}다'로 끝납니다. 　• 예시: 철수가 밥을 먹는다. / 꽃이 예쁘다. / 나는 학생이다.
\par\medskip\begin{itemize}[leftmargin=18pt, itemsep=2pt, topsep=2pt, label=\textbullet]
\item 주어(이/가) — '강아지가 짖는다'(누가→강아지가)
\item 주어(께서) — '할머니께서 오신다'(높임)
\item 서술어(어찌하다·어떠하다·무엇이다) — '눈이 내린다·차다·축복이다'
\end{itemize}

\kfSecHeader{kfAreaConcept}{문제}{}

\begin{kfQBlock}{1}{}
\kfQStem{다음 빈칸에 들어갈 주어로 알맞지 \uline{않은} 것은 무엇인가요?}
\begin{kfEvidence}[　　　　　]이/가 운동장에서 달립니다.\end{kfEvidence}
\begin{kfChoices}
\kfChoice 누나
\kfChoice 형
\kfChoice 친구
\kfChoice 책상
\end{kfChoices}
\end{kfQBlock}
\begin{kfQBlock}{2}{}
\kfQStem{빈칸에 들어갈 주어로 알맞지 \uline{않은} 것은 무엇인가요?}
\begin{kfEvidence}[　　　　　]은/는 맵지만 맛있습니다.\end{kfEvidence}
\begin{kfChoices}
\kfChoice 사탕
\kfChoice 떡볶이
\kfChoice 우유
\kfChoice 빵
\end{kfChoices}
\end{kfQBlock}
\begin{kfQBlock}{3}{}
\kfQStem{빈칸에 알맞은 주어는 무엇인가요?}
\begin{kfEvidence}[　　　　　]께서 용돈을 주셨습니다.\end{kfEvidence}
\begin{kfChoices}
\kfChoice 친구
\kfChoice 동생
\kfChoice 할머니
\kfChoice 강아지
\end{kfChoices}
\end{kfQBlock}
\begin{kfQBlock}{4}{}
\kfQStem{빈칸에 알맞은 주어는 무엇인가요?}
\begin{kfEvidence}[　　　　　]이/가 주룩주룩 내립니다.\end{kfEvidence}
\begin{kfChoices}
\kfChoice 바람
\kfChoice 비
\kfChoice 눈
\kfChoice 햇빛
\end{kfChoices}
\end{kfQBlock}
\begin{kfQBlock}{5}{}
\kfQStem{다음 빈칸에 들어갈 서술어로 알맞지 \uline{않은} 것은 무엇인가요?}
\begin{kfEvidence}친구들과 놀이공원에 가니 정말 [　　　　　].\end{kfEvidence}
\begin{kfChoices}
\kfChoice 신난다
\kfChoice 즐겁다
\kfChoice 행복하다
\kfChoice 차갑다
\end{kfChoices}
\end{kfQBlock}
\begin{kfQBlock}{6}{}
\kfQStem{빈칸에 알맞은 서술어는 무엇인가요?}
\begin{kfEvidence}우리 아빠는 [　　　　　].\end{kfEvidence}
\begin{kfChoices}
\kfChoice 회사원이다
\kfChoice 예쁘다
\kfChoice 달린다
\kfChoice 내린다
\end{kfChoices}
\end{kfQBlock}
\begin{kfQBlock}{7}{}
\kfQStem{빈칸에 알맞은 서술어는 무엇인가요?}
\begin{kfEvidence}토끼가 깡충깡충 [　　　　　].\end{kfEvidence}
\begin{kfChoices}
\kfChoice 먹는다
\kfChoice 잔다
\kfChoice 뛴다
\kfChoice 운다
\end{kfChoices}
\end{kfQBlock}
\begin{kfQBlock}{8}{}
\kfQStem{빈칸에 알맞은 서술어는 무엇인가요?}
\begin{kfEvidence}도서관에서는 조용히 [　　　　　].\end{kfEvidence}
\begin{kfChoices}
\kfChoice 뛰어야 한다
\kfChoice 떠들어야 한다
\kfChoice 읽어야 한다
\kfChoice 놀아야 한다
\end{kfChoices}
\end{kfQBlock}
\begin{kfQBlock}{1}{}
\kfQStem{문장에서 동작이나 상태의 주인이 되는 말을 [　　　　　]라고 합니다.}
\kfAnswerLines{1}
\end{kfQBlock}
\begin{kfQBlock}{2}{}
\kfQStem{문장에서 주어의 동작이나 상태를 풀이하는 말을 [　　　　　]라고 합니다.}
\kfAnswerLines{1}
\end{kfQBlock}
\begin{kfQBlock}{3}{}
\kfQStem{주어 뒤에는 주로 조사 '은, 는, 이, [　　　　　]'가 붙습니다.}
\kfAnswerLines{1}
\end{kfQBlock}
\begin{kfQBlock}{1}{}
\kfQStem{내 꿈은 훌륭한 요리사가 되고 싶다.}
\kfAnswerNote{답안}{4}
\end{kfQBlock}
\begin{kfQBlock}{2}{}
\kfQStem{내 취미는 피아노를 친다.}
\kfAnswerNote{답안}{4}
\end{kfQBlock}
\begin{kfQBlock}{3}{}
\kfQStem{내가 지각한 이유는 늦잠을 잤다.}
\kfAnswerNote{답안}{4}
\end{kfQBlock}
\begin{kfQBlock}{4}{}
\kfQStem{이 꽃의 특징은 향기가 아주 좋다.}
\kfAnswerNote{답안}{4}
\end{kfQBlock}
\begin{kfQBlock}{5}{}
\kfQStem{내 생각은 우리가 서로 도와야 한다.}
\kfAnswerNote{답안}{4}
\end{kfQBlock}

\kfSecHeader{kfAreaConcept}{활동}{문장 성분 (주어, 서술어)}
\noindent\textit{\small [1-10] 다음 문장에서 주어 또는 서술어를 찾아보세요.}\par\medskip
\begin{kfBlankList}
\kfBlankItem 아기가 방긋 웃는다. (주어)
\kfBlankItem 새가 하늘을 날아간다. (주어)
\kfBlankItem 맛있는 피자가 식탁 위에 있다. (주어)
\kfBlankItem 우리 학교 운동장은 매우 넓다. (주어)
\kfBlankItem 어머니께서 시장에 가셨다. (주어)
\kfBlankItem 하늘이 매우 파랗다. (서술어)
\kfBlankItem 선생님께서 동화책을 읽어 주신다. (서술어)
\kfBlankItem 내 동생은 초등학생이다. (서술어)
\kfBlankItem 나는 어제 공원에서 놀았다. (서술어)
\kfBlankItem 칠판에 글씨가 가득하다. (서술어)
\end{kfBlankList}

\kfSecHeader{kfAreaConcept}{활동}{문장 성분 (주어, 서술어)}
\noindent\textit{\small 다음 진술이 맞으면 ○, 틀리면 × 표시하시오.}\par\medskip


\kfStartLit{문학 영역 — 배경 지식 돋보기}


\kfSecHeader{kfAreaLiterature}{지문}{배경 지식 돋보기}
\kfPassageNote{배경 지식 돋보기}{%
토끼전이 가진 문학적 특징과 배경 지식을 확인하며 이야기를 더 깊이 이해해 봅시다.

이야기꾼이 들려주던 노래가 소설이 되었어요!

토끼전은 원래 소리꾼이 북을 치며 노래하던 '판소리'였어요. 이 노래가 아주 인기가 많아서 나중에 글로 적혀 소설이 되었는데, 이를 '판소리계 소설'이라고 불러요.

이 이야기에는 두 종류의 인물이 등장해요.

1. 용왕: 힘은 아주 세지만, 병을 고치려는 욕심 때문에 판단을 흐리는 어리석은 지배층(권력자)을 상징해요. 2. 토끼: 힘은 약하고 가진 것은 없지만, 지혜와 꾀가 뛰어나서 강한 용왕을 속이는 백성(서민)을 상징해요.

옛날 사람들은 힘든 생활 속에서, 토끼가 꾀를 내어 무서운 용왕을 속이고 위기를 탈출하는 모습을 보면서 통쾌함을 느꼈답니다.
}


\kfSecHeader{kfAreaLiterature}{지문}{작품}
\kfPassageNote{작품}{%
[앞부분 줄거리] 용왕의 병을 고치기 위해 별주부는 육지로 나와 토끼를 꾀어 용궁으로 데려간다. 용왕이 토끼의 배를 갈라 간을 꺼내려 하자, 토끼는 꾀를 내어 말하기 시작한다.

　용왕이 분부하였다.

　“토끼를 묶어라! 내 병이 위중하니 빨리 배를 갈라 간을 내오너라!”

　나졸들이 우르르 달려들어 토끼를 묶고 배를 가르려 하니, 토끼가 눈을 끔적이며 급히 여쭈었다.

　“전하, 잠깐 제 말을 들어보십시오. 세상의 모든 짐승은 간이 배 속에 붙어 있으나, \uline{저}만은 간을 넣었다 뺐다 할 수 있사옵니다.”

　용왕이 기가 막혀 물었다.

　“네 말이 당치 않도다. 어찌 간을 배 밖에 두고 다닌단 말이냐?”

　토끼가 태연하게 대답했다.

　“제가 사는 육지는 맹수들이 많아 간을 배 속에 넣고 다니면 썩거나 터지기 쉽습니다. 그래서 평소에는 깊은 산속 푸른 바위 밑에 감추어 두었다가, 필요할 때만 배 속에 넣어 다니옵니다. 이번에는 별주부님이 너무 급하게 재촉하는 바람에 미처 챙기지 못하고 그냥 왔나이다.”

　용왕이 \uline{그 말을} 듣고 보니 묘하게 그럴듯하였다.

　“그 말이 정말이냐?”

　“제가 어찌 감히 용왕님을 속이겠습니까? 저를 다시 육지로 보내주시면 간을 가져와 바치겠나이다. 만약 배를 갈랐는데 간이 없으면, 저는 헛된 죽음을 당하고 전하의 병도 고치지 못할 것이니 누구를 원망하시겠습니까?”

　용왕은 잠시 고민하다가 토끼의 말을 믿기로 하였다.
}


\kfSecHeader{kfAreaLiterature}{활동}{배경 지식 돋보기}
\noindent\textit{\small [1-2] 다음 인물이 상징하는 것을 바르게 선으로 연결해 보세요.}\par\medskip
\begin{kfConnect}
\kfPair{토끼}{㉡ 힘은 없지만 지혜로운 백성 (서민)}
\kfPair{용왕}{㉠ 힘은 세지만 어리석은 지배층 (권력자)}
\end{kfConnect}

\kfSecHeader{kfAreaLiterature}{활동}{작품 속 숨은 뜻 찾기}
\noindent\textit{\small [1-4] 다음 중 지문에 나타난 '사실'에는 O, 토끼가 꾸며낸 '거짓 주장'에는 X 하세요.}\par\medskip
\begin{kfOXList}
\kfOXItem 용왕은 병이 들어서 토끼의 간이 필요하다. \kfOX
\kfOXItem 토끼는 간을 넣었다 뺐다 할 수 있는 신비한 동물이다. \kfOX
\kfOXItem 토끼는 지금 용궁에 잡혀 와 있다. \kfOX
\kfOXItem 토끼의 간은 육지의 바위 밑에 숨겨져 있다. \kfOX
\end{kfOXList}

\kfSecHeader{kfAreaLiterature}{활동}{차근차근 내용 확인}
\noindent\textit{\small [1-6] 글의 내용과 일치하면 O, 일치하지 않으면 X를 쓰세요.}\par\medskip
\begin{kfOXList}
\kfOXItem 용왕은 토끼의 병을 낫게 해주려고 배를 가르라고 했다. \kfOX
\kfOXItem 토끼는 겁을 먹었지만, 곧바로 아무렇지 않은 듯 꾀를 내어 대답했다. \kfOX
\kfOXItem 토끼는 육지에 맹수들이 많아 간을 배 속에 넣고 다니면 위험하다고 말했다. \kfOX
\kfOXItem 토끼는 자신의 간을 깊은 산속 푸른 바위 밑에 숨겨두고 왔다고 했다. \kfOX
\kfOXItem 토끼는 만약 배를 갈랐는데 간이 없으면 용왕의 병도 고치지 못할 것이라고 경고했다. \kfOX
\kfOXItem 용왕은 토끼의 말을 끝내 믿지 않고 배를 갈랐다. \kfOX
\end{kfOXList}

\kfSecHeader{kfAreaLiterature}{활동}{나만의 생각 펼치기}
\noindent\textit{\small 다음 조건에 맞게 글을 써 보세요.}\par\medskip
\kfWritingPrompt{토끼처럼 위험하거나 곤란한 상황에서 재치를 발휘해 문제를 해결한 경험이 있나요? 곤란했던 상황: \_\_\_\_\_\_\_\_\_\_\_\_\_\_\_\_\_\_\_\_\_\_\_\_\_\_\_\_\_\_\_\_\_\_\_\_\_\_\_\_\_\_\_\_\_\_\_\_\_\_\_\_\_\_\_\_\_\_\_\_\_\_\_\_\_ 나의 해결 방법: \_\_\_\_\_\_\_\_\_\_\_\_\_\_\_\_\_\_\_\_\_\_\_\_\_\_\_\_\_\_\_\_\_\_\_\_\_\_\_\_\_\_\_\_\_\_\_\_\_\_\_\_\_\_\_\_\_\_\_\_\_\_\_\_\_}
\kfWriting{답안 작성}{8}
\kfWritingPrompt{만약 용왕이 토끼의 말을 믿지 않고 “말도 안 되는 소리 하지 마라!”라고 했다면, 토끼는 어떤 꾀를 내었을까요? 토끼가 되어 위기를 탈출할 새로운 방법을 상상해서 써 보세요. 나(토끼)의 새로운 작전: \_\_\_\_\_\_\_\_\_\_\_\_\_\_\_\_\_\_\_\_\_\_\_\_\_\_\_\_\_\_\_\_\_\_\_\_\_\_\_\_\_\_\_\_\_\_\_\_\_\_\_\_\_\_\_\_ 용왕에게 할 말: \_\_\_\_\_\_\_\_\_\_\_\_\_\_\_\_\_\_\_\_\_\_\_\_\_\_\_\_\_\_\_\_\_\_\_\_\_\_\_\_\_\_\_\_\_\_\_\_\_\_\_\_\_\_\_\_\_\_\_\_\_\_\_\_\_}
\kfWriting{답안 작성}{8}

\kfSecHeader{kfAreaLiterature}{문제}{}

\begin{kfQBlock}{1}{}
\kfQStem{'㉠저'가 가리키는 대상은 누구인가요?}
\begin{kfChoices}
\kfChoice 용왕
\kfChoice 별주부
\kfChoice 토끼
\kfChoice 나졸
\end{kfChoices}
\end{kfQBlock}
\begin{kfQBlock}{2}{}
\kfQStem{'㉡그 말'이 가리키는 내용으로 적절하지 \uline{않은} 것은 무엇인가요?}
\begin{kfChoices}
\kfChoice 용왕이 신하들에게 호통치며 토끼를 단단히 묶으라고 분부한 것
\kfChoice 용왕이 토끼의 거짓말을 의심 없이 그대로 끝까지 믿기로 한 것
\kfChoice 토끼가 간을 넣었다 뺐다 할 수 있고, 육지에 두고 왔다는 것
\kfChoice 별주부가 토끼에게 어서 빨리 가야 한다고 급하게 재촉한 것
\end{kfChoices}
\end{kfQBlock}
\begin{kfQBlock}{3}{}
\kfQStem{용왕은 토끼의 말이 말도 안 되는 거짓말인데도 결국 믿고 맙니다. 용왕이 올바른 판단을 하지 못한 이유는 무엇일까요?}
\begin{kfChoices}
\kfChoice 토끼의 처지가 너무나 불쌍해 보여서 동정심이 깊이 생겼기 때문에
\kfChoice 자신의 병을 고치고 싶다는 욕심이 너무 커서 의심을 멈췄기 때문에
\kfChoice 토끼가 용왕보다 힘이 더 세서 함부로 거스를 수 없었기 때문에
\kfChoice 토끼의 말을 듣고 보니 묘하게 설득력이 없었기 때문에
\end{kfChoices}
\end{kfQBlock}
\begin{kfQBlock}{1}{}
\kfQStem{용왕이 토끼를 잡아 배를 가르려고 한 이유는 무엇인가요?}
\begin{kfChoices}
\kfChoice 토끼가 용궁의 보물을 훔쳐서
\kfChoice 토끼의 간으로 자신의 병을 고치기 위해서
\kfChoice 토끼가 별주부를 괴롭혀서
\kfChoice 육지의 이야기를 듣고 싶어서
\end{kfChoices}
\end{kfQBlock}
\begin{kfQBlock}{2}{}
\kfQStem{토끼가 “배를 갈라라”는 용왕의 명령을 듣고 보인 반응은 무엇인가요?}
\begin{kfChoices}
\kfChoice 너무 무서워서 그저 엉엉 울기만 했다.
\kfChoice 달려드는 나졸들과 끝까지 힘껏 싸웠다.
\kfChoice 당황하지 않고 차분하게 꾀를 내었다.
\kfChoice 자기를 데려온 별주부에게 살려달라고 빌었다.
\end{kfChoices}
\end{kfQBlock}
\begin{kfQBlock}{3}{}
\kfQStem{토끼가 위기를 모면하기 위해 한 거짓말(주장)은 무엇인가요?}
\begin{kfChoices}
\kfChoice “나는 토끼가 아니라 호랑이입니다.”
\kfChoice “내 간은 너무 작아서 약이 되지 않습니다.”
\kfChoice “내 간은 배 속에 없고 육지에 두고 왔습니다.”
\kfChoice “별주부가 내 간을 이미 먹어버렸습니다.”
\end{kfChoices}
\end{kfQBlock}
\begin{kfQBlock}{4}{}
\kfQStem{토끼는 왜 간을 가져오지 못했다고 핑계를 댔나요?}
\begin{kfChoices}
\kfChoice “길을 잃어버려 헤매다가 잃어버렸습니다.”
\kfChoice “별주부가 너무 급하게 재촉해서 챙기지 못했습니다.”
\kfChoice “오는 길에 다른 동물에게 도둑맞았습니다.”
\kfChoice “용왕님께 드릴 선물로 포장해 두느라 늦었습니다.”
\end{kfChoices}
\end{kfQBlock}
\begin{kfQBlock}{5}{}
\kfQStem{토끼가 “만약 배를 갈랐는데 간이 없으면...”이라고 말한 이유는 무엇인가요?}
\begin{kfChoices}
\kfChoice 용왕을 겁주어 자신을 풀어주게 하려고
\kfChoice 자신의 말이 진짜라고 믿게 하려고
\kfChoice 배를 가르는 것이 아플까 봐 걱정되어서
\kfChoice 용왕의 병을 걱정하는 척하며 시간을 끌려고
\end{kfChoices}
\end{kfQBlock}
\begin{kfQBlock}{6}{}
\kfQStem{용왕이 토끼의 말을 믿고 내린 '판단'은 무엇인가요?}
\begin{kfChoices}
\kfChoice 토끼를 감옥에 가두기로 했다.
\kfChoice 토끼를 바로 죽이기로 했다.
\kfChoice 별주부를 벌주기로 했다.
\kfChoice 토끼를 다시 육지로 보내주기로 했다.
\end{kfChoices}
\end{kfQBlock}
\begin{kfQBlock}{7}{}
\kfQStem{이 글을 통해 알 수 있는 토끼의 성격은 무엇인가요?}
\begin{kfChoices}
\kfChoice 겁이 많고 소심하다.
\kfChoice 힘이 세고 용감하다.
\kfChoice 꾀가 많고 임기응변에 강하다.
\kfChoice 정직하고 거짓말을 못 한다.
\end{kfChoices}
\end{kfQBlock}
\begin{kfQBlock}{1}{}
\kfQStem{토끼전처럼 원래는 소리꾼이 부르던 (　　　　)였는데, 나중에 글로 적혀 소설이 된 것을 '(　　　　)계 소설'이라고 합니다.}
\kfAnswerLines{1}
\end{kfQBlock}
\begin{kfQBlock}{1}{}
\kfQStem{토끼의 주장: “제 간은 지금 배 속에 없고, 육지에 두고 왔습니다.” 근거 1: “육지에는 (　　　　)이/가 많아서 배 속에 넣고 다니면 위험하기 때문입니다.” 근거 2: “평소에는 깊은 산속 (　　　　　) 밑에 잘 감추어 둡니다.”}
\kfAnswerLines{1}
\end{kfQBlock}
\begin{kfQBlock}{1}{}
\kfQStem{용왕의 병을 고치기 위해 (　　　　)가 토끼를 꾀어 용궁으로 데려간다.}
\kfAnswerLines{1}
\end{kfQBlock}
\begin{kfQBlock}{2}{}
\kfQStem{용왕이 토끼의 (　　　　)을/를 꺼내려 하자, 나졸들이 토끼를 묶는다.}
\kfAnswerLines{1}
\end{kfQBlock}
\begin{kfQBlock}{3}{}
\kfQStem{토끼가 간을 넣었다 뺐다 할 수 있으며, 간을 (　　　　)에 두고 왔다는 꾀를 낸다.}
\kfAnswerLines{1}
\end{kfQBlock}
\begin{kfQBlock}{4}{}
\kfQStem{용왕은 토끼의 말을 믿고, 토끼를 다시 육지로 (　　　　) 주기로 한다.}
\kfAnswerLines{1}
\end{kfQBlock}
\begin{kfQBlock}{1}{}
\kfQStem{Q1. 간은 신체 기관입니다. 필요할 때마다 몸에서 떼어냈다가 다시 붙이는 것이 가능한가요? : ( 예 / 아니요 ) Q2. 그렇다면 토끼의 주장은 타당한가요? 알맞은 말에 O표 하고, 그 이유를 완성해 보세요. : ( 타당하다 / 타당하지 않다 ) : 왜냐하면 \_\_\_\_\_\_\_\_\_\_\_\_\_\_\_\_\_\_\_\_\_\_\_\_\_\_\_\_\_\_\_\_\_\_\_\_\_\_\_\_\_\_\_\_\_\_\_\_\_\_\_\_\_\_\_\_\_\_\_\_\_\_\_\_\_\_\_\_ 때문입니다.}
\kfAnswerNote{답안}{4}
\end{kfQBlock}


\kfStartNonlit{비문학 영역 — 사실과 의견, 무엇이 다를까?}


\kfSecHeader{kfAreaNonlit}{지문}{사실과 의견, 무엇이 다를까?}
\kfPassageNote{사실과 의견, 무엇이 다를까?}{%
　우리는 친구와 대화하거나 뉴스를 볼 때 수많은 말을 듣습니다. 이때 그 말이 '사실'인지 아니면 누군가의 '의견'인지 구분하는 것은 매우 중요합니다.

　'사실'이란 실제로 있었던 일이나 현재 있는 상태를 말합니다. 누가 보더라도 똑같이 인정할 수 있는 객관적인 내용입니다. 예를 들어 “이순신 장군은 거북선을 만들었다”, “물은 100도에서 끓는다”와 같은 문장은 사실입니다. 사실은 있는 그대로의 모습을 설명하며, 증거를 통해 참인지 거짓인지 입증할 수 있습니다.

\uline{반면에} '의견'은 어떤 대상이나 일에 대한 글쓴이의 생각이나 느낌을 말합니다. 사람마다 다를 수 있는 주관적인 내용입니다. “이순신 장군은 세상에서 가장 훌륭한 위인이다”, “라면은 김치와 먹어야 맛있다”와 같은 문장은 의견입니다. 의견을 말할 때는 보통 '나는 \textasciitilde{}라고 생각한다', '\textasciitilde{}인 것 같다', '\textasciitilde{}해야 한다'와 같은 표현을 자주 씁니다.

　우리가 글을 읽거나 말을 들을 때 사실과 의견을 혼동하면 올바른 판단을 내리기 어렵습니다. 사실이 아닌 것을 사실처럼 믿거나, 한 사람의 의견을 모든 사람의 생각인 것처럼 착각할 수 있기 때문입니다. 따라서 내용을 받아들일 때는 근거가 타당한지, \uline{그것}이 객관적인 사실인지 꼼꼼하게 따져보는 습관을 길러야 합니다.
}


\kfSecHeader{kfAreaNonlit}{활동}{차근차근 내용 확인}
\noindent\textit{\small [1-5] 글의 내용과 일치하면 O, 일치하지 않으면 X를 쓰세요.}\par\medskip
\begin{kfOXList}
\kfOXItem '사실'은 누가 보더라도 똑같이 인정할 수 있는 객관적인 내용이다. \kfOX
\kfOXItem '의견'은 사람마다 생각이 다를 수 있는 주관적인 내용이다. \kfOX
\kfOXItem 사실은 증거를 통해 참과 거짓을 증명하기 어렵다. \kfOX
\kfOXItem 글쓴이는 사실과 의견을 혼동하면 올바른 판단을 내리기 어렵다고 말했다. \kfOX
\kfOXItem “\textasciitilde{}라고 생각한다”, “\textasciitilde{}해야 한다”는 주로 의견을 말할 때 쓰는 표현이다. \kfOX
\end{kfOXList}

\kfSecHeader{kfAreaNonlit}{활동}{나만의 생각 펼치기}
\noindent\textit{\small 다음 조건에 맞게 글을 써 보세요.}\par\medskip
\kfWritingPrompt{다음 사실을 바탕으로 자신의 생각(의견)을 만들어 문장을 완성해 보세요. (사실): “우리 반 친구 30명 중 28명이 과학 탐구 올림픽에 참가했다.” (의견): “우리 반은 \_\_\_\_\_\_\_\_\_\_\_\_\_\_\_\_\_\_\_\_\_\_\_\_\_\_\_\_\_\_\_\_\_\_\_\_\_\_\_\_\_\_\_\_\_\_\_\_\_\_\_\_\_\_\_\_\_\_\_\_\_\_\_\_\_.”}
\kfWriting{답안 작성}{8}
\kfWritingPrompt{다음 사실을 바탕으로 자신의 생각(의견)을 만들어 문장을 완성해 보세요. (사실): “학교 급식에 매일 채소가 나오고 있다.” (의견): “나는 학교 급식에 \_\_\_\_\_\_\_\_\_\_\_\_\_\_\_\_\_\_\_\_\_\_\_\_\_\_\_\_\_\_\_\_\_\_\_\_\_\_\_\_\_\_\_\_\_\_\_\_\_\_\_\_\_\_\_\_\_\_\_.”}
\kfWriting{답안 작성}{8}
\kfWritingPrompt{내가 좋아하는 책이나 영화를 친구에게 추천하는 글을 한 문장으로 써보세요. (단, 사실과 의견이 각각 한 문장씩 포함해야 합니다.)}
\kfWriting{답안 작성}{8}

\kfSecHeader{kfAreaNonlit}{문제}{}

\begin{kfQBlock}{1}{}
\kfQStem{'㉠반면에'가 연결하는 앞문단과 뒷문단의 관계는 무엇인가요?}
\begin{kfChoices}
\kfChoice 앞의 내용에 대한 구체적인 예를 든다.
\kfChoice 앞의 내용이 원인이 되어 뒤의 결과가 나온다.
\kfChoice 앞의 내용과 뒤의 내용이 서로 반대되거나 대조된다.
\kfChoice 앞의 내용에 더해 비슷한 내용을 하나 더 추가한다.
\end{kfChoices}
\end{kfQBlock}
\begin{kfQBlock}{2}{}
\kfQStem{'㉡그것'이 가리키는 내용으로 알맞지 \uline{않은} 것은 무엇인가요?}
\begin{kfChoices}
\kfChoice 우리가 받아들이려는 내용
\kfChoice 글쓴이의 주관적인 느낌
\kfChoice 사람마다 다른 생각
\kfChoice 올바르지 못한 판단
\end{kfChoices}
\end{kfQBlock}
\begin{kfQBlock}{3}{}
\kfQStem{'원인'과 '결과'로 내용을 바르게 나눈 것은 무엇인가요?}
\begin{kfEvidence}사실과 의견을 혼동하면 올바른 판단을 내리기 어렵습니다.\end{kfEvidence}
\begin{kfChoices}
\kfChoice (원인) 올바른 판단을 내리기 어렵다 / (결과) 사실과 의견을 혼동한다
\kfChoice (원인) 사실과 의견을 혼동한다 / (결과) 올바른 판단을 내리기 어렵다
\kfChoice (원인) 사실과 의견을 구분한다 / (결과) 올바른 판단을 내릴 수 있다
\kfChoice (원인) 뉴스를 보며 수많은 말을 듣는다 / (결과) 사실과 의견을 혼동한다
\end{kfChoices}
\end{kfQBlock}
\begin{kfQBlock}{1}{}
\kfQStem{다음 중 '사실'을 설명하는 문장으로 가장 적절하지 \uline{않은} 것은 무엇인가요?}
\begin{kfChoices}
\kfChoice 누가 보더라도 똑같이 인정할 수 있는 내용이다.
\kfChoice 증거를 통해 참인지 거짓인지 입증할 수 있다.
\kfChoice '\textasciitilde{}인 것 같다'와 같은 표현을 자주 사용한다.
\kfChoice 있는 그대로의 모습을 설명하는 객관적인 내용이다.
\end{kfChoices}
\end{kfQBlock}
\begin{kfQBlock}{2}{}
\kfQStem{글쓴이가 사실과 의견을 구분해야 한다고 강조한 이유로 적절하지 \uline{않은} 것은 무엇인가요?}
\begin{kfChoices}
\kfChoice 글을 쓸 때 더 재미있는 표현을 사용하기 위해서
\kfChoice 사실이 아닌 것을 사실처럼 믿는 착각을 막기 위해서
\kfChoice 대화할 때 상대방의 기분을 좋게 만들어 주기 위해서
\kfChoice 의견을 말할 때 더 자신감을 가지기 위해서
\end{kfChoices}
\end{kfQBlock}
\begin{kfQBlock}{1}{}
\kfQStem{증거를 들어 옳거나 사실임을 증명함:}
\begin{kfEvidence}주관적\end{kfEvidence}
\kfAnswerLines{1}
\end{kfQBlock}
\begin{kfQBlock}{2}{}
\kfQStem{개인의 의견이나 감정 없이 사실을 있는 그대로 나타냄:}
\begin{kfEvidence}입증\end{kfEvidence}
\kfAnswerLines{1}
\end{kfQBlock}
\begin{kfQBlock}{3}{}
\kfQStem{개인적인 생각이나 감정에 바탕을 둠:}
\begin{kfEvidence}혼동\end{kfEvidence}
\kfAnswerLines{1}
\end{kfQBlock}
\begin{kfQBlock}{4}{}
\kfQStem{여러 가지를 뒤섞어 갈피를 잡지 못함:}
\kfAnswerLines{1}
\end{kfQBlock}
\begin{kfQBlock}{1}{}
\kfQStem{| 구분 | (1) 사실 | (2) 의견 | | :--- | :--- | :--- | | 뜻 | 실제로 있었던 일이나 현재 있는 상태 | 대상이나 일에 대한 개인의 ( ① )이나 느낌 | | 성격 | ( ② )적인 내용 (누구나 인정함) | 주관적인 내용 (사람마다 다름) | | 특징 | 증거를 통해 ( ③ )할 수 있음. | '\textasciitilde{}라고 생각한다', '\textasciitilde{}해야 한다' 등의 표현 사용 |}
\kfAnswerLines{1}
\end{kfQBlock}
\begin{kfQBlock}{1}{}
\kfQStem{“경주는 신라의 수도였다.”}
\kfAnswerLines{1}
\end{kfQBlock}
\begin{kfQBlock}{2}{}
\kfQStem{“경주는 우리나라에서 가장 아름다운 도시다.”}
\kfAnswerLines{1}
\end{kfQBlock}
\begin{kfQBlock}{3}{}
\kfQStem{“축구는 11명의 선수가 한 팀이 되어 경기한다.”}
\kfAnswerLines{1}
\end{kfQBlock}
\begin{kfQBlock}{4}{}
\kfQStem{“축구는 야구보다 훨씬 재미있는 스포츠다.”}
\kfAnswerLines{1}
\end{kfQBlock}
\begin{kfQBlock}{5}{}
\kfQStem{“오늘은 어제보다 기온이 5도 높다.”}
\kfAnswerLines{1}
\end{kfQBlock}
\begin{kfQBlock}{6}{}
\kfQStem{“오늘 날씨는 소풍 가기에 딱 좋은 날씨다.”}
\kfAnswerLines{1}
\end{kfQBlock}
\begin{kfQBlock}{1}{}
\kfQStem{“우리 가게의 신메뉴 '왕새우 피자'가 출시되었습니다. 이 피자의 지름은 30cm이며, 통새우가 10마리 들어갑니다. 여러분이 먹어본 피자 중에서 가장 환상적인 맛일 것입니다.” 사실: \_\_\_\_\_\_\_\_\_\_\_\_\_\_\_\_\_\_\_\_\_\_\_\_\_\_\_\_\_\_\_\_\_\_\_\_\_\_\_\_\_\_\_\_\_\_\_\_\_\_\_\_\_\_\_\_\_\_\_\_\_\_\_\_\_\_\_ 의견: \_\_\_\_\_\_\_\_\_\_\_\_\_\_\_\_\_\_\_\_\_\_\_\_\_\_\_\_\_\_\_\_\_\_\_\_\_\_\_\_\_\_\_\_\_\_\_\_\_\_\_\_\_\_\_\_\_\_\_\_\_\_\_\_\_\_\_}
\kfAnswerNote{답안}{4}
\end{kfQBlock}
\begin{kfQBlock}{2}{}
\kfQStem{“저는 어제 A브랜드의 러닝화를 샀습니다. 이 신발의 무게는 200g이고, 가격은 5만 원입니다. 신발이 가벼워서 마치 구름 위를 걷는 것처럼 편안한 느낌이 듭니다.” 사실: \_\_\_\_\_\_\_\_\_\_\_\_\_\_\_\_\_\_\_\_\_\_\_\_\_\_\_\_\_\_\_\_\_\_\_\_\_\_\_\_\_\_\_\_\_\_\_\_\_\_\_\_\_\_\_\_\_\_\_\_\_\_\_\_\_\_\_ 의견: \_\_\_\_\_\_\_\_\_\_\_\_\_\_\_\_\_\_\_\_\_\_\_\_\_\_\_\_\_\_\_\_\_\_\_\_\_\_\_\_\_\_\_\_\_\_\_\_\_\_\_\_\_\_\_\_\_\_\_\_\_\_\_\_\_\_\_}
\kfAnswerNote{답안}{4}
\end{kfQBlock}
\begin{kfQBlock}{3}{}
\kfQStem{“소설 <소나기>는 황순원 작가가 쓴 단편 소설입니다. 시골 소년과 도시 소녀의 짧고 순수한 사랑 이야기를 담고 있습니다. 읽으면 누구나 눈물을 흘릴 만큼 슬프고 아름다운 이야기입니다.” 사실: \_\_\_\_\_\_\_\_\_\_\_\_\_\_\_\_\_\_\_\_\_\_\_\_\_\_\_\_\_\_\_\_\_\_\_\_\_\_\_\_\_\_\_\_\_\_\_\_\_\_\_\_\_\_\_\_\_\_\_\_\_\_\_\_\_\_\_ 의견: \_\_\_\_\_\_\_\_\_\_\_\_\_\_\_\_\_\_\_\_\_\_\_\_\_\_\_\_\_\_\_\_\_\_\_\_\_\_\_\_\_\_\_\_\_\_\_\_\_\_\_\_\_\_\_\_\_\_\_\_\_\_\_\_\_\_\_}
\kfAnswerNote{답안}{4}
\end{kfQBlock}
\begin{kfQBlock}{4}{}
\kfQStem{“드림 놀이공원은 오전 9시에 문을 엽니다. 이곳에는 높이 50m에서 떨어지는 롤러코스터가 있습니다. 이 롤러코스터는 세상에서 가장 짜릿하고 무서운 놀이기구입니다.” 사실: \_\_\_\_\_\_\_\_\_\_\_\_\_\_\_\_\_\_\_\_\_\_\_\_\_\_\_\_\_\_\_\_\_\_\_\_\_\_\_\_\_\_\_\_\_\_\_\_\_\_\_\_\_\_\_\_\_\_\_\_\_\_\_\_\_\_\_ 의견: \_\_\_\_\_\_\_\_\_\_\_\_\_\_\_\_\_\_\_\_\_\_\_\_\_\_\_\_\_\_\_\_\_\_\_\_\_\_\_\_\_\_\_\_\_\_\_\_\_\_\_\_\_\_\_\_\_\_\_\_\_\_\_\_\_\_\_}
\kfAnswerNote{답안}{4}
\end{kfQBlock}
\begin{kfQBlock}{5}{}
\kfQStem{“지난 주말에 개봉한 영화 <우주 대탐험>을 보았습니다. 상영 시간은 총 2시간이었고, 우주 비행사 3명이 주인공으로 나옵니다. 올해 본 영화 중에서 지루하지 않고 제일 재미있었습니다.” 사실: \_\_\_\_\_\_\_\_\_\_\_\_\_\_\_\_\_\_\_\_\_\_\_\_\_\_\_\_\_\_\_\_\_\_\_\_\_\_\_\_\_\_\_\_\_\_\_\_\_\_\_\_\_\_\_\_\_\_\_\_\_\_\_\_\_\_\_ 의견: \_\_\_\_\_\_\_\_\_\_\_\_\_\_\_\_\_\_\_\_\_\_\_\_\_\_\_\_\_\_\_\_\_\_\_\_\_\_\_\_\_\_\_\_\_\_\_\_\_\_\_\_\_\_\_\_\_\_\_\_\_\_\_\_\_\_\_}
\kfAnswerNote{답안}{4}
\end{kfQBlock}


\kfStartWeekly{실력확인 영역 — 주간 실력 확인}


\noindent\textit{\small 제한시간: 30분 / 총 20문항}\par\medskip
\par\noindent\textbf{[6-10] 다음 글을 읽고, 물음에 답하세요.}\par\medskip
\kfPassageNote{지문 6}{%
[앞부분 줄거리] 용왕의 병을 고치기 위해 별주부는 육지로 나와 토끼를 꾀어 용궁으로 데려간다. 용왕이 토끼의 배를 갈라 간을 꺼내려 하자, 토끼는 꾀를 내어 말하기 시작한다.

용왕이 분부하였다. “토끼를 묶어라! 내 병이 위중하니 빨리 배를 갈라 간을 내오너라!” 나졸들이 우르르 달려들어 토끼를 묶고 배를 가르려 하니, 토끼가 눈을 끔적이며 급히 여쭈었다. “전하, 잠깐 제 말을 들어보십시오. 세상의 모든 짐승은 간이 배 속에 붙어 있으나, \uline{저}만은 간을 넣었다 뺐다 할 수 있사옵니다.” 용왕이 기가 막혀 물었다. “네 말이 당치 않도다. 어찌 간을 배 밖에 두고 다닌단 말이냐?” 토끼가 태연하게 대답했다. “제가 사는 육지는 맹수들이 많아 간을 배 속에 넣고 다니면 썩거나 터지기 쉽습니다. 그래서 평소에는 깊은 산속 푸른 바위 밑에 감추어 두었다가, 필요할 때만 배 속에 넣어 다니옵니다. 이번에는 별주부님이 너무 급하게 재촉하는 바람에 미처 챙기지 못하고 그냥 왔나이다.” 용왕은 잠시 고민하다가 토끼의 말을 믿기로 하였다.
}
\par\noindent\textbf{[11-15] 다음 글을 읽고, 물음에 답하세요.}\par\medskip
\kfPassageNote{지문 11}{%
우리는 친구와 대화하거나 뉴스를 볼 때 수많은 말을 듣습니다. 이때 그 말이 '사실'인지 아니면 누군가의 '의견'인지 구분하는 것은 매우 중요합니다.

'사실'이란 실제로 있었던 일이나 현재에 있는 일이나 상태를 말합니다. 누가 보더라도 똑같이 인정할 수 있는 객관적인 내용입니다. 예를 들어 “이순신 장군은 거북선을 만들었다”, “물은 100도에서 끓는다”와 같은 문장은 사실입니다. 사실은 있는 그대로의 모습을 설명하며, 증거를 통해 참인지 거짓인지 입증할 수 있습니다.

반면에 '의견'은 어떤 대상이나 일에 대한 글쓴이의 생각이나 느낌을 말합니다. 사람마다 다를 수 있는 주관적인 내용입니다. “이순신 장군은 세상에서 가장 훌륭한 위인이다”, “라면은 김치와 먹어야 맛있다”와 같은 문장은 의견입니다. 의견을 말할 때는 보통 '나는 \textasciitilde{}라고 생각한다', '\textasciitilde{}인 것 같다', '\textasciitilde{}해야 한다'와 같은 표현을 자주 씁니다.

우리가 글을 읽거나 말을 들을 때 사실과 의견을 혼동하면 올바른 판단을 내리기 어렵습니다. 사실이 아닌 것을 사실처럼 믿거나, 한 사람의 의견을 모든 사람의 생각인 것처럼 착각할 수 있기 때문입니다. 따라서 내용을 받아들일 때는 근거가 타당한지, 그것이 객관적인 사실인지 꼼꼼하게 따져보는 습관을 길러야 합니다.
}
\begin{kfQBlock}{1}{}
\kfQStem{다음 글의 빈칸 ㉠, ㉡에 들어갈 말로 가장 적절한 것은 무엇인가요?}
\begin{kfEvidence}토론을 할 때 상대방의 의견을 무조건 비난하는 것은 옳지 않다. 상대방 주장의 잘못된 점을 분석하고 대안을 제시하는 건전한 ㉠[　　　　　]을/를 해야 한다. 또한 자신의 감정보다는 ㉡[　　　　　]인 사실에 바탕을 두어야 한다.\end{kfEvidence}
\begin{kfChoices}
\kfChoice ㉠ 근거 - ㉡ 주관적
\kfChoice ㉠ 입증 - ㉡ 논리적
\kfChoice ㉠ 비판 - ㉡ 객관적
\kfChoice ㉠ 타당성 - ㉡ 편견
\end{kfChoices}
\end{kfQBlock}

\begin{kfQBlock}{2}{}
\kfQStem{다음 문장의 빈칸에 들어갈 가장 알맞은 단어는 무엇인가요?}
\begin{kfEvidence}토론에서 중요한 것은 감정에 치우쳐 상대방을 [　　　　　]하는 것이 아니라, 논리적인 근거를 제시하는 것입니다.\end{kfEvidence}
\begin{kfChoices}
\kfChoice 의견
\kfChoice 주장
\kfChoice 사실
\kfChoice 비난
\end{kfChoices}
\end{kfQBlock}

\begin{kfQBlock}{3}{}
\kfQStem{다음 대화의 빈칸에 들어갈 단어로 가장 적절한 것은 무엇인가요?}
\begin{kfEvidence}변호사: “피고인이 죄가 없다는 것은 분명한 사실입니다.” 판사: “그 주장을 믿으려면 증거를 들어서 [　　　　　]해야 합니다. 말만으로는 부족합니다.”\end{kfEvidence}
\begin{kfChoices}
\kfChoice 비판
\kfChoice 입증
\kfChoice 결론
\kfChoice 편견
\end{kfChoices}
\end{kfQBlock}

\begin{kfQBlock}{4}{}
\kfQStem{'사실'과 '의견'의 관계처럼, 단어의 성격이 서로 반대되는 것은 무엇인가요?}
\begin{kfChoices}
\kfChoice 근거 - 이유
\kfChoice 입증 - 증명
\kfChoice 타당성 - 올바름
\kfChoice 객관적 - 주관적
\end{kfChoices}
\end{kfQBlock}

\begin{kfQBlock}{5}{}
\kfQStem{다음 문장에서 밑줄 친 단어의 사용이 \uline{어색한} 것은 무엇인가요?}
\begin{kfChoices}
\kfChoice 사실 그대로를 바라보는 \uline{객관적인} 태도가 필요하다.
\kfChoice 올바른 \uline{비판}은 친구가 잘못을 고칠 기회를 준다.
\kfChoice 그는 자신의 무죄를 밝히기 위해 확실한 \uline{근거}를 댔다.
\kfChoice 이 주장은 앞뒤가 전혀 맞지 않아 매우 \uline{논리적}이야.
\end{kfChoices}
\end{kfQBlock}

\begin{kfQBlock}{6}{}
\kfQStem{위 글의 내용을 바르게 이해한 것은 무엇인가요?}
\begin{kfChoices}
\kfChoice 용왕은 토끼가 미워서 죽이려고 했다.
\kfChoice 토끼는 간을 실제로 집에 두고 왔다.
\kfChoice 별주부는 토끼에게 간을 챙길 시간을 충분히 주었다.
\kfChoice 토끼는 죽을 위기에 처했지만 꾀를 내어 대응했다.
\end{kfChoices}
\end{kfQBlock}

\begin{kfQBlock}{7}{}
\kfQStem{㉠에 담긴 토끼의 의도로 가장 적절한 것은 무엇인가요?}
\begin{kfChoices}
\kfChoice 자신이 특별한 동물임을 자랑하기 위해서
\kfChoice 간이 배 속에 없다는 사실을 고백하기 위해서
\kfChoice 용왕을 혼란스럽게 만들어 시간을 벌기 위해서
\kfChoice 용왕의 병을 고칠 다른 방법을 알려주기 위해서
\end{kfChoices}
\end{kfQBlock}

\begin{kfQBlock}{8}{}
\kfQStem{토끼가 말한 '거짓말의 근거'로 적절하지 \uline{않은} 것은 무엇인가요?}
\begin{kfChoices}
\kfChoice 육지에는 맹수들이 많아서 위험하다.
\kfChoice 간을 배 속에 계속 넣고 다니면 썩을 수 있다.
\kfChoice 별주부가 너무 서두르는 바람에 챙겨오지 못했다.
\kfChoice 내 간은 너무 무거워서 평소에는 빼놓고 다닌다.
\end{kfChoices}
\end{kfQBlock}

\begin{kfQBlock}{9}{}
\kfQStem{<보기>를 참고하여 위 글을 감상한 내용으로 가장 적절한 것은 무엇인가요?}
\begin{kfEvidence}<보기> 　「토끼전」은 조선 후기 평범한 백성들의 생각을 담은 이야기입니다. 이 이야기 속에 등장하는 동물들은 당시 사회의 다양한 사람들의 모습을 상징합니다. 　* 토끼: 힘은 없지만 꾀와 지혜를 발휘하여 위기를 극복하는 '힘없는 백성' 　* 용왕: 자신의 욕심을 채우기 위해 백성을 희생시키려는 '무능하고 부패한 지배층'\end{kfEvidence}
\begin{kfChoices}
\kfChoice 토끼가 별주부의 호의에 끌려간 모습은, 백성들이 높은 사람의 도움에 의지해 살아갔음을 보여주는군.
\kfChoice 별주부가 끝까지 임무를 완수하려 한 모습은, 백성이라면 위에서 시킨 일을 의심 없이 따라야 함을 강조한 것이군.
\kfChoice 용왕이 토끼의 말을 곧이 믿은 모습은, 지배층이 백성보다 더 어리석을 수 있다는 점만을 풍자한 것이군.
\kfChoice 약한 토끼가 강한 용왕을 꾀로 속여 살아 돌아오는 모습은, 힘센 지배층의 횡포에 맞서 이기고 싶었던 백성들의 마음을 담은 것이군.
\end{kfChoices}
\end{kfQBlock}

\begin{kfQBlock}{10}{}
\kfQStem{위 글의 '토끼'가 처한 상황과 대처 방법을 가장 잘 나타낸 속담을 <보기>에서 고른 것은 무엇인가요?}
\begin{kfEvidence}<보기> 　ㄱ. 닭 쫓던 개 지붕 쳐다본다. 　ㄴ. 원숭이도 나무에서 떨어질 때가 있다. 　ㄷ. 얌전한 고양이 부뚜막에 먼저 올라간다. 　ㄹ. 호랑이에게 물려가도 정신만 차리면 산다.\end{kfEvidence}
\begin{kfChoices}
\kfChoice ㄱ
\kfChoice ㄴ
\kfChoice ㄷ
\kfChoice ㄹ
\end{kfChoices}
\end{kfQBlock}

\begin{kfQBlock}{11}{}
\kfQStem{위 글의 내용을 바탕으로 할 때, <보기>의 학생들 중 적절하지 \uline{않은} 생각을 가진 학생은 누구인가요?}
\begin{kfEvidence}<보기> 　* 지수: “사실은 누가 봐도 똑같아야 하니까 증거를 통해 진짜인지 가짜인지 증명할 수 있어야 해.” 　* 민호: “의견은 사람마다 생각이 다를 수 있으니까 친구의 의견이 내 생각과 조금 달라도 이해해야 해.” 　* 영희: “사실과 의견을 헷갈리면 한 사람의 주관적인 생각을 마치 모두가 인정하는 사실인 것처럼 잘못 믿을 수도 있어.” 　* 철수: “의견은 자기만의 느낌이나 생각을 자유롭게 말하는 거니까 그 의견이 타당한지 근거를 따져 볼 필요는 없어.”\end{kfEvidence}
\begin{kfChoices}
\kfChoice 지수
\kfChoice 민호
\kfChoice 영희
\kfChoice 철수
\end{kfChoices}
\end{kfQBlock}

\begin{kfQBlock}{12}{}
\kfQStem{위 글을 참고하여 <보기>의 상황을 비판한 내용으로 가장 적절한 것은 무엇인가요?}
\begin{kfEvidence}<보기> 　민수: “야, 이 건강음료 좀 봐! 광고에 '세상에서 가장 맛있는 음료'라고 적혀 있어. 게다가 '마시면 누구나 힘이 불끈 솟는다'고 하네. 이건 정말 대단한 음료수가 분명해! 꼭 사 먹어야지.”\end{kfEvidence}
\begin{kfChoices}
\kfChoice 민수는 광고에 나온 내용이 모두 거짓말이라고 의심하고 있어.
\kfChoice 민수는 다른 사람의 의견보다 자신의 경험을 더 중요하게 생각하고 있어.
\kfChoice 민수는 근거의 타당성을 꼼꼼하게 따져본 후 음료수를 선택하고 있어.
\kfChoice 민수는 주관적인 의견을 객관적인 사실인 것처럼 착각하고 있어.
\end{kfChoices}
\end{kfQBlock}

\begin{kfQBlock}{13}{}
\kfQStem{(가)\textasciitilde{}(라)를 '사실'과 '의견'으로 바르게 분류한 것은 무엇인가요?}
\begin{kfEvidence}(가) 독도는 대한민국의 가장 동쪽에 있는 섬이다. 행정 구역상으로는 경상북도 울릉군 울릉읍 독도리에 속한다.

(나) 식물의 잎은 햇빛을 이용하여 스스로 양분을 만든다. 이러한 작용을 과학 용어로 '광합성'이라고 한다.

(다) 아침밥을 못 먹고 허겁지겁 학교에 오는 친구들이 많다. 학생들의 건강을 위해서 학교 등교 시간을 9시 30분으로 늦추는 것이 좋다고 생각한다.

(라) 가을 하늘은 사계절 중 가장 아름답다. 특히 노을이 질 때의 하늘 색깔은 그 어떤 보석보다도 신비롭고 황홀하다.\end{kfEvidence}
\begin{kfChoices}
\kfChoice 사실: (가), (나) / 의견: (다), (라)
\kfChoice 사실: (가), (다) / 의견: (나), (라)
\kfChoice 사실: (나), (다) / 의견: (가), (라)
\kfChoice 사실: (다), (라) / 의견: (가), (나)
\end{kfChoices}
\end{kfQBlock}

\begin{kfQBlock}{14}{}
\kfQStem{다음 중 객관적인 사실을 말하고 있는 사람은 누구인가요?}
\begin{kfChoices}
\kfChoice 준호: “저 사과는 껍질이 아주 빨간 걸 보니, 분명 맛있을 거야.”
\kfChoice 민수: “환경을 보호하기 위해서는 주방 세제를 사용해서는 안 된다.”
\kfChoice 지민: “수학은 논리적인 사고력을 기르는 데 가장 유익한 과목이야.”
\kfChoice 수진: “임진왜란은 1592년에 일본이 조선을 침략하면서 일어난 전쟁이야.”
\end{kfChoices}
\end{kfQBlock}

\begin{kfQBlock}{15}{}
\kfQStem{다음 '의견' 문장을 객관적인 '사실' 문장으로 바꾸어 표현할 때, 가장 적절하지 \uline{않은} 것은?}
\begin{kfEvidence}의견: “흰수염고래(대왕고래)는 지구상에서 가장 거대하고 놀라운 동물이다.”\end{kfEvidence}
\begin{kfChoices}
\kfChoice 흰수염고래의 몸길이는 최대 33미터까지 자란다.
\kfChoice 흰수염고래가 헤엄치는 모습은 마치 우아한 무용수처럼 아름답다.
\kfChoice 흰수염고래의 심장 무게는 약 600kg에서 900kg 사이로 소형 자동차 한 대의 무게와 비슷하다.
\kfChoice 흰수염고래는 물속에서 시속 5\textasciitilde{}20km로 이동하며, 위급할 때는 시속 40km 이상의 속도를 낼 수 있다.
\end{kfChoices}
\end{kfQBlock}

\begin{kfQBlock}{16}{}
\kfQStem{다음 문장에서 '주어'에 해당하는 말은 무엇인가요?}
\begin{kfEvidence}파란 하늘에 하얀 구름이 둥실 떠 있다.\end{kfEvidence}
\begin{kfChoices}
\kfChoice 파란
\kfChoice 하늘에
\kfChoice 구름이
\kfChoice 떠 있다
\end{kfChoices}
\end{kfQBlock}

\begin{kfQBlock}{17}{}
\kfQStem{다음 문장에서 서술어 '없다'의 주어로 알맞은 것은 무엇인가요?}
\begin{kfEvidence}나는 어제 지우개를 잃어버려서, 지금 필통 속에 지우개가 하나도 없다.\end{kfEvidence}
\begin{kfChoices}
\kfChoice 나는
\kfChoice 어제
\kfChoice 지우개가
\kfChoice 필통 속에
\end{kfChoices}
\end{kfQBlock}

\begin{kfQBlock}{18}{}
\kfQStem{다음 문장에서 주어의 동작이나 상태를 풀이하는 말은 무엇인가요?}
\begin{kfEvidence}이순신 장군은 거북선을 만들었다.\end{kfEvidence}
\begin{kfChoices}
\kfChoice 이순신
\kfChoice 장군은
\kfChoice 거북선을
\kfChoice 만들었다
\end{kfChoices}
\end{kfQBlock}

\begin{kfQBlock}{19}{}
\kfQStem{다음 <보기>의 설명에 해당하는 문장 성분이 빠져 있는 문장은 무엇인가요?}
\begin{kfEvidence}<보기> 　이 성분은 문장의 주인입니다. 동작이나 상태의 주체가 되며, 주로 '누가', '무엇이'에 해당합니다.\end{kfEvidence}
\begin{kfChoices}
\kfChoice 선생님께서 오신다.
\kfChoice 어제 도서관에 갔어.
\kfChoice 꽃이 예쁘게 피었다.
\kfChoice 우리 아빠는 회사원이다.
\end{kfChoices}
\end{kfQBlock}

\begin{kfQBlock}{20}{}
\kfQStem{다음 중 주어와 서술어의 호응 관계가 어색한 문장을 모두 고른 것은 무엇인가요?}
\begin{kfEvidence}　㉠ 내 장래 희망은 훌륭한 요리사가 되고 싶다. 　㉡ 내가 학교에 늦은 까닭은 버스를 놓쳤기 때문이다. 　㉢ 이 가방의 장점은 주머니가 많아서 편리하다는 점이다. 　㉣ 우리가 해야 할 일은 교실을 깨끗이 청소하자.\end{kfEvidence}
\begin{kfChoices}
\kfChoice ㉠, ㉡
\kfChoice ㉠, ㉣
\kfChoice ㉡, ㉢
\kfChoice ㉢, ㉣
\end{kfChoices}
\end{kfQBlock}


\kfSecHeader{kfAreaWeekly}{문제}{}


\kfSecHeader{kfAreaWeekly}{활동}{주간 종합 평가}
\noindent\textit{\small 다음 빈칸에 알맞은 말을 골라 채우시오.}\par\medskip



\end{document}
